%% Volume III, Chapter 4: Uniform Convergence
\chapter{Uniform Convergence}
\label{ch:uniform-convergence}

\abstract{Almost every computation of analysis exchanges two limits: a series is differentiated term by term, a limit is taken under an integral sign, a double sum is summed in the other order. Convergence of functions at each point justifies none of these steps. A limit of continuous functions can jump, and the integrals of a sequence of functions can stay away from the integral of its limit. The answer that the nineteenth century found is uniform convergence, which is convergence in the supremum norm; it is also how numerical work measures the error of a partial sum that replaces a function. At the end of the chapter, uniform limits produce continuous functions that have a derivative nowhere.}

In his \emph{Cours d'analyse} of 1821, Cauchy stated as a theorem that the sum of a convergent series of continuous functions is continuous~\cite{cauchy1821}. Five years later Abel observed, in his paper on the binomial series, that the theorem ``admits exceptions'', and gave the series
\[
\sin x-\frac12\sin2x+\frac13\sin3x-\cdots,
\]
whose terms are continuous and whose sum jumps at the odd multiples of $\pi$~\cite{abel1826}. The flaw in Cauchy's argument is not obvious. The proof that a sum $f=\sum_kf_k$ is continuous at $a$ passes from the partial sums $s_n$, which are continuous, to $f$, and it needs that one index $n$ makes $s_n$ close to $f$ at $a$ and at all points near $a$ at the same time. Convergence at each point gives, for each point, an index of its own, and nothing prevents these indices from growing without bound as the points approach $a$. The missing hypothesis, that a single index serves all points at once, is \emph{uniform convergence}.

The same issue appears whenever two limiting processes meet. Continuity of the limit means
\[
\lim_{x\to a}\lim_{n\to\infty}s_n(x)=\lim_{n\to\infty}\lim_{x\to a}s_n(x);
\]
integration term by term exchanges a limit with the limit of Riemann sums; differentiation term by term exchanges a limit with the limit of difference quotients; and summing a double series by rows or by columns exchanges two summations. Each exchange can fail, and \cref{sec:pointwise-uniform} collects the standard counterexamples in \cref{tab:interchange-failures}. Uniform convergence, sometimes with a little more, makes each exchange legitimate. Since $\sup_{x\in S}|f_n(x)-f(x)|$ is the distance of $f_n$ from $f$ in the supremum norm of \cref{sec:metrics-norms}, uniform convergence is convergence in the space $\ell^\infty(S)$, and the completeness of that space, \cref{thm:bounded-functions-complete}, and the M-test after \cref{thm:absolute-series} are already the first theorems of the subject. \Cref{ch:metric-spaces} proved one exchange theorem, \cref{lem:uniform-limit}: a uniform limit of continuous functions is continuous. This chapter develops the general theory that the note on uniform convergence in \vcref{1}{sec:power-series} and \vcref{2}{lem:vector-functions} announced.

\Cref{sec:pointwise-uniform} defines pointwise and uniform convergence and shows how to decide which of them holds. \Cref{sec:limit-continuity} proves the general theorem on the exchange of two limits, of which the continuity of uniform limits is a special case, and Dini's theorem, by which compactness turns monotone pointwise convergence into uniform convergence; as an application it approximates $\sqrt x$ by polynomials. \Cref{sec:limit-integration} proves that integrals and, under a condition on the derivatives, derivatives pass to uniform limits. \Cref{sec:function-series} applies the results to series of functions, with the M-test of Weierstrass and the tests of Abel and Dirichlet, and \cref{sec:double-series} treats double sequences and double series. \Cref{sec:power-series-revisited} returns to the power series of \vcref{1}{sec:power-series}: it proves that their sums are analytic at every point of the interval of convergence and that they are continuous up to an end at which the series converges, which is Abel's theorem, and it compares Abel's method of summation with Ces\`aro's and with Tauber's converse. \Cref{sec:nowhere-differentiable} constructs continuous functions that are differentiable nowhere, as uniform limits of smooth or piecewise linear functions, and a continuous curve that fills a square. \Cref{sec:uniform-computation} measures uniform errors in computation.

The conventions of \cref{ch:metric-spaces} remain in force. $\mathbb F$ is $\mathbb R$ or $\mathbb C$, functions on a set $S$ take values in $\mathbb F$, in a normed space or in a metric space as stated, and sequences of functions $(f_n)$ are indexed from $0$ unless stated otherwise. For a function $g$ on $S$ with values in a normed space we write $\|g\|_\infty\coloneqq\sup_{s\in S}\|g(s)\|$, now allowing the value $\infty$; we say on which set the supremum is taken when it is not clear from the context. The integrals are the Riemann integrals of \vcref{1}{ch:integration}; the Lebesgue integral of Part~III has convergence theorems of its own, and the chapter indicates where they go further.

\section{Pointwise and Uniform Convergence}
\label{sec:pointwise-uniform}

\subsection{Two Notions of Convergence}

\begin{definition}[Pointwise and uniform convergence]
\label{def:pointwise-uniform}
Let $S$ be a set, $(Y,d)$ a metric space, and $f_n\colon S\to Y$ ($n\in\mathbb N$) and $f\colon S\to Y$ functions.
\begin{enumerate}
\item The sequence $(f_n)$ \term{converges pointwise} to $f$ on $S$ if $f_n(s)\to f(s)$ for every $s\in S$: for every $s\in S$ and every $\varepsilon>0$ there is $N$ with $d(f_n(s),f(s))<\varepsilon$ for all $n\ge N$.
\item The sequence $(f_n)$ \term{converges uniformly} to $f$ on $S$ if for every $\varepsilon>0$ there is $N$ with $d(f_n(s),f(s))<\varepsilon$ for all $n\ge N$ and all $s\in S$; equivalently, if
\[
\sup_{s\in S}d\bigl(f_n(s),f(s)\bigr)\to0\qquad(n\to\infty),
\]
where the suprema lie in $[0,\infty]$ and are finite for all large $n$.
\end{enumerate}
\end{definition}

The two definitions differ only in the place of one quantifier. In (i) the index $N$ may depend on $s$ as well as on $\varepsilon$; in (ii) it depends on $\varepsilon$ alone, and one index serves all points at once. The relation is that of continuity and uniform continuity in \vcref{1}{tab:quantifiers}. Uniform convergence implies pointwise convergence to the same function, and the pointwise limit, if it exists, is unique; so the only candidate for a uniform limit is the pointwise limit, and the question in every example is whether the convergence to it is uniform. Uniform convergence on $S$ implies uniform convergence on every subset of $S$, and uniform convergence on each of finitely many sets implies uniform convergence on their union, with the largest of the finitely many indices.

When $Y$ is $\mathbb F$ or a normed space, the supremum in (ii) is the norm $\|f_n-f\|_\infty$ of \cref{sec:metrics-norms}, extended to functions that need not be bounded.

\begin{proposition}[Uniform convergence and the supremum norm]
\label{prop:uniform-sup-norm}
Let $S$ be a nonempty set, $Y$ a normed space, and $f_n,f\colon S\to Y$.
\begin{enumerate}
\item $f_n\to f$ uniformly on $S$ if and only if $\|f_n-f\|_\infty\to0$.
\item If the $f_n$ are bounded and $f_n\to f$ uniformly, then $f$ is bounded, and $f_n\to f$ in the normed space $\ell^\infty(S,Y)$ of bounded functions $S\to Y$ with the norm $\|\cdot\|_\infty$; conversely, convergence in $\ell^\infty(S,Y)$ is uniform convergence.
\end{enumerate}
\end{proposition}

\begin{proof}
(i) is the definition. (ii) Choose $N$ with $\|f_N-f\|_\infty\le1$; then $\|f(s)\|\le\|f_N\|_\infty+1$ for every $s$, so $f$ is bounded, and $f_n\to f$ in $\ell^\infty(S,Y)$ by (i). The converse is (i) again.
\end{proof}

For $Y=\mathbb F$ the space $\ell^\infty(S,\mathbb F)$ is the space $\ell^\infty(S)$ of \cref{sec:metrics-norms}, and the same verification shows that $\ell^\infty(S,Y)$ is a normed space for every normed space $Y$. A uniform limit of unbounded functions may be unbounded, as $f_n(x)=x+1/(n+1)$ on $\mathbb R$ shows; then $\|f_n-f\|_\infty=1/(n+1)$, although no $f_n$ has a finite norm. Uniform convergence is therefore slightly more general than convergence in $\ell^\infty(S,Y)$, and we use the extended notation $\|g\|_\infty\in[0,\infty]$ throughout.

\subsection{First Examples}

The following four sequences on intervals of the real line are the standard examples; \cref{fig:basic-examples} shows them.

\begin{enumerate}
\item The monomials $f_n(x)\coloneqq x^n$ on $[0,1]$ converge pointwise to the function $f$ with $f(x)=0$ for $0\le x<1$ and $f(1)=1$, since $x^n\to0$ for $0\le x<1$ (\vcref{1}{thm:standard-limits}). The convergence is not uniform: $\sup_{0\le x\le1}|x^n-f(x)|=\sup_{0\le x<1}x^n=1$ for every $n$. It is not uniform on $[0,1)$ either, where the limit is continuous; but it is uniform on every interval $[0,a]$ with $a<1$, since $\sup_{0\le x\le a}x^n=a^n\to0$. The trouble is concentrated at the point $1$: for $x=2^{-1/n}$, close to $1$, the value $x^n$ is still $1/2$.
\item The functions $f_n(x)\coloneqq nxe^{-nx}$ on $[0,\infty)$ converge pointwise to $0$: $f_n(0)=0$, and for $x>0$, $nxe^{-nx}\to0$ as $n\to\infty$, since an exponential beats a power. The derivative $n(1-nx)e^{-nx}$ vanishes at $x=1/n$, where $f_n$ attains its maximum $e^{-1}$; so $\|f_n\|_\infty=e^{-1}$ for $n\ge1$, and the convergence is not uniform. On $[\delta,\infty)$ with $\delta>0$ it is uniform, since $f_n$ decreases on $[1/n,\infty)$ and therefore $\sup_{x\ge\delta}f_n(x)=n\delta e^{-n\delta}$ as soon as $1/n\le\delta$, and $n\delta e^{-n\delta}\to0$. The graphs have a hump of constant height that moves into the point $0$.
\item The tents $f_n(x)\coloneqq\max(0,1-|nx-1|)$ on $[0,1]$, for $n\ge1$, vanish outside $[0,2/n]$ and rise linearly to the value $1$ at $x=1/n$. They converge pointwise to $0$: $f_n(0)=0$, and $f_n(x)=0$ once $2/n<x$. Again $\|f_n\|_\infty=1$, so the convergence is not uniform on $[0,1]$, although it is uniform on every $[\delta,1]$ with $\delta>0$, where $f_n=0$ for $n>2/\delta$. The tent narrows as it moves into $0$, and its area $1/n$ tends to $0$; the taller tents $nf_n$, of height $n$ and area $1$, will serve as a counterexample for integrals.
\item The functions $f_n(x)\coloneqq\sin(nx)/n$ on $\mathbb R$, $n\ge1$, converge uniformly to $0$, since $\|f_n\|_\infty=1/n$. Their derivatives $f_n'(x)=\cos(nx)$ do not converge at all, except at the points $x\in2\pi\mathbb Z$, where they are equal to $1$.
\end{enumerate}

\begin{figure}
\sdpanel{A}{\begin{tikzpicture}
\begin{axis}[sdplot, width=0.46\textwidth, height=0.34\textwidth, xmin=0, xmax=1.02, ymin=0, ymax=1.08, xtick={0,0.5,1}, ytick={0,0.5,1}, xlabel={$x$}, legend pos=north west]
\addplot[black, line width=0.7pt, domain=0:1, samples=60] {x};
\addlegendentry{$n=1$}
\addplot[black, dashed, line width=0.7pt, domain=0:1, samples=80] {x^4};
\addlegendentry{$n=4$}
\addplot[black, densely dotted, line width=0.8pt, domain=0:1, samples=160] {x^16};
\addlegendentry{$n=16$}
\addplot[black!45, line width=0.7pt, domain=0:1, samples=200] {x^64};
\addlegendentry{$n=64$}
\end{axis}
\end{tikzpicture}}\hfill
\sdpanel{B}{\begin{tikzpicture}
\begin{axis}[sdplot, width=0.46\textwidth, height=0.34\textwidth, xmin=0, xmax=2, ymin=0, ymax=0.62, xtick={0,1,2}, ytick={0,0.2,0.4,0.6}, xlabel={$x$}]
\addplot[black, line width=0.7pt, domain=0:2, samples=100] {x*exp(-x)};
\addlegendentry{$n=1$}
\addplot[black, dashed, line width=0.7pt, domain=0:2, samples=120] {4*x*exp(-4*x)};
\addlegendentry{$n=4$}
\addplot[black, densely dotted, line width=0.8pt, domain=0:2, samples=200] {16*x*exp(-16*x)};
\addlegendentry{$n=16$}
\draw[black!45, densely dashed, line width=0.4pt] (axis cs:0,0.3679) -- (axis cs:2,0.3679);
\end{axis}
\end{tikzpicture}}
\sdpanelrow
\sdpanel{C}{\begin{tikzpicture}
\begin{axis}[sdplot, width=0.46\textwidth, height=0.34\textwidth, xmin=0, xmax=1.02, ymin=0, ymax=1.08, xtick={0,0.5,1}, ytick={0,0.5,1}, xlabel={$x$}]
\addplot[black, line width=0.7pt] coordinates {(0,0) (0.5,1) (1,0)};
\addlegendentry{$n=2$}
\addplot[black, dashed, line width=0.7pt] coordinates {(0,0) (0.125,1) (0.25,0) (1,0)};
\addlegendentry{$n=8$}
\addplot[black, densely dotted, line width=0.8pt] coordinates {(0,0) (0.03125,1) (0.0625,0) (1,0)};
\addlegendentry{$n=32$}
\end{axis}
\end{tikzpicture}}\hfill
\sdpanel{D}{\begin{tikzpicture}
\begin{axis}[sdplot, width=0.46\textwidth, height=0.34\textwidth, xmin=0, xmax=6.2832, ymin=-0.55, ymax=0.88, xtick={0,3.1416,6.2832}, xticklabels={$0$,$\pi$,$2\pi$}, ytick={-0.5,0,0.5}, xlabel={$x$}, legend columns=-1, legend style={at={(0.5,0.98)}, anchor=north}]
\addplot[black, line width=0.7pt, domain=0:6.2832, samples=100] {sin(deg(2*x))/2};
\addlegendentry{$n=2$}
\addplot[black, dashed, line width=0.7pt, domain=0:6.2832, samples=200] {sin(deg(6*x))/6};
\addlegendentry{$n=6$}
\addplot[black, densely dotted, line width=0.8pt, domain=0:6.2832, samples=300] {sin(deg(20*x))/20};
\addlegendentry{$n=20$}
\end{axis}
\end{tikzpicture}}
\caption{Four sequences of functions that converge pointwise. (A) $x^n$ on $[0,1]$, with a discontinuous limit; the convergence is uniform on $[0,a]$ for $a<1$. (B) $nxe^{-nx}$ on $[0,\infty)$: a hump of constant height $e^{-1}$ (dashed line) moves into the origin. (C) Tents of height $1$ over $[0,2/n]$, narrowing into the origin. (D) $\sin(nx)/n$, which converges uniformly to $0$ while the slopes of the graphs do not converge. In (A)--(C) the limit is not uniform on the whole interval, and in each case the failure is concentrated near a single point.}
\label{fig:basic-examples}
\end{figure}

In (i)--(iii), the failure of uniform convergence comes from a feature of the graphs, a steep rise or a hump, that does not disappear but moves into a point; away from that point the convergence is uniform. This is typical, and it suggests the following notions.

\subsection{Local Uniformity}

\begin{definition}[Locally uniform convergence]
\label{def:locally-uniform}
Let $X$ be a metric space, $S\subseteq X$, $Y$ a metric space, and $f_n,f\colon S\to Y$. The sequence $(f_n)$ converges \term{locally uniformly} to $f$ on $S$ if every point $a\in S$ has a ball $B(a,r)$ on whose intersection with $S$ the convergence is uniform. It converges \term{uniformly on compact sets} if it converges uniformly on every compact subset of $S$.
\end{definition}

\begin{proposition}
\label{prop:locally-uniform}
In the setting of \cref{def:locally-uniform}, locally uniform convergence implies uniform convergence on compact sets. If every point of $S$ has a compact neighbourhood in $S$, that is, if for every $a\in S$ some set $\bar B(a,r)\cap S$ with $r>0$ is compact, then the two notions coincide. This is the case when $S$ is an interval of $\mathbb R$ or an open subset of $\mathbb R^n$.
\end{proposition}

\begin{proof}
Let the convergence be locally uniform and $K\subseteq S$ compact. The balls $B(a,r_a)$, $a\in K$, on which the convergence is uniform cover $K$, and by \cref{prop:compact-subset} finitely many of them do; on each of these finitely many sets the convergence is uniform, hence on their union and on $K$. Conversely, if the convergence is uniform on compact sets and $\bar B(a,r)\cap S$ is compact, it is uniform on $B(a,r)\cap S$. For an interval $S$ and $a\in S$, the set $\bar B(a,r)\cap S$ is a closed bounded interval for every $r>0$, compact by \cref{thm:heine-borel}; for an open $S\subseteq\mathbb R^n$ and $\bar B(a,r)\subseteq S$, the ball itself is closed and bounded.
\end{proof}

In examples (i)--(iii), the convergence is locally uniform on $[0,1)$, on $(0,\infty)$ and on $(0,1]$, respectively, although it is not uniform there. The exponential series is another example: its partial sums $p_n(x)=\sum_{k\le n}x^k/k!$ converge to $e^x$ uniformly on every interval $[-r,r]$, since the remainder there is at most $\sum_{k>n}r^k/k!$, the remainder of the series for $e^r$; but not uniformly on $\mathbb R$, since $e^x-p_n(x)$ is unbounded for every $n$, as an exponential beats a polynomial. On $\mathbb R$, local uniformity is the natural notion for power series, and \cref{sec:power-series-revisited} shows that it holds for all of them.

\subsection{The Cauchy Criterion}

A uniform limit can be recognised without knowing it, by the criterion of Cauchy.

\begin{theorem}[Uniform Cauchy criterion]
\label{thm:uniform-cauchy}
Let $S$ be a set, $(Y,d)$ a complete metric space, and $f_n\colon S\to Y$. The sequence $(f_n)$ converges uniformly on $S$ to some function if and only if it is \term{uniformly Cauchy}: for every $\varepsilon>0$ there is $N$ with
\[
d\bigl(f_m(s),f_n(s)\bigr)<\varepsilon\qquad\text{for all }m,n\ge N\text{ and all }s\in S .
\]
\end{theorem}

\begin{proof}
If $f_n\to f$ uniformly and $N$ is such that $d(f_n(s),f(s))<\varepsilon/2$ for $n\ge N$ and all $s$, the triangle inequality gives the condition. Conversely, let the condition hold. For each $s\in S$, $(f_n(s))$ is a Cauchy sequence in $Y$, with a limit $f(s)$ because $Y$ is complete. Given $\varepsilon>0$ and $N$ as in the condition, fix $n\ge N$ and $s\in S$, and let $m\to\infty$ in $d(f_m(s),f_n(s))<\varepsilon$; since the metric is continuous (\cref{prop:distance-lipschitz}), $d(f(s),f_n(s))\le\varepsilon$. This holds for all $n\ge N$ and all $s$, so $f_n\to f$ uniformly.
\end{proof}

For bounded functions with values in $\mathbb F$ the theorem is the completeness of $\ell^\infty(S)$, \cref{thm:bounded-functions-complete}, whose proof is the same. The theorem is used in two ways. To prove uniform convergence without a candidate for the limit, as for series of functions in \cref{sec:function-series}, one estimates differences $f_m-f_n$. To disprove it, one finds $\varepsilon>0$ and, for infinitely many $n$, an index $m>n$ and a point $s$ with $d(f_m(s),f_n(s))\ge\varepsilon$.

\subsection{Deciding Uniform Convergence}

In practice the question is settled in two steps. First the pointwise limit $f$ is found. Then the supremum $\|f_n-f\|_\infty$ is computed or estimated: from above, to prove uniform convergence, and from below, to disprove it. For the estimate from below a single well-chosen point for each $n$ suffices.

\begin{proposition}[Sequential test]
\label{prop:sequential-test}
Let $S$ be a set, $Y$ a metric space, and $f_n,f\colon S\to Y$. Then $f_n\to f$ uniformly on $S$ if and only if $d(f_n(s_n),f(s_n))\to0$ for every sequence $(s_n)$ of points of $S$.
\end{proposition}

\begin{proof}
If the convergence is uniform, $d(f_n(s_n),f(s_n))\le\sup_sd(f_n(s),f(s))\to0$. If it is not, there are $\varepsilon>0$ and infinitely many indices $n$ with $\sup_sd(f_n(s),f(s))>\varepsilon$; for each of them choose $s_n$ with $d(f_n(s_n),f(s_n))>\varepsilon$, and for the other indices choose $s_n$ arbitrarily. Then $d(f_n(s_n),f(s_n))$ does not tend to $0$.
\end{proof}

The points $s_n$ should be chosen where $f_n$ is still far from $f$: $s_n=2^{-1/n}$ for $x^n$, where $|f_n(s_n)-f(s_n)|=1/2$, and $s_n=1/n$ for $nxe^{-nx}$ and for the tents. Here are three further examples, which show the method at work.

The functions $f_n(x)\coloneqq x/(1+nx^2)$ on $\mathbb R$ converge pointwise to $0$. By the inequality of the means, $1+nx^2\ge2\sqrt n\,|x|$, so $|f_n(x)|\le1/(2\sqrt n)$, with equality at $x=1/\sqrt n$; thus $\|f_n\|_\infty=1/(2\sqrt n)\to0$, and the convergence is uniform on $\mathbb R$.

The functions $f_n(x)\coloneqq\sqrt{x^2+1/n}$, $n\ge1$, converge to $|x|$ uniformly on $\mathbb R$:
\[
0\le\sqrt{x^2+1/n}-|x|=\frac{1/n}{\sqrt{x^2+1/n}+|x|}\le\frac{1/n}{\sqrt{1/n}}=\frac1{\sqrt n}.
\]
Each $f_n$ is infinitely differentiable, and the limit is not differentiable at $0$.

The functions $f_n(x)\coloneqq n^2x(1-x)^n$ on $[0,1]$ converge pointwise to $0$, since $f_n(1)=0$ and $n^2q^n\to0$ for $0\le q<1$. At $x_n=1/(n+1)$, where $f_n$ has its maximum,
\[
f_n(x_n)=\frac{n^2}{n+1}\Bigl(1-\frac1{n+1}\Bigr)^n\ge\frac{n^2}{n+1}\cdot\frac14
\]
for $n\ge1$, because $(1-1/(n+1))^n=1/(1+1/n)^n$ is greater than $1/e>1/4$ by \vcref{1}{thm:e}. So $\|f_n\|_\infty\to\infty$: the functions converge pointwise to $0$ while their maxima grow without bound.

\subsection{Operations with Uniformly Convergent Sequences}

\begin{proposition}[Operations]
\label{prop:uniform-operations}
Let $S$ be a set and let $f_n\to f$ and $g_n\to g$ uniformly on $S$.
\begin{enumerate}
\item If the functions take values in a normed space, then $f_n+g_n\to f+g$ and $cf_n\to cf$ uniformly, for every scalar $c$.
\item If the functions take values in $\mathbb F$ and $f$ and $g$ are bounded, then $f_ng_n\to fg$ uniformly.
\item If the functions take values in $\mathbb F$ and $|f(s)|\ge\mu>0$ for all $s\in S$, then $f_n(s)\ne0$ for all $s$ and all large $n$, and $1/f_n\to1/f$ uniformly.
\item If the functions take values in a metric space $Y$ and $\varphi\colon Y\to Z$ is uniformly continuous, then $\varphi\circ f_n\to\varphi\circ f$ uniformly.
\item For every map $h\colon T\to S$ from a set $T$, $f_n\circ h\to f\circ h$ uniformly on $T$.
\end{enumerate}
\end{proposition}

\begin{proof}
(i) follows from $\|(f_n+g_n)-(f+g)\|_\infty\le\|f_n-f\|_\infty+\|g_n-g\|_\infty$ and $\|cf_n-cf\|_\infty=|c|\,\|f_n-f\|_\infty$. (ii) For large $n$, $\|g_n\|_\infty\le\|g\|_\infty+1$, and
\[
|f_ng_n-fg|\le|f_n-f|\,|g_n|+|f|\,|g_n-g|\le\|f_n-f\|_\infty\bigl(\|g\|_\infty+1\bigr)+\|f\|_\infty\|g_n-g\|_\infty .
\]
(iii) For $n$ with $\|f_n-f\|_\infty\le\mu/2$, $|f_n|\ge\mu/2$ everywhere, and $|1/f_n-1/f|=|f-f_n|/(|f_n|\,|f|)\le2\|f_n-f\|_\infty/\mu^2$. (iv) Given $\varepsilon>0$, choose $\delta>0$ with $d_Z(\varphi(y),\varphi(y'))<\varepsilon$ whenever $d_Y(y,y')<\delta$, and then $N$ with $d_Y(f_n(s),f(s))<\delta$ for $n\ge N$ and all $s$. (v) The supremum over $T$ of $d(f_n(h(t)),f(h(t)))$ is at most the supremum over $S$.
\end{proof}

Boundedness in (ii) is needed. On $\mathbb R$, $x+1/n\to x$ uniformly, while $(x+1/n)^2-x^2=2x/n+1/n^2$ is unbounded for every $n$; for the same reason, uniform continuity in (iv) cannot be weakened to continuity, as $\varphi(y)=y^2$ shows. Composition on the other side, (v), needs no hypothesis, because it only relabels the points.

\subsection{Limits That Cannot Be Exchanged}

\Cref{tab:interchange-failures} collects the ways in which a property of the terms can fail to pass to the limit. In each row, the functions converge, and the property of the limit differs from the limit of the properties. The table also records whether the convergence is uniform; where it is, the failure shows that uniform convergence alone is not enough.

\begin{table}
\sdfloat{%
\begin{sdtable}{colspec={lllll}, row{1}={font=\small\bfseries}}
\toprule
property & example & for $f_n$ & for $\lim f_n$ & uniform? \\
\midrule
continuity & $x^n$ on $[0,1]$ & continuous & jumps at $1$ & no \\
integral & $nf_n$, tents of height $n$ & $\int_0^1=1$ & $\int_0^1=0$ & no \\
improper integral & $e^{-x/n}/n$ on $[0,\infty)$ & $\int_0^\infty=1$ & $\int_0^\infty=0$ & yes \\
derivative & $x/(1+nx^2)$ at $0$ & $f_n'(0)=1$ & $f'(0)=0$ & yes \\
differentiability & $\sqrt{x^2+1/n}$ & smooth & $|x|$ & yes \\
boundedness & $\min(1/x,n)$ on $(0,1]$ & bounded & $1/x$ & no \\
maxima & $n^2x(1-x)^n$ on $[0,1]$ & $\|f_n\|_\infty\to\infty$ & $0$ & no \\
arc length & $\sin(nx)/n$ on $[0,2\pi]$ & $7.6404$ & $2\pi$ & yes \\
\bottomrule
\end{sdtable}}{Properties that do not pass to the limit of a sequence of functions, with an example for each. The tents are those of \cref{fig:basic-examples}C. The arc length of the graph of $\sin(nx)/n$ over $[0,2\pi]$ is $\int_0^{2\pi}(1+\cos^2nx)^{1/2}\,dx=\int_0^{2\pi}(1+\cos^2t)^{1/2}\,dt=7.6404$ for every $n\ge1$ (\vcref{1}{thm:arc-length}), while the limit $0$ has a graph of length $2\pi=6.2832$.\label{tab:interchange-failures}}
\end{table}

The rows of the table are answered by the theorems of the next two sections: continuity and integrals pass to uniform limits, derivatives pass when the derivatives converge uniformly, improper integrals when the functions are dominated by an integrable function, and boundedness when the convergence is uniform (\cref{prop:uniform-sup-norm}). The last row shows that arc length is not a continuous function of the curve in the supremum norm; it is lower semicontinuous, since each polygonal approximation of the length is continuous, but uniform closeness of graphs does not control their lengths. The value $7.6404$ was computed by the numerical integration of the library and confirmed in ball arithmetic; it is $4\sqrt2$ times the complete elliptic integral $\int_0^{\pi/2}(1-\sin^2t/2)^{1/2}\,dt$, a special function of Volume~XI.

\section{Continuity of the Limit}
\label{sec:limit-continuity}

\subsection{The Exchange of Two Limits}

Limits of functions between metric spaces are defined as for functions of a real variable (\vcref{1}{def:function-limit}). Let $X$ and $Y$ be metric spaces, $S\subseteq X$, and $a\in X$ an accumulation point of $S$ (\cref{def:closure}). A function $f\colon S\to Y$ has the \term{limit} $L\in Y$ as $x\to a$, written $\lim_{x\to a}f(x)=L$, if for every $\varepsilon>0$ there is $\delta>0$ with $d_Y(f(x),L)<\varepsilon$ for all $x\in S$ with $0<d_X(x,a)<\delta$. If $S\subseteq\mathbb R$ is unbounded above, $\lim_{x\to\infty}f(x)=L$ means that for every $\varepsilon>0$ there is $R$ with $d_Y(f(x),L)<\varepsilon$ for all $x\in S$ with $x>R$. A limit is unique when it exists, because $a$ is an accumulation point, and a function defined at $a$ is continuous at $a$ if and only if $a$ is isolated in $S$ or $\lim_{x\to a}f(x)=f(a)$.

\begin{theorem}[Exchange of limits]
\label{thm:exchange-limits}
Let $X$ be a metric space, $S\subseteq X$, $a\in X$ an accumulation point of $S$, and $Y$ a complete metric space. Let $f_n\colon S\to Y$ converge uniformly on $S$ to $f$, and suppose that for every $n$ the limit $L_n\coloneqq\lim_{x\to a}f_n(x)$ exists. Then $(L_n)$ converges, $\lim_{x\to a}f(x)$ exists, and
\[
\lim_{x\to a}\lim_{n\to\infty}f_n(x)=\lim_{n\to\infty}\lim_{x\to a}f_n(x) .
\]
The same holds for $x\to\infty$ when $S\subseteq\mathbb R$ is unbounded above.
\end{theorem}

\begin{proof}
Let $\varepsilon>0$. By \cref{thm:uniform-cauchy} there is $N$ with $d(f_m(x),f_n(x))<\varepsilon$ for all $m,n\ge N$ and all $x\in S$. Letting $x\to a$ and using the continuity of the metric (\cref{prop:distance-lipschitz}) gives $d(L_m,L_n)\le\varepsilon$ for $m,n\ge N$. So $(L_n)$ is a Cauchy sequence, and it has a limit $L$ in the complete space $Y$. Now choose $n$ with $\sup_{x\in S}d(f_n(x),f(x))<\varepsilon$ and $d(L_n,L)<\varepsilon$, and then $\delta>0$ with $d(f_n(x),L_n)<\varepsilon$ for $x\in S$, $0<d(x,a)<\delta$. For such $x$,
\[
d\bigl(f(x),L\bigr)\le d\bigl(f(x),f_n(x)\bigr)+d\bigl(f_n(x),L_n\bigr)+d(L_n,L)<3\varepsilon .
\]
So $\lim_{x\to a}f(x)=L$. For $x\to\infty$ replace the condition $0<d(x,a)<\delta$ by $x>R$.
\end{proof}

The proof is the $\varepsilon/3$ argument of \cref{lem:uniform-limit}, with the limits $L_n$ in place of the values $f_n(a)$; completeness of $Y$ provides the limit of the $L_n$, which in the setting of the lemma is the value $f(a)$. The theorem is often called the Moore--Osgood theorem. Without uniform convergence it fails: for $x^n$ on $S=[0,1)$ and $a=1$, $\lim_{x\to1}\lim_nx^n=0$, while $\lim_n\lim_{x\to1}x^n=1$.

\begin{corollary}[Continuity of uniform limits]
\label{cor:limit-continuous}
Let $X$ and $Y$ be metric spaces, $S\subseteq X$, and $f_n\colon S\to Y$ functions that converge pointwise to $f$.
\begin{enumerate}
\item If every $f_n$ is continuous at $a\in S$ and the convergence is uniform on $B(a,r)\cap S$ for some $r>0$, then $f$ is continuous at $a$.
\item If every $f_n$ is continuous and the convergence is locally uniform, then $f$ is continuous. In particular, a limit of continuous functions on an interval of $\mathbb R$, or on an open subset of $\mathbb R^n$, that is uniform on compact sets is continuous.
\end{enumerate}
\end{corollary}

\begin{proof}
(i) The proof of \cref{lem:uniform-limit}, applied to the restrictions of $f_n$ and $f$ to the set $B(a,r)\cap S$, uses continuity of $f_n$ at the point $a$ only and shows that the restriction of $f$ is continuous at $a$; since $B(a,r)\cap S$ contains all points of $S$ close to $a$, so is $f$. (ii) follows from (i) and \cref{prop:locally-uniform}.
\end{proof}

No completeness is needed here, since the value $f(a)$ is given. The corollary explains the examples of \cref{sec:pointwise-uniform}: the limit of $x^n$ is continuous wherever the convergence is locally uniform, on $[0,1)$, and it can jump only at the point $1$, where uniformity fails. Without uniformity, a pointwise limit of continuous functions can still be continuous, as the tents show, but it need not be, and \cref{thm:baire-one} shows that it is at least continuous on a residual set when the domain is complete.

The theorem also settles double sequences. A family $(a_{mn})_{m,n\in\mathbb N}$ of points of a metric space is a \term{double sequence}, and its \term{iterated limits} are $\lim_m\lim_na_{mn}$ and $\lim_n\lim_ma_{mn}$, when the inner and outer limits exist.

\begin{corollary}[Iterated limits]
\label{cor:iterated-limits}
Let $(a_{mn})$ be a double sequence in a complete metric space $Y$. Suppose that $b_m\coloneqq\lim_na_{mn}$ exists for every $m$, and that $c_n\coloneqq\lim_ma_{mn}$ exists uniformly in $n$, that is, $\sup_nd(a_{mn},c_n)\to0$ as $m\to\infty$. Then both iterated limits exist, and $\lim_mb_m=\lim_nc_n$.
\end{corollary}

\begin{proof}
Apply \cref{thm:exchange-limits} with $S=\mathbb N\subseteq\mathbb R$, the limit $n\to\infty$, and the functions $f_m(n)\coloneqq a_{mn}$, which converge to $c$ uniformly on $\mathbb N$ as $m\to\infty$, with $\lim_{n\to\infty}f_m(n)=b_m$.
\end{proof}

For $a_{mn}\coloneqq m/(m+n)$, with $m,n\ge1$, $\lim_na_{mn}=0$ for every $m$ and $\lim_ma_{mn}=1$ for every $n$, so the iterated limits are $0$ and $1$. The convergence in $m$ is not uniform in $n$: $\sup_n|a_{mn}-1|=\sup_nn/(m+n)=1$. \Cref{sec:double-series} returns to double sequences and introduces their double limits.

\subsection{Uniform Continuity and Lipschitz Constants}

\begin{proposition}
\label{prop:limit-uniformly-continuous}
Let $X$ and $Y$ be metric spaces and $f_n\colon X\to Y$.
\begin{enumerate}
\item If every $f_n$ is uniformly continuous and $f_n\to f$ uniformly, then $f$ is uniformly continuous.
\item If every $f_n$ is Lipschitz with the constant $K$ and $f_n\to f$ pointwise, then $f$ is Lipschitz with the constant $K$.
\end{enumerate}
\end{proposition}

\begin{proof}
(i) Given $\varepsilon>0$, choose $n$ with $d(f_n(x),f(x))<\varepsilon/3$ for all $x$, and $\delta>0$ with $d(f_n(x),f_n(y))<\varepsilon/3$ whenever $d(x,y)<\delta$. For such $x,y$ the triangle inequality gives $d(f(x),f(y))<\varepsilon$. (ii) $d(f(x),f(y))=\lim_nd(f_n(x),f_n(y))\le Kd(x,y)$, by the continuity of the metric.
\end{proof}

Part (ii) needs only pointwise convergence, because a Lipschitz condition with a fixed constant is a closed condition on finitely many values at a time; uniform continuity is not, and the monomials $x^n$, each uniformly continuous on $[0,1]$, have a discontinuous pointwise limit. If the constants are not bounded, the limit need not be Lipschitz at all: $\sqrt{x+1/n}$ is Lipschitz on $[0,1]$ with the constant $\sqrt n/2$, by the mean value theorem, and it converges uniformly to $\sqrt x$, which is not Lipschitz near $0$. The functions $\sqrt{x^2+1/n}$ have the constant $1$, and so has their limit $|x|$.

\subsection{Dini's Theorem}

Pointwise convergence on a compact set becomes uniform under a monotonicity hypothesis. This is the theorem that the note on uniform convergence on compact sets in \cref{sec:compact-continuous} announced.

\begin{theorem}[Dini]
\label{thm:dini}
Let $K$ be a compact metric space and $f_n\colon K\to\mathbb R$ continuous functions with $f_0\le f_1\le f_2\le\cdots$ that converge pointwise to a continuous function $f$. Then $f_n\to f$ uniformly on $K$. The same holds for decreasing sequences.
\end{theorem}

\begin{proof}
Let $g_n\coloneqq f-f_n$. The functions $g_n$ are continuous, $g_n\ge g_{n+1}\ge0$, and $g_n(x)\to0$ for every $x$. Let $\varepsilon>0$, and let $U_n\coloneqq\{x\in K: g_n(x)<\varepsilon\}$, which is open by \cref{thm:continuity-characterisations}. Since $g_{n+1}\le g_n$, the sets increase, $U_n\subseteq U_{n+1}$, and since $g_n(x)\to0$, every $x$ lies in some $U_n$. By compactness finitely many of them cover $K$, and the one with the largest index, $U_N$, is then all of $K$. For $n\ge N$ and $x\in K$, $0\le g_n(x)\le g_N(x)<\varepsilon$. For a decreasing sequence apply this to $-f_n$.
\end{proof}

The proof is the open-cover argument in its purest form: each point $x$ has a neighbourhood on which the error is eventually below $\varepsilon$, because $g_n(x)<\varepsilon$ is an open condition, and monotonicity makes ``eventually'' last for ever. Each of the four hypotheses is needed.
\begin{enumerate}
\item Compactness: $x^n$ decreases to $0$ on $[0,1)$, a set that is not compact, and the convergence is not uniform.
\item Monotonicity: the tents of \cref{fig:basic-examples}C converge pointwise to the continuous function $0$ on the compact set $[0,1]$, not monotonically, and not uniformly.
\item Continuity of the limit: $x^n$ decreases on $[0,1]$ to a discontinuous function, and by \cref{cor:limit-continuous} the convergence cannot be uniform.
\item Continuity of the $f_n$: the indicator functions $\mathbf 1_{(0,1/n)}$, $n\ge1$, decrease to $0$ on $[0,1]$, and their supremum is $1$ for every $n$.
\end{enumerate}

A second theorem replaces monotonicity in $n$ by monotonicity in $x$, and compactness by the order structure of an interval. It is due to P\'olya~\cite{polya1920}, and Volume~IV applies it to distribution functions.

\begin{proposition}[Monotone functions with a continuous limit]
\label{prop:polya}
Let $f_n\colon[a,b]\to\mathbb R$ be increasing functions (not necessarily continuous) that converge pointwise to a continuous function $f$. Then $f_n\to f$ uniformly on $[a,b]$.
\end{proposition}

\begin{proof}
The limit $f$ is increasing, as a limit of increasing functions, and uniformly continuous by \vcref{1}{thm:heine-cantor}. Given $\varepsilon>0$, choose points $a=x_0<x_1<\cdots<x_k=b$ with $f(x_i)-f(x_{i-1})<\varepsilon$ for all $i$, and then $N$ with $|f_n(x_i)-f(x_i)|<\varepsilon$ for $n\ge N$ and $i=0,\ldots,k$. Let $n\ge N$ and $x\in[x_{i-1},x_i]$. Then, since $f_n$ and $f$ are increasing,
\[
f_n(x)\le f_n(x_i)<f(x_i)+\varepsilon<f(x_{i-1})+2\varepsilon\le f(x)+2\varepsilon ,
\]
and in the same way $f_n(x)\ge f_n(x_{i-1})>f(x_{i-1})-\varepsilon>f(x_i)-2\varepsilon\ge f(x)-2\varepsilon$. So $|f_n(x)-f(x)|<2\varepsilon$ for all $x$ and $n\ge N$.
\end{proof}

\subsection{Square Roots by Polynomials}

Dini's theorem proves that $\sqrt x$ is a uniform limit of polynomials on $[0,1]$ by an explicit iteration, and with it that $|x|$ is a uniform limit of polynomials on $[-1,1]$. Chapter~5 uses the second fact in the proof of the Stone--Weierstrass theorem.

\begin{proposition}[Polynomial approximation of the square root]
\label{prop:sqrt-polynomials}
Define polynomials by $p_0\coloneqq0$ and
\[
p_{n+1}(x)\coloneqq p_n(x)+\frac{x-p_n(x)^2}2\qquad(n\in\mathbb N).
\]
Then $0\le p_n(x)\le p_{n+1}(x)\le\sqrt x$ for all $x\in[0,1]$ and $n\in\mathbb N$, and
\[
0\le\sqrt x-p_n(x)\le\sqrt x\Bigl(1-\frac{\sqrt x}2\Bigr)^n\le\frac2{n+1}\qquad(0\le x\le1).
\]
In particular $p_n\to\sqrt x$ uniformly on $[0,1]$, and the polynomials $q_n(x)\coloneqq p_n(x^2)$ converge to $|x|$ uniformly on $[-1,1]$, with $\||x|-q_n\|_\infty\le2/(n+1)$.
\end{proposition}

\begin{proof}
Fix $x\in[0,1]$ and write $t\coloneqq\sqrt x\in[0,1]$ and $e_n\coloneqq t-p_n(x)$. The recursion gives
\[
e_{n+1}=e_n-\frac{t^2-p_n(x)^2}2=e_n\Bigl(1-\frac{t+p_n(x)}2\Bigr).
\]
Suppose that $0\le p_n(x)\le t$, which holds for $n=0$. Then $0\le(t+p_n(x))/2\le t\le1$, so $0\le e_{n+1}\le e_n(1-t/2)$; that is, $p_n(x)\le p_{n+1}(x)\le t$, and by induction these inequalities hold for all $n$, with $0\le e_n\le t(1-t/2)^n$, since $e_0=t$. This is the first bound; it shows already that $p_n(x)\to\sqrt x$ for every $x$, since $(1-t/2)^n\to0$ for $t>0$ and $e_n=0$ for $t=0$.

The convergence is uniform by Dini's theorem, since the $p_n$ increase and are continuous, and the limit is continuous. The explicit bound gives more. The function $h(t)\coloneqq t(1-t/2)^n$ on $[0,1]$ has the derivative $(1-t/2)^{n-1}\bigl(1-(n+1)t/2\bigr)$, which vanishes only at $t=2/(n+1)$; so its maximum on $[0,1]$, for $n\ge1$, is
\[
h\Bigl(\frac2{n+1}\Bigr)=\frac2{n+1}\Bigl(1-\frac1{n+1}\Bigr)^n\le\frac2{n+1},
\]
and for $n=0$ it is $1\le2$. Finally, for $x\in[-1,1]$, $|x|=\sqrt{x^2}$ with $x^2\in[0,1]$, so $0\le|x|-q_n(x)\le2/(n+1)$ by \cref{prop:uniform-operations}\,(v) and the bound.
\end{proof}

The iteration is the fixed point iteration $p\mapsto p+(x-p^2)/2$ for the equation $p^2=x$, started at $0$; at each point it converges monotonically, and Dini's theorem makes the convergence uniform. \Cref{fig:sqrt-polynomials} shows the first polynomials. Their degrees double at each step, $p_1(x)=x/2$ and $p_2(x)=x-x^2/8$, and the convergence is slowest near $0$, where the factor $1-\sqrt x/2$ is close to $1$; there the square root has its infinite slope, and no polynomial can follow it closely. \Cref{sec:uniform-computation} shows that the true error is a little more than a quarter of the bound.

\begin{figure}
\sdfloat{%
\begin{tikzpicture}
\begin{axis}[sdplot, width=0.56\textwidth, height=0.42\textwidth, xmin=0, xmax=1.02, ymin=0, ymax=1.22, xtick={0,0.25,0.5,0.75,1}, ytick={0,0.5,1}, xlabel={$x$}, sdlegendtop]
\addplot[black, line width=1.2pt, domain=0:1, samples=200] {sqrt(x)};
\addlegendentry{$\sqrt x$}
\addplot[black, dashed, line width=0.7pt] coordinates {(0.0000,0.0000) (0.0000,0.0000) (0.0000,0.0000) (0.0001,0.0000) (0.0001,0.0001) (0.0002,0.0001) (0.0002,0.0001) (0.0003,0.0002) (0.0004,0.0002) (0.0005,0.0003) (0.0006,0.0003) (0.0008,0.0004) (0.0009,0.0004) (0.0011,0.0005) (0.0012,0.0006) (0.0014,0.0007) (0.0016,0.0008) (0.0018,0.0009) (0.0020,0.0010) (0.0023,0.0011) (0.0025,0.0013) (0.0028,0.0014) (0.0030,0.0015) (0.0033,0.0017) (0.0036,0.0018) (0.0039,0.0020) (0.0042,0.0021) (0.0046,0.0023) (0.0049,0.0025) (0.0053,0.0026) (0.0056,0.0028) (0.0060,0.0030) (0.0064,0.0032) (0.0068,0.0034) (0.0072,0.0036) (0.0077,0.0038) (0.0081,0.0040) (0.0086,0.0043) (0.0090,0.0045) (0.0095,0.0048) (0.0100,0.0050) (0.0105,0.0053) (0.0110,0.0055) (0.0116,0.0058) (0.0121,0.0060) (0.0127,0.0063) (0.0132,0.0066) (0.0138,0.0069) (0.0144,0.0072) (0.0150,0.0075) (0.0156,0.0078) (0.0163,0.0081) (0.0169,0.0085) (0.0176,0.0088) (0.0182,0.0091) (0.0189,0.0095) (0.0196,0.0098) (0.0203,0.0102) (0.0210,0.0105) (0.0218,0.0109) (0.0225,0.0112) (0.0233,0.0116) (0.0240,0.0120) (0.0248,0.0124) (0.0256,0.0128) (0.0264,0.0132) (0.0272,0.0136) (0.0281,0.0140) (0.0289,0.0145) (0.0298,0.0149) (0.0306,0.0153) (0.0315,0.0158) (0.0324,0.0162) (0.0333,0.0167) (0.0342,0.0171) (0.0352,0.0176) (0.0361,0.0181) (0.0371,0.0185) (0.0380,0.0190) (0.0390,0.0195) (0.0400,0.0200) (0.0410,0.0205) (0.0420,0.0210) (0.0431,0.0215) (0.0441,0.0220) (0.0452,0.0226) (0.0462,0.0231) (0.0473,0.0237) (0.0484,0.0242) (0.0495,0.0248) (0.0506,0.0253) (0.0518,0.0259) (0.0529,0.0265) (0.0541,0.0270) (0.0552,0.0276) (0.0564,0.0282) (0.0576,0.0288) (0.0588,0.0294) (0.0600,0.0300) (0.0613,0.0306) (0.0625,0.0312) (0.0638,0.0319) (0.0650,0.0325) (0.0663,0.0332) (0.0676,0.0338) (0.0689,0.0345) (0.0702,0.0351) (0.0716,0.0358) (0.0729,0.0365) (0.0743,0.0371) (0.0756,0.0378) (0.0770,0.0385) (0.0784,0.0392) (0.0798,0.0399) (0.0812,0.0406) (0.0827,0.0413) (0.0841,0.0420) (0.0856,0.0428) (0.0870,0.0435) (0.0885,0.0443) (0.0900,0.0450) (0.0915,0.0458) (0.0930,0.0465) (0.0946,0.0473) (0.0961,0.0481) (0.0977,0.0488) (0.0992,0.0496) (0.1008,0.0504) (0.1024,0.0512) (0.1040,0.0520) (0.1056,0.0528) (0.1073,0.0536) (0.1089,0.0545) (0.1106,0.0553) (0.1122,0.0561) (0.1139,0.0570) (0.1156,0.0578) (0.1173,0.0587) (0.1190,0.0595) (0.1208,0.0604) (0.1225,0.0612) (0.1243,0.0621) (0.1260,0.0630) (0.1278,0.0639) (0.1296,0.0648) (0.1314,0.0657) (0.1332,0.0666) (0.1351,0.0675) (0.1369,0.0684) (0.1388,0.0694) (0.1406,0.0703) (0.1425,0.0713) (0.1444,0.0722) (0.1463,0.0732) (0.1482,0.0741) (0.1502,0.0751) (0.1521,0.0761) (0.1541,0.0770) (0.1560,0.0780) (0.1580,0.0790) (0.1600,0.0800) (0.1620,0.0810) (0.1640,0.0820) (0.1661,0.0830) (0.1681,0.0840) (0.1702,0.0851) (0.1722,0.0861) (0.1743,0.0872) (0.1764,0.0882) (0.1785,0.0893) (0.1806,0.0903) (0.1828,0.0914) (0.1849,0.0924) (0.1871,0.0935) (0.1892,0.0946) (0.1914,0.0957) (0.1936,0.0968) (0.1958,0.0979) (0.1980,0.0990) (0.2003,0.1001) (0.2025,0.1013) (0.2048,0.1024) (0.2070,0.1035) (0.2093,0.1047) (0.2116,0.1058) (0.2139,0.1070) (0.2162,0.1081) (0.2186,0.1093) (0.2209,0.1104) (0.2233,0.1116) (0.2256,0.1128) (0.2280,0.1140) (0.2304,0.1152) (0.2328,0.1164) (0.2352,0.1176) (0.2377,0.1188) (0.2401,0.1200) (0.2426,0.1213) (0.2450,0.1225) (0.2475,0.1238) (0.2500,0.1250) (0.2525,0.1263) (0.2550,0.1275) (0.2576,0.1288) (0.2601,0.1300) (0.2627,0.1313) (0.2652,0.1326) (0.2678,0.1339) (0.2704,0.1352) (0.2730,0.1365) (0.2756,0.1378) (0.2783,0.1391) (0.2809,0.1405) (0.2836,0.1418) (0.2862,0.1431) (0.2889,0.1445) (0.2916,0.1458) (0.2943,0.1472) (0.2970,0.1485) (0.2998,0.1499) (0.3025,0.1513) (0.3053,0.1526) (0.3080,0.1540) (0.3108,0.1554) (0.3136,0.1568) (0.3164,0.1582) (0.3192,0.1596) (0.3221,0.1610) (0.3249,0.1624) (0.3278,0.1639) (0.3306,0.1653) (0.3335,0.1668) (0.3364,0.1682) (0.3393,0.1697) (0.3422,0.1711) (0.3452,0.1726) (0.3481,0.1740) (0.3511,0.1755) (0.3540,0.1770) (0.3570,0.1785) (0.3600,0.1800) (0.3630,0.1815) (0.3660,0.1830) (0.3691,0.1845) (0.3721,0.1860) (0.3752,0.1876) (0.3782,0.1891) (0.3813,0.1907) (0.3844,0.1922) (0.3875,0.1938) (0.3906,0.1953) (0.3938,0.1969) (0.3969,0.1985) (0.4001,0.2000) (0.4032,0.2016) (0.4064,0.2032) (0.4096,0.2048) (0.4128,0.2064) (0.4160,0.2080) (0.4193,0.2096) (0.4225,0.2113) (0.4258,0.2129) (0.4290,0.2145) (0.4323,0.2162) (0.4356,0.2178) (0.4389,0.2195) (0.4422,0.2211) (0.4456,0.2228) (0.4489,0.2245) (0.4523,0.2261) (0.4556,0.2278) (0.4590,0.2295) (0.4624,0.2312) (0.4658,0.2329) (0.4692,0.2346) (0.4727,0.2363) (0.4761,0.2380) (0.4796,0.2398) (0.4830,0.2415) (0.4865,0.2433) (0.4900,0.2450) (0.4935,0.2468) (0.4970,0.2485) (0.5006,0.2503) (0.5041,0.2520) (0.5077,0.2538) (0.5112,0.2556) (0.5148,0.2574) (0.5184,0.2592) (0.5220,0.2610) (0.5256,0.2628) (0.5293,0.2646) (0.5329,0.2664) (0.5366,0.2683) (0.5402,0.2701) (0.5439,0.2720) (0.5476,0.2738) (0.5513,0.2757) (0.5550,0.2775) (0.5588,0.2794) (0.5625,0.2812) (0.5663,0.2831) (0.5700,0.2850) (0.5738,0.2869) (0.5776,0.2888) (0.5814,0.2907) (0.5852,0.2926) (0.5891,0.2945) (0.5929,0.2964) (0.5968,0.2984) (0.6006,0.3003) (0.6045,0.3023) (0.6084,0.3042) (0.6123,0.3062) (0.6162,0.3081) (0.6202,0.3101) (0.6241,0.3121) (0.6281,0.3140) (0.6320,0.3160) (0.6360,0.3180) (0.6400,0.3200) (0.6440,0.3220) (0.6480,0.3240) (0.6521,0.3260) (0.6561,0.3281) (0.6602,0.3301) (0.6642,0.3321) (0.6683,0.3342) (0.6724,0.3362) (0.6765,0.3383) (0.6806,0.3403) (0.6848,0.3424) (0.6889,0.3444) (0.6931,0.3465) (0.6972,0.3486) (0.7014,0.3507) (0.7056,0.3528) (0.7098,0.3549) (0.7140,0.3570) (0.7183,0.3591) (0.7225,0.3612) (0.7268,0.3634) (0.7310,0.3655) (0.7353,0.3677) (0.7396,0.3698) (0.7439,0.3720) (0.7482,0.3741) (0.7526,0.3763) (0.7569,0.3785) (0.7613,0.3806) (0.7656,0.3828) (0.7700,0.3850) (0.7744,0.3872) (0.7788,0.3894) (0.7832,0.3916) (0.7877,0.3938) (0.7921,0.3961) (0.7966,0.3983) (0.8010,0.4005) (0.8055,0.4028) (0.8100,0.4050) (0.8145,0.4073) (0.8190,0.4095) (0.8236,0.4118) (0.8281,0.4141) (0.8327,0.4163) (0.8372,0.4186) (0.8418,0.4209) (0.8464,0.4232) (0.8510,0.4255) (0.8556,0.4278) (0.8603,0.4301) (0.8649,0.4325) (0.8696,0.4348) (0.8742,0.4371) (0.8789,0.4395) (0.8836,0.4418) (0.8883,0.4442) (0.8930,0.4465) (0.8978,0.4489) (0.9025,0.4512) (0.9073,0.4536) (0.9120,0.4560) (0.9168,0.4584) (0.9216,0.4608) (0.9264,0.4632) (0.9312,0.4656) (0.9361,0.4680) (0.9409,0.4704) (0.9458,0.4729) (0.9506,0.4753) (0.9555,0.4778) (0.9604,0.4802) (0.9653,0.4827) (0.9702,0.4851) (0.9752,0.4876) (0.9801,0.4900) (0.9851,0.4925) (0.9900,0.4950) (0.9950,0.4975) (1.0000,0.5000)};
\addlegendentry{$p_1$}
\addplot[black, densely dotted, line width=0.8pt] coordinates {(0.0000,0.0000) (0.0000,0.0000) (0.0000,0.0000) (0.0001,0.0001) (0.0001,0.0001) (0.0002,0.0002) (0.0002,0.0002) (0.0003,0.0003) (0.0004,0.0004) (0.0005,0.0005) (0.0006,0.0006) (0.0008,0.0008) (0.0009,0.0009) (0.0011,0.0011) (0.0012,0.0012) (0.0014,0.0014) (0.0016,0.0016) (0.0018,0.0018) (0.0020,0.0020) (0.0023,0.0023) (0.0025,0.0025) (0.0028,0.0028) (0.0030,0.0030) (0.0033,0.0033) (0.0036,0.0036) (0.0039,0.0039) (0.0042,0.0042) (0.0046,0.0046) (0.0049,0.0049) (0.0053,0.0053) (0.0056,0.0056) (0.0060,0.0060) (0.0064,0.0064) (0.0068,0.0068) (0.0072,0.0072) (0.0077,0.0076) (0.0081,0.0081) (0.0086,0.0085) (0.0090,0.0090) (0.0095,0.0095) (0.0100,0.0100) (0.0105,0.0105) (0.0110,0.0110) (0.0116,0.0115) (0.0121,0.0121) (0.0127,0.0126) (0.0132,0.0132) (0.0138,0.0138) (0.0144,0.0144) (0.0150,0.0150) (0.0156,0.0156) (0.0163,0.0162) (0.0169,0.0169) (0.0176,0.0175) (0.0182,0.0182) (0.0189,0.0189) (0.0196,0.0196) (0.0203,0.0203) (0.0210,0.0210) (0.0218,0.0217) (0.0225,0.0224) (0.0233,0.0232) (0.0240,0.0240) (0.0248,0.0247) (0.0256,0.0255) (0.0264,0.0263) (0.0272,0.0271) (0.0281,0.0280) (0.0289,0.0288) (0.0298,0.0296) (0.0306,0.0305) (0.0315,0.0314) (0.0324,0.0323) (0.0333,0.0332) (0.0342,0.0341) (0.0352,0.0350) (0.0361,0.0359) (0.0371,0.0369) (0.0380,0.0378) (0.0390,0.0388) (0.0400,0.0398) (0.0410,0.0408) (0.0420,0.0418) (0.0431,0.0428) (0.0441,0.0439) (0.0452,0.0449) (0.0462,0.0460) (0.0473,0.0470) (0.0484,0.0481) (0.0495,0.0492) (0.0506,0.0503) (0.0518,0.0514) (0.0529,0.0526) (0.0541,0.0537) (0.0552,0.0548) (0.0564,0.0560) (0.0576,0.0572) (0.0588,0.0584) (0.0600,0.0596) (0.0613,0.0608) (0.0625,0.0620) (0.0638,0.0632) (0.0650,0.0645) (0.0663,0.0658) (0.0676,0.0670) (0.0689,0.0683) (0.0702,0.0696) (0.0716,0.0709) (0.0729,0.0722) (0.0743,0.0736) (0.0756,0.0749) (0.0770,0.0763) (0.0784,0.0776) (0.0798,0.0790) (0.0812,0.0804) (0.0827,0.0818) (0.0841,0.0832) (0.0856,0.0846) (0.0870,0.0861) (0.0885,0.0875) (0.0900,0.0890) (0.0915,0.0905) (0.0930,0.0919) (0.0946,0.0934) (0.0961,0.0949) (0.0977,0.0965) (0.0992,0.0980) (0.1008,0.0995) (0.1024,0.1011) (0.1040,0.1027) (0.1056,0.1042) (0.1073,0.1058) (0.1089,0.1074) (0.1106,0.1090) (0.1122,0.1107) (0.1139,0.1123) (0.1156,0.1139) (0.1173,0.1156) (0.1190,0.1173) (0.1208,0.1189) (0.1225,0.1206) (0.1243,0.1223) (0.1260,0.1240) (0.1278,0.1258) (0.1296,0.1275) (0.1314,0.1292) (0.1332,0.1310) (0.1351,0.1328) (0.1369,0.1346) (0.1388,0.1363) (0.1406,0.1382) (0.1425,0.1400) (0.1444,0.1418) (0.1463,0.1436) (0.1482,0.1455) (0.1502,0.1473) (0.1521,0.1492) (0.1541,0.1511) (0.1560,0.1530) (0.1580,0.1549) (0.1600,0.1568) (0.1620,0.1587) (0.1640,0.1607) (0.1661,0.1626) (0.1681,0.1646) (0.1702,0.1665) (0.1722,0.1685) (0.1743,0.1705) (0.1764,0.1725) (0.1785,0.1745) (0.1806,0.1765) (0.1828,0.1786) (0.1849,0.1806) (0.1871,0.1827) (0.1892,0.1847) (0.1914,0.1868) (0.1936,0.1889) (0.1958,0.1910) (0.1980,0.1931) (0.2003,0.1952) (0.2025,0.1974) (0.2048,0.1995) (0.2070,0.2017) (0.2093,0.2038) (0.2116,0.2060) (0.2139,0.2082) (0.2162,0.2104) (0.2186,0.2126) (0.2209,0.2148) (0.2233,0.2170) (0.2256,0.2193) (0.2280,0.2215) (0.2304,0.2238) (0.2328,0.2260) (0.2352,0.2283) (0.2377,0.2306) (0.2401,0.2329) (0.2426,0.2352) (0.2450,0.2375) (0.2475,0.2398) (0.2500,0.2422) (0.2525,0.2445) (0.2550,0.2469) (0.2576,0.2493) (0.2601,0.2516) (0.2627,0.2540) (0.2652,0.2564) (0.2678,0.2588) (0.2704,0.2613) (0.2730,0.2637) (0.2756,0.2661) (0.2783,0.2686) (0.2809,0.2710) (0.2836,0.2735) (0.2862,0.2760) (0.2889,0.2785) (0.2916,0.2810) (0.2943,0.2835) (0.2970,0.2860) (0.2998,0.2885) (0.3025,0.2911) (0.3053,0.2936) (0.3080,0.2962) (0.3108,0.2987) (0.3136,0.3013) (0.3164,0.3039) (0.3192,0.3065) (0.3221,0.3091) (0.3249,0.3117) (0.3278,0.3143) (0.3306,0.3170) (0.3335,0.3196) (0.3364,0.3223) (0.3393,0.3249) (0.3422,0.3276) (0.3452,0.3303) (0.3481,0.3330) (0.3511,0.3357) (0.3540,0.3384) (0.3570,0.3411) (0.3600,0.3438) (0.3630,0.3465) (0.3660,0.3493) (0.3691,0.3520) (0.3721,0.3548) (0.3752,0.3576) (0.3782,0.3603) (0.3813,0.3631) (0.3844,0.3659) (0.3875,0.3687) (0.3906,0.3716) (0.3938,0.3744) (0.3969,0.3772) (0.4001,0.3801) (0.4032,0.3829) (0.4064,0.3858) (0.4096,0.3886) (0.4128,0.3915) (0.4160,0.3944) (0.4193,0.3973) (0.4225,0.4002) (0.4258,0.4031) (0.4290,0.4060) (0.4323,0.4089) (0.4356,0.4119) (0.4389,0.4148) (0.4422,0.4178) (0.4456,0.4207) (0.4489,0.4237) (0.4523,0.4267) (0.4556,0.4297) (0.4590,0.4327) (0.4624,0.4357) (0.4658,0.4387) (0.4692,0.4417) (0.4727,0.4447) (0.4761,0.4478) (0.4796,0.4508) (0.4830,0.4539) (0.4865,0.4569) (0.4900,0.4600) (0.4935,0.4631) (0.4970,0.4661) (0.5006,0.4692) (0.5041,0.4723) (0.5077,0.4754) (0.5112,0.4786) (0.5148,0.4817) (0.5184,0.4848) (0.5220,0.4879) (0.5256,0.4911) (0.5293,0.4942) (0.5329,0.4974) (0.5366,0.5006) (0.5402,0.5037) (0.5439,0.5069) (0.5476,0.5101) (0.5513,0.5133) (0.5550,0.5165) (0.5588,0.5197) (0.5625,0.5229) (0.5663,0.5262) (0.5700,0.5294) (0.5738,0.5326) (0.5776,0.5359) (0.5814,0.5392) (0.5852,0.5424) (0.5891,0.5457) (0.5929,0.5490) (0.5968,0.5522) (0.6006,0.5555) (0.6045,0.5588) (0.6084,0.5621) (0.6123,0.5654) (0.6162,0.5688) (0.6202,0.5721) (0.6241,0.5754) (0.6281,0.5787) (0.6320,0.5821) (0.6360,0.5854) (0.6400,0.5888) (0.6440,0.5922) (0.6480,0.5955) (0.6521,0.5989) (0.6561,0.6023) (0.6602,0.6057) (0.6642,0.6091) (0.6683,0.6125) (0.6724,0.6159) (0.6765,0.6193) (0.6806,0.6227) (0.6848,0.6261) (0.6889,0.6296) (0.6931,0.6330) (0.6972,0.6365) (0.7014,0.6399) (0.7056,0.6434) (0.7098,0.6468) (0.7140,0.6503) (0.7183,0.6538) (0.7225,0.6572) (0.7268,0.6607) (0.7310,0.6642) (0.7353,0.6677) (0.7396,0.6712) (0.7439,0.6747) (0.7482,0.6782) (0.7526,0.6818) (0.7569,0.6853) (0.7613,0.6888) (0.7656,0.6924) (0.7700,0.6959) (0.7744,0.6994) (0.7788,0.7030) (0.7832,0.7065) (0.7877,0.7101) (0.7921,0.7137) (0.7966,0.7172) (0.8010,0.7208) (0.8055,0.7244) (0.8100,0.7280) (0.8145,0.7316) (0.8190,0.7352) (0.8236,0.7388) (0.8281,0.7424) (0.8327,0.7460) (0.8372,0.7496) (0.8418,0.7532) (0.8464,0.7569) (0.8510,0.7605) (0.8556,0.7641) (0.8603,0.7678) (0.8649,0.7714) (0.8696,0.7750) (0.8742,0.7787) (0.8789,0.7823) (0.8836,0.7860) (0.8883,0.7897) (0.8930,0.7933) (0.8978,0.7970) (0.9025,0.8007) (0.9073,0.8044) (0.9120,0.8081) (0.9168,0.8117) (0.9216,0.8154) (0.9264,0.8191) (0.9312,0.8228) (0.9361,0.8265) (0.9409,0.8302) (0.9458,0.8339) (0.9506,0.8377) (0.9555,0.8414) (0.9604,0.8451) (0.9653,0.8488) (0.9702,0.8526) (0.9752,0.8563) (0.9801,0.8600) (0.9851,0.8638) (0.9900,0.8675) (0.9950,0.8713) (1.0000,0.8750)};
\addlegendentry{$p_2$}
\addplot[black, dashdotted, line width=0.7pt] coordinates {(0.0000,0.0000) (0.0000,0.0000) (0.0000,0.0000) (0.0001,0.0001) (0.0001,0.0002) (0.0002,0.0003) (0.0002,0.0004) (0.0003,0.0006) (0.0004,0.0008) (0.0005,0.0010) (0.0006,0.0012) (0.0008,0.0015) (0.0009,0.0018) (0.0011,0.0021) (0.0012,0.0024) (0.0014,0.0028) (0.0016,0.0032) (0.0018,0.0036) (0.0020,0.0040) (0.0023,0.0045) (0.0025,0.0050) (0.0028,0.0055) (0.0030,0.0060) (0.0033,0.0066) (0.0036,0.0072) (0.0039,0.0078) (0.0042,0.0084) (0.0046,0.0091) (0.0049,0.0098) (0.0053,0.0105) (0.0056,0.0112) (0.0060,0.0119) (0.0064,0.0127) (0.0068,0.0135) (0.0072,0.0144) (0.0077,0.0152) (0.0081,0.0161) (0.0086,0.0170) (0.0090,0.0179) (0.0095,0.0189) (0.0100,0.0198) (0.0105,0.0208) (0.0110,0.0218) (0.0116,0.0229) (0.0121,0.0239) (0.0127,0.0250) (0.0132,0.0261) (0.0138,0.0273) (0.0144,0.0284) (0.0150,0.0296) (0.0156,0.0308) (0.0163,0.0321) (0.0169,0.0333) (0.0176,0.0346) (0.0182,0.0359) (0.0189,0.0372) (0.0196,0.0385) (0.0203,0.0399) (0.0210,0.0413) (0.0218,0.0427) (0.0225,0.0441) (0.0233,0.0456) (0.0240,0.0471) (0.0248,0.0486) (0.0256,0.0501) (0.0264,0.0516) (0.0272,0.0532) (0.0281,0.0548) (0.0289,0.0564) (0.0298,0.0580) (0.0306,0.0596) (0.0315,0.0613) (0.0324,0.0630) (0.0333,0.0647) (0.0342,0.0664) (0.0352,0.0682) (0.0361,0.0700) (0.0371,0.0718) (0.0380,0.0736) (0.0390,0.0754) (0.0400,0.0773) (0.0410,0.0791) (0.0420,0.0810) (0.0431,0.0830) (0.0441,0.0849) (0.0452,0.0868) (0.0462,0.0888) (0.0473,0.0908) (0.0484,0.0928) (0.0495,0.0949) (0.0506,0.0969) (0.0518,0.0990) (0.0529,0.1011) (0.0541,0.1032) (0.0552,0.1053) (0.0564,0.1074) (0.0576,0.1096) (0.0588,0.1118) (0.0600,0.1140) (0.0613,0.1162) (0.0625,0.1184) (0.0638,0.1207) (0.0650,0.1229) (0.0663,0.1252) (0.0676,0.1275) (0.0689,0.1298) (0.0702,0.1322) (0.0716,0.1345) (0.0729,0.1369) (0.0743,0.1393) (0.0756,0.1417) (0.0770,0.1441) (0.0784,0.1465) (0.0798,0.1490) (0.0812,0.1515) (0.0827,0.1539) (0.0841,0.1564) (0.0856,0.1589) (0.0870,0.1615) (0.0885,0.1640) (0.0900,0.1666) (0.0915,0.1691) (0.0930,0.1717) (0.0946,0.1743) (0.0961,0.1769) (0.0977,0.1796) (0.0992,0.1822) (0.1008,0.1849) (0.1024,0.1875) (0.1040,0.1902) (0.1056,0.1929) (0.1073,0.1956) (0.1089,0.1984) (0.1106,0.2011) (0.1122,0.2039) (0.1139,0.2066) (0.1156,0.2094) (0.1173,0.2122) (0.1190,0.2150) (0.1208,0.2178) (0.1225,0.2206) (0.1243,0.2234) (0.1260,0.2263) (0.1278,0.2291) (0.1296,0.2320) (0.1314,0.2349) (0.1332,0.2378) (0.1351,0.2407) (0.1369,0.2436) (0.1388,0.2465) (0.1406,0.2494) (0.1425,0.2524) (0.1444,0.2553) (0.1463,0.2583) (0.1482,0.2613) (0.1502,0.2643) (0.1521,0.2673) (0.1541,0.2703) (0.1560,0.2733) (0.1580,0.2763) (0.1600,0.2793) (0.1620,0.2823) (0.1640,0.2854) (0.1661,0.2884) (0.1681,0.2915) (0.1702,0.2946) (0.1722,0.2976) (0.1743,0.3007) (0.1764,0.3038) (0.1785,0.3069) (0.1806,0.3100) (0.1828,0.3131) (0.1849,0.3162) (0.1871,0.3194) (0.1892,0.3225) (0.1914,0.3256) (0.1936,0.3288) (0.1958,0.3319) (0.1980,0.3351) (0.2003,0.3383) (0.2025,0.3414) (0.2048,0.3446) (0.2070,0.3478) (0.2093,0.3510) (0.2116,0.3542) (0.2139,0.3574) (0.2162,0.3606) (0.2186,0.3638) (0.2209,0.3670) (0.2233,0.3702) (0.2256,0.3734) (0.2280,0.3766) (0.2304,0.3799) (0.2328,0.3831) (0.2352,0.3863) (0.2377,0.3896) (0.2401,0.3928) (0.2426,0.3960) (0.2450,0.3993) (0.2475,0.4025) (0.2500,0.4058) (0.2525,0.4090) (0.2550,0.4123) (0.2576,0.4156) (0.2601,0.4188) (0.2627,0.4221) (0.2652,0.4254) (0.2678,0.4286) (0.2704,0.4319) (0.2730,0.4352) (0.2756,0.4384) (0.2783,0.4417) (0.2809,0.4450) (0.2836,0.4483) (0.2862,0.4515) (0.2889,0.4548) (0.2916,0.4581) (0.2943,0.4614) (0.2970,0.4647) (0.2998,0.4679) (0.3025,0.4712) (0.3053,0.4745) (0.3080,0.4778) (0.3108,0.4811) (0.3136,0.4843) (0.3164,0.4876) (0.3192,0.4909) (0.3221,0.4942) (0.3249,0.4975) (0.3278,0.5007) (0.3306,0.5040) (0.3335,0.5073) (0.3364,0.5106) (0.3393,0.5139) (0.3422,0.5171) (0.3452,0.5204) (0.3481,0.5237) (0.3511,0.5269) (0.3540,0.5302) (0.3570,0.5335) (0.3600,0.5367) (0.3630,0.5400) (0.3660,0.5432) (0.3691,0.5465) (0.3721,0.5498) (0.3752,0.5530) (0.3782,0.5563) (0.3813,0.5595) (0.3844,0.5627) (0.3875,0.5660) (0.3906,0.5692) (0.3938,0.5725) (0.3969,0.5757) (0.4001,0.5789) (0.4032,0.5822) (0.4064,0.5854) (0.4096,0.5886) (0.4128,0.5918) (0.4160,0.5950) (0.4193,0.5982) (0.4225,0.6014) (0.4258,0.6046) (0.4290,0.6078) (0.4323,0.6110) (0.4356,0.6142) (0.4389,0.6174) (0.4422,0.6206) (0.4456,0.6238) (0.4489,0.6269) (0.4523,0.6301) (0.4556,0.6333) (0.4590,0.6364) (0.4624,0.6396) (0.4658,0.6427) (0.4692,0.6459) (0.4727,0.6490) (0.4761,0.6522) (0.4796,0.6553) (0.4830,0.6584) (0.4865,0.6616) (0.4900,0.6647) (0.4935,0.6678) (0.4970,0.6709) (0.5006,0.6740) (0.5041,0.6771) (0.5077,0.6802) (0.5112,0.6833) (0.5148,0.6864) (0.5184,0.6894) (0.5220,0.6925) (0.5256,0.6956) (0.5293,0.6986) (0.5329,0.7017) (0.5366,0.7048) (0.5402,0.7078) (0.5439,0.7108) (0.5476,0.7139) (0.5513,0.7169) (0.5550,0.7199) (0.5588,0.7229) (0.5625,0.7260) (0.5663,0.7290) (0.5700,0.7320) (0.5738,0.7350) (0.5776,0.7380) (0.5814,0.7409) (0.5852,0.7439) (0.5891,0.7469) (0.5929,0.7499) (0.5968,0.7528) (0.6006,0.7558) (0.6045,0.7587) (0.6084,0.7617) (0.6123,0.7646) (0.6162,0.7675) (0.6202,0.7705) (0.6241,0.7734) (0.6281,0.7763) (0.6320,0.7792) (0.6360,0.7821) (0.6400,0.7850) (0.6440,0.7879) (0.6480,0.7908) (0.6521,0.7937) (0.6561,0.7965) (0.6602,0.7994) (0.6642,0.8023) (0.6683,0.8051) (0.6724,0.8080) (0.6765,0.8108) (0.6806,0.8137) (0.6848,0.8165) (0.6889,0.8193) (0.6931,0.8221) (0.6972,0.8250) (0.7014,0.8278) (0.7056,0.8306) (0.7098,0.8334) (0.7140,0.8362) (0.7183,0.8390) (0.7225,0.8417) (0.7268,0.8445) (0.7310,0.8473) (0.7353,0.8501) (0.7396,0.8528) (0.7439,0.8556) (0.7482,0.8583) (0.7526,0.8611) (0.7569,0.8638) (0.7613,0.8666) (0.7656,0.8693) (0.7700,0.8720) (0.7744,0.8747) (0.7788,0.8774) (0.7832,0.8801) (0.7877,0.8828) (0.7921,0.8855) (0.7966,0.8882) (0.8010,0.8909) (0.8055,0.8936) (0.8100,0.8963) (0.8145,0.8990) (0.8190,0.9016) (0.8236,0.9043) (0.8281,0.9070) (0.8327,0.9096) (0.8372,0.9123) (0.8418,0.9149) (0.8464,0.9175) (0.8510,0.9202) (0.8556,0.9228) (0.8603,0.9254) (0.8649,0.9281) (0.8696,0.9307) (0.8742,0.9333) (0.8789,0.9359) (0.8836,0.9385) (0.8883,0.9411) (0.8930,0.9437) (0.8978,0.9463) (0.9025,0.9489) (0.9073,0.9515) (0.9120,0.9541) (0.9168,0.9566) (0.9216,0.9592) (0.9264,0.9618) (0.9312,0.9644) (0.9361,0.9669) (0.9409,0.9695) (0.9458,0.9720) (0.9506,0.9746) (0.9555,0.9771) (0.9604,0.9797) (0.9653,0.9822) (0.9702,0.9848) (0.9752,0.9873) (0.9801,0.9899) (0.9851,0.9924) (0.9900,0.9949) (0.9950,0.9974) (1.0000,1.0000)};
\addlegendentry{$p_4$}
\addplot[black!45, line width=0.7pt] coordinates {(0.0000,0.0000) (0.0000,0.0000) (0.0000,0.0002) (0.0001,0.0004) (0.0001,0.0008) (0.0002,0.0012) (0.0002,0.0018) (0.0003,0.0024) (0.0004,0.0032) (0.0005,0.0040) (0.0006,0.0049) (0.0008,0.0060) (0.0009,0.0071) (0.0011,0.0083) (0.0012,0.0096) (0.0014,0.0110) (0.0016,0.0124) (0.0018,0.0140) (0.0020,0.0156) (0.0023,0.0173) (0.0025,0.0191) (0.0028,0.0209) (0.0030,0.0229) (0.0033,0.0249) (0.0036,0.0269) (0.0039,0.0291) (0.0042,0.0313) (0.0046,0.0335) (0.0049,0.0358) (0.0053,0.0382) (0.0056,0.0406) (0.0060,0.0431) (0.0064,0.0456) (0.0068,0.0482) (0.0072,0.0508) (0.0077,0.0534) (0.0081,0.0561) (0.0086,0.0589) (0.0090,0.0616) (0.0095,0.0644) (0.0100,0.0672) (0.0105,0.0701) (0.0110,0.0730) (0.0116,0.0759) (0.0121,0.0788) (0.0127,0.0817) (0.0132,0.0847) (0.0138,0.0876) (0.0144,0.0906) (0.0150,0.0936) (0.0156,0.0967) (0.0163,0.0997) (0.0169,0.1027) (0.0176,0.1058) (0.0182,0.1088) (0.0189,0.1118) (0.0196,0.1149) (0.0203,0.1180) (0.0210,0.1210) (0.0218,0.1241) (0.0225,0.1271) (0.0233,0.1302) (0.0240,0.1332) (0.0248,0.1363) (0.0256,0.1393) (0.0264,0.1424) (0.0272,0.1454) (0.0281,0.1485) (0.0289,0.1515) (0.0298,0.1545) (0.0306,0.1575) (0.0315,0.1605) (0.0324,0.1635) (0.0333,0.1665) (0.0342,0.1695) (0.0352,0.1725) (0.0361,0.1754) (0.0371,0.1784) (0.0380,0.1814) (0.0390,0.1843) (0.0400,0.1872) (0.0410,0.1901) (0.0420,0.1931) (0.0431,0.1960) (0.0441,0.1988) (0.0452,0.2017) (0.0462,0.2046) (0.0473,0.2075) (0.0484,0.2103) (0.0495,0.2132) (0.0506,0.2160) (0.0518,0.2188) (0.0529,0.2216) (0.0541,0.2244) (0.0552,0.2272) (0.0564,0.2300) (0.0576,0.2328) (0.0588,0.2356) (0.0600,0.2384) (0.0613,0.2411) (0.0625,0.2439) (0.0638,0.2466) (0.0650,0.2493) (0.0663,0.2521) (0.0676,0.2548) (0.0689,0.2575) (0.0702,0.2602) (0.0716,0.2629) (0.0729,0.2656) (0.0743,0.2683) (0.0756,0.2709) (0.0770,0.2736) (0.0784,0.2763) (0.0798,0.2789) (0.0812,0.2816) (0.0827,0.2842) (0.0841,0.2869) (0.0856,0.2895) (0.0870,0.2921) (0.0885,0.2948) (0.0900,0.2974) (0.0915,0.3000) (0.0930,0.3026) (0.0946,0.3052) (0.0961,0.3078) (0.0977,0.3104) (0.0992,0.3130) (0.1008,0.3156) (0.1024,0.3182) (0.1040,0.3208) (0.1056,0.3234) (0.1073,0.3259) (0.1089,0.3285) (0.1106,0.3311) (0.1122,0.3336) (0.1139,0.3362) (0.1156,0.3388) (0.1173,0.3413) (0.1190,0.3439) (0.1208,0.3464) (0.1225,0.3490) (0.1243,0.3515) (0.1260,0.3541) (0.1278,0.3566) (0.1296,0.3592) (0.1314,0.3617) (0.1332,0.3643) (0.1351,0.3668) (0.1369,0.3693) (0.1388,0.3719) (0.1406,0.3744) (0.1425,0.3769) (0.1444,0.3795) (0.1463,0.3820) (0.1482,0.3845) (0.1502,0.3870) (0.1521,0.3896) (0.1541,0.3921) (0.1560,0.3946) (0.1580,0.3971) (0.1600,0.3997) (0.1620,0.4022) (0.1640,0.4047) (0.1661,0.4072) (0.1681,0.4097) (0.1702,0.4122) (0.1722,0.4148) (0.1743,0.4173) (0.1764,0.4198) (0.1785,0.4223) (0.1806,0.4248) (0.1828,0.4273) (0.1849,0.4298) (0.1871,0.4323) (0.1892,0.4348) (0.1914,0.4374) (0.1936,0.4399) (0.1958,0.4424) (0.1980,0.4449) (0.2003,0.4474) (0.2025,0.4499) (0.2048,0.4524) (0.2070,0.4549) (0.2093,0.4574) (0.2116,0.4599) (0.2139,0.4624) (0.2162,0.4649) (0.2186,0.4674) (0.2209,0.4699) (0.2233,0.4724) (0.2256,0.4749) (0.2280,0.4774) (0.2304,0.4800) (0.2328,0.4825) (0.2352,0.4850) (0.2377,0.4875) (0.2401,0.4900) (0.2426,0.4925) (0.2450,0.4950) (0.2475,0.4975) (0.2500,0.5000) (0.2525,0.5025) (0.2550,0.5050) (0.2576,0.5075) (0.2601,0.5100) (0.2627,0.5125) (0.2652,0.5150) (0.2678,0.5175) (0.2704,0.5200) (0.2730,0.5225) (0.2756,0.5250) (0.2783,0.5275) (0.2809,0.5300) (0.2836,0.5325) (0.2862,0.5350) (0.2889,0.5375) (0.2916,0.5400) (0.2943,0.5425) (0.2970,0.5450) (0.2998,0.5475) (0.3025,0.5500) (0.3053,0.5525) (0.3080,0.5550) (0.3108,0.5575) (0.3136,0.5600) (0.3164,0.5625) (0.3192,0.5650) (0.3221,0.5675) (0.3249,0.5700) (0.3278,0.5725) (0.3306,0.5750) (0.3335,0.5775) (0.3364,0.5800) (0.3393,0.5825) (0.3422,0.5850) (0.3452,0.5875) (0.3481,0.5900) (0.3511,0.5925) (0.3540,0.5950) (0.3570,0.5975) (0.3600,0.6000) (0.3630,0.6025) (0.3660,0.6050) (0.3691,0.6075) (0.3721,0.6100) (0.3752,0.6125) (0.3782,0.6150) (0.3813,0.6175) (0.3844,0.6200) (0.3875,0.6225) (0.3906,0.6250) (0.3938,0.6275) (0.3969,0.6300) (0.4001,0.6325) (0.4032,0.6350) (0.4064,0.6375) (0.4096,0.6400) (0.4128,0.6425) (0.4160,0.6450) (0.4193,0.6475) (0.4225,0.6500) (0.4258,0.6525) (0.4290,0.6550) (0.4323,0.6575) (0.4356,0.6600) (0.4389,0.6625) (0.4422,0.6650) (0.4456,0.6675) (0.4489,0.6700) (0.4523,0.6725) (0.4556,0.6750) (0.4590,0.6775) (0.4624,0.6800) (0.4658,0.6825) (0.4692,0.6850) (0.4727,0.6875) (0.4761,0.6900) (0.4796,0.6925) (0.4830,0.6950) (0.4865,0.6975) (0.4900,0.7000) (0.4935,0.7025) (0.4970,0.7050) (0.5006,0.7075) (0.5041,0.7100) (0.5077,0.7125) (0.5112,0.7150) (0.5148,0.7175) (0.5184,0.7200) (0.5220,0.7225) (0.5256,0.7250) (0.5293,0.7275) (0.5329,0.7300) (0.5366,0.7325) (0.5402,0.7350) (0.5439,0.7375) (0.5476,0.7400) (0.5513,0.7425) (0.5550,0.7450) (0.5588,0.7475) (0.5625,0.7500) (0.5663,0.7525) (0.5700,0.7550) (0.5738,0.7575) (0.5776,0.7600) (0.5814,0.7625) (0.5852,0.7650) (0.5891,0.7675) (0.5929,0.7700) (0.5968,0.7725) (0.6006,0.7750) (0.6045,0.7775) (0.6084,0.7800) (0.6123,0.7825) (0.6162,0.7850) (0.6202,0.7875) (0.6241,0.7900) (0.6281,0.7925) (0.6320,0.7950) (0.6360,0.7975) (0.6400,0.8000) (0.6440,0.8025) (0.6480,0.8050) (0.6521,0.8075) (0.6561,0.8100) (0.6602,0.8125) (0.6642,0.8150) (0.6683,0.8175) (0.6724,0.8200) (0.6765,0.8225) (0.6806,0.8250) (0.6848,0.8275) (0.6889,0.8300) (0.6931,0.8325) (0.6972,0.8350) (0.7014,0.8375) (0.7056,0.8400) (0.7098,0.8425) (0.7140,0.8450) (0.7183,0.8475) (0.7225,0.8500) (0.7268,0.8525) (0.7310,0.8550) (0.7353,0.8575) (0.7396,0.8600) (0.7439,0.8625) (0.7482,0.8650) (0.7526,0.8675) (0.7569,0.8700) (0.7613,0.8725) (0.7656,0.8750) (0.7700,0.8775) (0.7744,0.8800) (0.7788,0.8825) (0.7832,0.8850) (0.7877,0.8875) (0.7921,0.8900) (0.7966,0.8925) (0.8010,0.8950) (0.8055,0.8975) (0.8100,0.9000) (0.8145,0.9025) (0.8190,0.9050) (0.8236,0.9075) (0.8281,0.9100) (0.8327,0.9125) (0.8372,0.9150) (0.8418,0.9175) (0.8464,0.9200) (0.8510,0.9225) (0.8556,0.9250) (0.8603,0.9275) (0.8649,0.9300) (0.8696,0.9325) (0.8742,0.9350) (0.8789,0.9375) (0.8836,0.9400) (0.8883,0.9425) (0.8930,0.9450) (0.8978,0.9475) (0.9025,0.9500) (0.9073,0.9525) (0.9120,0.9550) (0.9168,0.9575) (0.9216,0.9600) (0.9264,0.9625) (0.9312,0.9650) (0.9361,0.9675) (0.9409,0.9700) (0.9458,0.9725) (0.9506,0.9750) (0.9555,0.9775) (0.9604,0.9800) (0.9653,0.9825) (0.9702,0.9850) (0.9752,0.9875) (0.9801,0.9900) (0.9851,0.9925) (0.9900,0.9950) (0.9950,0.9975) (1.0000,1.0000)};
\addlegendentry{$p_{16}$}
\end{axis}
\end{tikzpicture}}{The polynomials $p_1$, $p_2$, $p_4$ and $p_{16}$ of \cref{prop:sqrt-polynomials} (of degrees $1$, $2$, $8$ and $2^{15}$) increase towards $\sqrt x$ (thick). The error is largest near $0$, where the square root rises steeply.\label{fig:sqrt-polynomials}}
\end{figure}

\section{Integration and Differentiation of the Limit}
\label{sec:limit-integration}

\subsection{Integrals of Uniform Limits}

\begin{theorem}[Integration of uniform limits]
\label{thm:uniform-integral}
Let $f_n\colon[a,b]\to\mathbb R$ be integrable functions that converge uniformly to $f$ on $[a,b]$. Then $f$ is integrable,
\[
\Bigl|\int_a^bf_n-\int_a^bf\Bigr|\le(b-a)\,\|f_n-f\|_\infty\to0 ,
\]
and the functions $F_n(x)\coloneqq\int_a^xf_n$ converge uniformly on $[a,b]$ to $F(x)\coloneqq\int_a^xf$.
\end{theorem}

\begin{proof}
Let $\varepsilon_n\coloneqq\|f_n-f\|_\infty$, which is finite for large $n$ and tends to $0$. Then $f$ is bounded, since $|f|\le|f_n|+\varepsilon_n$. Let $P$ be a partition of $[a,b]$, and on its $k$th subinterval let $M_k$ and $m_k$ be the supremum and the infimum of $f$, and $M_k^{(n)}$ and $m_k^{(n)}$ those of $f_n$. From $f_n-\varepsilon_n\le f\le f_n+\varepsilon_n$ we get $M_k\le M_k^{(n)}+\varepsilon_n$ and $m_k\ge m_k^{(n)}-\varepsilon_n$, and so, in the notation of \vcref{1}{def:darboux-sums},
\[
U(f,P)-L(f,P)\le U(f_n,P)-L(f_n,P)+2\varepsilon_n(b-a) .
\]
Given $\varepsilon>0$, choose $n$ with $2\varepsilon_n(b-a)<\varepsilon/2$, and then, by Riemann's criterion (\vcref{1}{thm:riemann-criterion}) for the integrable $f_n$, a partition $P$ with $U(f_n,P)-L(f_n,P)<\varepsilon/2$. Then $U(f,P)-L(f,P)<\varepsilon$, and $f$ is integrable by the same criterion. By the linearity and the monotonicity of the integral (\vcref{1}{thm:integral-linear,thm:integral-monotone}),
\[
\Bigl|\int_a^xf_n-\int_a^xf\Bigr|\le\int_a^x|f_n-f|\le(x-a)\,\varepsilon_n\le(b-a)\,\varepsilon_n
\]
for every $x\in[a,b]$, which proves both remaining statements.
\end{proof}

For complex-valued functions, integrated through their real and imaginary parts, the theorem holds with the same bound, applied to the two parts separately with the constant $2(b-a)$, or with $b-a$ when the inequality $|\int u|\le\int|u|$ of \vcref{2}{lem:complex-calculus} is available. In the language of \cref{ch:metric-spaces}, the theorem says that the integrable functions form a closed subspace of $\ell^\infty([a,b])$, and that the integral is a bounded linear functional on it, with norm $b-a$; on $C[a,b]$ this is the case $K=1$ of \cref{prop:integral-operator-norm}.

\begin{corollary}[Integration term by term]
\label{cor:termwise-integration}
If $f_k\colon[a,b]\to\mathbb R$ are integrable and the series $\sum_kf_k$ converges uniformly on $[a,b]$, then its sum is integrable and $\int_a^b\sum_kf_k=\sum_k\int_a^bf_k$, the series on the right being convergent.
\end{corollary}

\begin{proof}
Apply \cref{thm:uniform-integral} to the partial sums, whose integrals are the partial sums of $\sum_k\int_a^bf_k$ by the linearity of the integral.
\end{proof}

Uniform convergence is sufficient for the exchange, not necessary. The monomials $x^n$ have integrals $1/(n+1)$ over $[0,1]$, which tend to $0$, the integral of their pointwise limit; the tents of \cref{fig:basic-examples}C also have integrals tending to the integral $0$ of their limit. The taller tents $nf_n$ do not: their integrals are all $1$. What distinguishes the two cases is boundedness. The monomials and the low tents are bounded by $1$, while the tall tents have the heights $n$, and the following theorem of Arzel\`a~\cite{arzela1885} shows that boundedness is all that is needed, as long as the limit is integrable.

\begin{theorem}[Bounded convergence]
\label{thm:bounded-convergence}
Let $f_n\colon[a,b]\to\mathbb R$ be integrable functions with $|f_n(x)|\le M$ for all $n$ and $x$, and suppose that $f_n\to f$ pointwise, where $f$ is integrable. Then $\int_a^bf_n\to\int_a^bf$.
\end{theorem}

\begin{proof}[sketch]
Replacing $f_n$ by $|f_n-f|$, which is integrable and bounded by $2M$, reduces the statement to nonnegative functions $g_n\to0$, for which $\int_a^bg_n\to0$ is to be shown. In Part~III, every Riemann integrable function is Lebesgue integrable with the same integral, and the statement becomes a special case of the dominated convergence theorem of Chapter~18, with the constant $2M$ as the dominating function. Proofs within the Riemann integral exist, but they are longer: they replace $\sup_{k\ge n}g_k$, which need not be integrable, by functions below it whose integrals can be controlled, and conclude with a compactness argument of the kind used for Dini's theorem. Luxemburg~\cite{luxemburg1971} and Lewin~\cite{lewin1986} give such proofs.
\end{proof}

The hypothesis that $f$ is integrable cannot be dropped. Enumerate the rational numbers of $[0,1]$ as $r_0,r_1,\ldots$, and let $f_n$ be the indicator function of $\{r_0,\ldots,r_n\}$. Each $f_n$ vanishes except at finitely many points, so it is integrable with integral $0$ (\vcref{1}{cor:change-finitely-many}), the $f_n$ are bounded by $1$, and they increase pointwise to the Dirichlet function $\mathbf 1_{\mathbb Q}$, which is not integrable (\vcref{1}{prop:dirichlet}). The Riemann integral is not closed under bounded pointwise limits, not even under increasing ones, and this is one of the defects that the Lebesgue integral removes.

\subsection{Improper Integrals}

On an unbounded interval, uniform convergence is not enough. The functions $f_n(x)\coloneqq e^{-x/n}/n$, $n\ge1$, converge to $0$ uniformly on $[0,\infty)$, since $\|f_n\|_\infty=1/n$, but $\int_0^\infty f_n=1$ for every $n$: a fixed amount of area is spread ever more thinly over a growing interval. No integrable function bounds all of them, since for $x\ge1$ and $n=\lceil x\rceil$ the value $f_n(x)\ge e^{-1}/(x+1)$, and $1/(x+1)$ is not integrable on $[0,\infty)$. A dominating function is exactly what is needed.

\begin{theorem}[Dominated limits of improper integrals]
\label{thm:improper-dominated}
Let $f_n\colon[a,\infty)\to\mathbb R$ be integrable on every interval $[a,R]$, and suppose that $f_n\to f$ uniformly on every $[a,R]$, and that there is a function $g$, integrable on every $[a,R]$ and with $\int_a^\infty g$ convergent, such that $|f_n(x)|\le g(x)$ for all $n$ and $x\ge a$. Then the improper integrals $\int_a^\infty f_n$ and $\int_a^\infty f$ converge absolutely, and
\[
\lim_{n\to\infty}\int_a^\infty f_n=\int_a^\infty f .
\]
\end{theorem}

\begin{proof}
By \cref{thm:uniform-integral}, $f$ is integrable on every $[a,R]$, and $|f|\le g$ as the pointwise limit of functions with $|f_n|\le g$. So $|f_n|$ and $|f|$ are integrable on every $[a,R]$ (\vcref{1}{thm:integrable-algebra}) and bounded by $g$, and the integrals converge absolutely by the comparison test and \vcref{1}{thm:comparison-test,thm:absolute-convergence}. Let $\varepsilon>0$, and choose $R$ with $\int_R^\infty g<\varepsilon$. Then
\[
\Bigl|\int_a^\infty f_n-\int_a^\infty f\Bigr|\le\int_a^R|f_n-f|+\int_R^\infty|f_n|+\int_R^\infty|f|\le(R-a)\sup_{[a,R]}|f_n-f|+2\varepsilon ,
\]
and the first term tends to $0$ as $n\to\infty$.
\end{proof}

The theorem is the sequential counterpart of \vcref{1}{thm:parameter-improper}\,(i), where a parameter $t$ took the place of the index $n$ and the uniform convergence of the integral was obtained from a dominating function by \vcref{1}{prop:weierstrass-test}. It holds, with the same proof, for integrals that are improper at a finite end. In Part~III, the dominated convergence theorem of Chapter~18 needs only pointwise convergence. A classical application is a formula of Gauss for the gamma function of \vcref{1}{def:gamma}, which the note on the gamma function in \vcref{1}{sec:general-powers} mentions and which Chapter~15 extends to complex arguments.

\begin{proposition}[Gauss's formula]
\label{prop:gauss-gamma}
For every $s>0$,
\[
\Gamma(s)=\lim_{n\to\infty}\frac{n!\,n^s}{s(s+1)\cdots(s+n)} .
\]
\end{proposition}

\begin{proof}
Let first $s\ge1$, so that $x^{s-1}$ is continuous on $[0,\infty)$, and let $f_n(x)\coloneqq(1-x/n)^nx^{s-1}$ for $0\le x\le n$ and $f_n(x)\coloneqq0$ for $x>n$, a continuous function. By \vcref{1}{prop:ln-inequalities}, $n\ln(1-x/n)\le-x$ for $0\le x<n$, so $0\le f_n(x)\le g(x)\coloneqq x^{s-1}e^{-x}$, whose integral over $[0,\infty)$ is $\Gamma(s)$ (\vcref{1}{thm:gamma}). For $0\le t\le1/2$,
\[
0\le-\ln(1-t)-t=\sum_{k\ge2}\frac{t^k}k\le\frac{t^2}2\sum_{j\ge0}t^j\le t^2
\]
by \vcref{1}{thm:log-series,thm:ln-lambda}, so for $R\ge1$, $n\ge2R$ and $0\le x\le R$ we have $-x-x^2/n\le n\ln(1-x/n)\le-x$, and since $1-e^{-y}\le y$ for $y\ge0$,
\[
0\le g(x)-f_n(x)=x^{s-1}e^{-x}\bigl(1-e^{n\ln(1-x/n)+x}\bigr)\le x^{s-1}e^{-x}\frac{x^2}n\le\frac{R^{s+1}}n .
\]
So $f_n\to g$ uniformly on $[0,R]$, and \cref{thm:improper-dominated} gives $\int_0^nf_n\to\Gamma(s)$. With $x=nu$,
\[
\int_0^n\Bigl(1-\frac xn\Bigr)^nx^{s-1}\,dx=n^s\,I(n,s),\qquad I(n,s)\coloneqq\int_0^1(1-u)^nu^{s-1}\,du .
\]
Integration by parts (\vcref{1}{thm:parts}) with $u^{s-1}=(u^s/s)'$ gives $I(n,s)=(n/s)\,I(n-1,s+1)$ for $n\ge1$, the boundary terms vanishing, and $I(0,s+n)=1/(s+n)$; hence $I(n,s)=n!/(s(s+1)\cdots(s+n))$, which proves the formula for $s\ge1$. For $0<s<1$, write $G_n(s)$ for the quotient in the formula. Then $G_n(s)=G_n(s+1)\,(s+n+1)/(sn)$, and since $G_n(s+1)\to\Gamma(s+1)=s\,\Gamma(s)$ (\vcref{1}{thm:gamma}) and $(s+n+1)/n\to1$, $G_n(s)\to\Gamma(s)$.
\end{proof}

\subsection{Derivatives of Uniform Limits}

Differentiation is where uniform convergence of the functions is least useful. The functions $\sqrt{x^2+1/n}$ are smooth and converge uniformly to $|x|$, which is not differentiable at $0$; the functions $x/(1+nx^2)$ converge uniformly to $0$ on $\mathbb R$, but their derivatives at $0$ are all equal to $1$; and $\sin(nx)/n$ converges uniformly to $0$ while its derivatives do not converge anywhere except on $2\pi\mathbb Z$. A derivative is a limit of difference quotients, $(f_n(x+h)-f_n(x))/h$, in which the uniformly small error $f_n-f$ is divided by the small number $h$, and nothing bounds the quotient. The right hypothesis concerns the derivatives themselves.

\begin{theorem}[Differentiation of limits]
\label{thm:uniform-derivative}
Let $I$ be a bounded interval and $f_n\colon I\to\mathbb F$ differentiable functions (one-sidedly at an end that belongs to $I$). Suppose that the derivatives $f_n'$ converge uniformly on $I$ to a function $g$, and that $(f_n(x_0))$ converges for some $x_0\in I$. Then $(f_n)$ converges uniformly on $I$ to a function $f$, $f$ is differentiable on $I$, and $f'=g$; that is, $(\lim_nf_n)'=\lim_nf_n'$.
\end{theorem}

\begin{proof}
Suppose first that the functions are real-valued. For $m,n\in\mathbb N$, the mean value theorem (\vcref{1}{thm:mvt}), applied to $h\coloneqq f_n-f_m$ on the interval between two points $s,t\in I$, gives
\begin{equation}
\bigl|(f_n-f_m)(s)-(f_n-f_m)(t)\bigr|\le|s-t|\,\|f_n'-f_m'\|_\infty .
\label{eq:mvt-differences}
\end{equation}
With $t=x_0$ and $\ell$ the length of $I$, $|f_n(s)-f_m(s)|\le|f_n(x_0)-f_m(x_0)|+\ell\,\|f_n'-f_m'\|_\infty$ for all $s\in I$. The right side is small for large $m$ and $n$, since $(f_n(x_0))$ converges and $(f_n')$ is uniformly Cauchy by \cref{thm:uniform-cauchy}. So $(f_n)$ is uniformly Cauchy, and it converges uniformly to some $f$.

Fix $x\in I$ and let $S\coloneqq I\setminus\{x\}$, of which $x$ is an accumulation point. On $S$ define the difference quotients
\[
\varphi_n(t)\coloneqq\frac{f_n(t)-f_n(x)}{t-x},\qquad\varphi(t)\coloneqq\frac{f(t)-f(x)}{t-x} .
\]
By \eqref{eq:mvt-differences} with $s=x$, $|\varphi_n(t)-\varphi_m(t)|\le\|f_n'-f_m'\|_\infty$ for all $t\in S$, so $(\varphi_n)$ is uniformly Cauchy on $S$; its pointwise limit is $\varphi$, so $\varphi_n\to\varphi$ uniformly on $S$. Moreover $\lim_{t\to x}\varphi_n(t)=f_n'(x)$. By \cref{thm:exchange-limits}, $\lim_{t\to x}\varphi(t)$ exists and equals $\lim_nf_n'(x)=g(x)$; that is, $f'(x)=g(x)$.

For complex-valued $f_n$, the real parts satisfy the hypotheses, since $(\operatorname{Re}f_n)'=\operatorname{Re}f_n'$ and $|\operatorname{Re}z|\le|z|$, and so do the imaginary parts; the conclusion for $f_n$ follows from those for the two parts.
\end{proof}

The derivatives $f_n'$ are not assumed to be continuous, and the limit $g$ need not be continuous either; the proof uses only the mean value theorem. On an unbounded interval, or on an open interval with uniform convergence of the derivatives only on compact subintervals, the theorem applies on every bounded subinterval containing $x_0$: then $f_n\to f$ uniformly on bounded subintervals, and $f'=g$ everywhere. The condition at a single point $x_0$ cannot be dropped, as $f_n=n$ shows, but it only fixes the constants of integration.

For continuously differentiable functions the theorem has a shorter proof, which deserves to be seen. If $f_n\in C^1[a,b]$, $f_n'\to g$ uniformly and $f_n(a)\to c$, then $g$ is continuous by \cref{lem:uniform-limit}, and $f_n(x)=f_n(a)+\int_a^xf_n'$ by the fundamental theorem (\vcref{1}{thm:ftc2}). By \cref{thm:uniform-integral} the right side converges uniformly to $c+\int_a^xg$, which is differentiable with derivative $g$ by \vcref{1}{thm:ftc1}. So the limit $f(x)=c+\int_a^xg$ belongs to $C^1[a,b]$, and $f'=g$.

\begin{corollary}[Limits in $C^k$]
\label{cor:ck-limits}
Let $k\ge1$ and $f_n\in C^k[a,b]$. If the derivatives $f_n^{(j)}$ converge uniformly on $[a,b]$ to functions $g_j$, for $j=0,1,\ldots,k$, then $g_0\in C^k[a,b]$ and $g_0^{(j)}=g_j$. With the norm
\[
\|f\|_{C^k}\coloneqq\sum_{j=0}^k\|f^{(j)}\|_\infty ,
\]
$C^k[a,b]$ is a Banach space.
\end{corollary}

\begin{proof}
The first statement follows by induction on $j$ from \cref{thm:uniform-derivative} and \cref{lem:uniform-limit}: if $g_0\in C^{j}[a,b]$ with $g_0^{(j)}=g_j$, apply the theorem to the functions $f_n^{(j)}$, whose derivatives converge uniformly to $g_{j+1}$, to get $g_j'=g_{j+1}$, which is continuous. The space $C^k[a,b]$ is a vector space, and $\|\cdot\|_{C^k}$ is a norm on it. A Cauchy sequence for this norm has each derivative $f_n^{(j)}$ uniformly Cauchy, hence uniformly convergent to some $g_j$ by \cref{thm:uniform-cauchy}; by the first statement $g_0\in C^k[a,b]$ and $\|f_n-g_0\|_{C^k}=\sum_j\|f_n^{(j)}-g_j\|_\infty\to0$.
\end{proof}

For $k=1$ this is the norm of \cref{prop:differentiation-unbounded}\,(ii), in which differentiation $D\colon C^1[0,1]\to C[0,1]$ is bounded. In the supremum norm alone, $D$ is unbounded, by part~(i) of that proposition, and $(C^1[a,b],\|\cdot\|_\infty)$ is not complete: the functions $\sqrt{x^2+1/n}$ on $[-1,1]$ form a Cauchy sequence for $\|\cdot\|_\infty$ whose uniform limit $|x|$ is not differentiable. What survives of $D$ in the supremum norm is the statement of the theorem, that if $f_n\to f$ and $Df_n\to g$ uniformly, then $f$ is differentiable and $Df=g$. In the language of operator theory, $D$ is an unbounded but \term{closed} operator: its graph $\{(f,f'): f\in C^1[a,b]\}$ is closed in $C[a,b]\times C[a,b]$. This is the setting in which differential operators are studied in functional analysis, as the Notes explain.

\section{Series of Functions}
\label{sec:function-series}

\subsection{Uniform Convergence of Series}

A series $\sum_kf_k$ of functions $f_k\colon S\to Y$ with values in a normed space \term{converges pointwise}, \term{uniformly} or \term{locally uniformly} if the sequence of its partial sums $s_n\coloneqq\sum_{k=0}^nf_k$ does; it \term{converges absolutely} on $S$ if $\sum_k\|f_k(s)\|<\infty$ for every $s\in S$. The theorems of the preceding sections apply to the partial sums. By \cref{thm:uniform-cauchy}, a series of functions with values in a Banach space converges uniformly on $S$ if and only if
\begin{equation}
\sup_{s\in S}\Bigl\|\sum_{k=m}^nf_k(s)\Bigr\|\to0\qquad(n\ge m\to\infty),
\label{eq:series-uniform-cauchy}
\end{equation}
that is, if for every $\varepsilon>0$ there is $N$ such that the norms of all blocks $\sum_{k=m}^nf_k(s)$ with $n\ge m\ge N$ are less than $\varepsilon$, at all points $s$ at once. Taking $m=n$ shows that the terms of a uniformly convergent series tend to $0$ uniformly, $\|f_n\|_\infty\to0$, the uniform form of the term test. The remainder $r_n\coloneqq\sum_{k>n}f_k$ of a uniformly convergent series tends to $0$ uniformly, and conversely, since $s_n-s=-r_n$.

\begin{theorem}[Weierstrass M-test]
\label{thm:m-test}
Let $S$ be a set, $Y$ a Banach space, and $f_k\colon S\to Y$ functions with $\|f_k(s)\|\le M_k$ for all $s\in S$, where $\sum_kM_k<\infty$. Then $\sum_kf_k$ converges absolutely and uniformly on $S$, and its remainders satisfy
\[
\Bigl\|\sum_{k>n}f_k(s)\Bigr\|\le\sum_{k>n}M_k\qquad\text{for all }s\in S .
\]
\end{theorem}

\begin{proof}
For $n\ge m$ and $s\in S$, $\|\sum_{k=m}^nf_k(s)\|\le\sum_{k=m}^nM_k$, which is small for large $m$ by the Cauchy criterion for the convergent series $\sum M_k$ (\vcref{1}{thm:cauchy-series}). So \eqref{eq:series-uniform-cauchy} holds, and the series converges uniformly. Absolute convergence at each point follows by comparison with $\sum M_k$, and the bound for the remainder from \cref{thm:absolute-series}, applied in $Y$ to the series $\sum_{k>n}f_k(s)$.
\end{proof}

Two earlier results are special cases. For bounded continuous functions with values in $\mathbb F$ the theorem is the statement after \cref{thm:absolute-series}, which takes $M_k=\|f_k\|_\infty$ and works in the Banach space $C_b(X)$. For continuous functions on a closed bounded interval with values in $\mathbb C^n$ and the norm $\|\cdot\|_\infty$, it is part (ii) of \vcref{2}{lem:vector-functions}, together with the continuity of the sum, which now follows from \cref{lem:uniform-limit}, and the integration term by term, which follows from \cref{cor:termwise-integration} applied to each component. In every case the best choice is $M_k=\|f_k\|_\infty$, and the test then says: if $\sum_k\|f_k\|_\infty<\infty$, the series converges uniformly. Such series are sometimes called \term{normally convergent}.

Here are some examples. The series $\sum_{k\ge1}\sin(kx)/k^2$ and $\sum_{k\ge1}\cos(kx)/k^2$ converge uniformly on $\mathbb R$, with $M_k=1/k^2$, and their sums are continuous. The exponential series converges uniformly on every interval $[-r,r]$, with $M_k=r^k/k!$, as observed in \cref{sec:pointwise-uniform}. The series $\sum_{k\ge1}x^k/k^2$ converges uniformly on $[-1,1]$, with $M_k=1/k^2$.

Not every uniformly convergent series can be detected by the M-test. For example, the series $\sum_{k\ge1}(-1)^k/(k+x)$ converges uniformly on $[0,\infty)$: for each $x$ it is an alternating series with decreasing terms, and by \vcref{1}{thm:leibniz} its remainder after the term with index $n$ is at most the next term, $1/(n+1+x)\le1/(n+1)$, for every $x\ge0$. But $\|(-1)^k/(k+x)\|_\infty=1/k$, and $\sum1/k$ diverges. Conversely, absolute convergence at every point does not imply uniform convergence. The terms of
\[
\sum_{k=0}^\infty\frac{x^2}{(1+x^2)^k}
\]
are nonnegative, and for $x\ne0$ the series is geometric with the ratio $1/(1+x^2)<1$ and the sum $x^2/(1-1/(1+x^2))=1+x^2$, while for $x=0$ every term vanishes. The sum is discontinuous at $0$, and by \cref{cor:limit-continuous} the convergence cannot be uniform near $0$, although it is absolute everywhere.

\subsection{The Tests of Abel and Dirichlet}

For series that converge only conditionally, the tests of Abel and Dirichlet, which in their numerical form prove the convergence of series such as $\sum_k\sin(k)/k$, have uniform versions. Both rest on summation by parts, the discrete form of integration by parts of \vcref{2}{prop:summation-by-parts}. With $A_k\coloneqq a_0+\cdots+a_k$ and $A_{-1}\coloneqq0$, so that $a_k=A_k-A_{k-1}$, a rearrangement of the terms gives, for $0\le m\le n$,
\begin{equation}
\sum_{k=m}^na_kb_k=A_nb_n-A_{m-1}b_m+\sum_{k=m}^{n-1}A_k(b_k-b_{k+1}) .
\label{eq:summation-by-parts}
\end{equation}
The identity also holds with $A_k$ replaced by $A_k-c$ for any constant $c$, since the added terms $-cb_n+cb_m-c\sum_{k=m}^{n-1}(b_k-b_{k+1})$ cancel.

\begin{theorem}[Abel and Dirichlet tests]
\label{thm:abel-dirichlet}
Let $S$ be a set, and $a_k\colon S\to\mathbb F$ and $b_k\colon S\to\mathbb R$ functions such that, for every $s\in S$, the sequence $(b_k(s))_k$ is monotone.
\begin{enumerate}
\item (Dirichlet) If the partial sums $A_n(s)=\sum_{k\le n}a_k(s)$ satisfy $|A_n(s)|\le M$ for all $n$ and $s$, and $b_k\to0$ uniformly on $S$, then $\sum_ka_kb_k$ converges uniformly on $S$.
\item (Abel) If $\sum_ka_k$ converges uniformly on $S$ and $|b_k(s)|\le M$ for all $k$ and $s$, then $\sum_ka_kb_k$ converges uniformly on $S$.
\end{enumerate}
\end{theorem}

\begin{proof}
Fix $s$ and omit it from the notation. Since $(b_k)$ is monotone, the differences $b_k-b_{k+1}$ all have the same sign, and $\sum_{k=m}^{n-1}|b_k-b_{k+1}|=|b_m-b_n|\le|b_m|+|b_n|$. (i) By \eqref{eq:summation-by-parts},
\[
\Bigl|\sum_{k=m}^na_kb_k\Bigr|\le M|b_n|+M|b_m|+M\bigl(|b_m|+|b_n|\bigr)\le4M\sup_{k\ge m}\|b_k\|_\infty ,
\]
which tends to $0$ as $m\to\infty$ independently of $s$; so \eqref{eq:series-uniform-cauchy} holds. (ii) Let $A(s)\coloneqq\sum_ka_k(s)$ and $\rho_k\coloneqq A_k-A$, which tends to $0$ uniformly. Given $\varepsilon>0$, choose $N$ with $|\rho_k(s)|<\varepsilon$ for $k\ge N$ and all $s$. For $n\ge m>N$, \eqref{eq:summation-by-parts} with $\rho_k$ in place of $A_k$ gives
\[
\Bigl|\sum_{k=m}^na_kb_k\Bigr|\le\varepsilon|b_n|+\varepsilon|b_m|+\varepsilon\bigl(|b_m|+|b_n|\bigr)\le4M\varepsilon ,
\]
for every $s$, and \eqref{eq:series-uniform-cauchy} holds again.
\end{proof}

With constant functions, the theorem contains the ordinary tests: Dirichlet's, that $\sum a_kb_k$ converges when the partial sums of $\sum a_k$ are bounded and $b_k$ decreases to $0$, and Abel's, that it converges when $\sum a_k$ converges and $(b_k)$ is monotone and bounded. Leibniz's test for alternating series, \vcref{1}{thm:leibniz}, is Dirichlet's test with $a_k=(-1)^k$.

\subsection{A Sawtooth}

The series of Abel's counterexample is closely related to
\begin{equation}
g(x)\coloneqq\sum_{k=1}^\infty\frac{\sin kx}k ,
\label{eq:sawtooth}
\end{equation}
which shows all the phenomena of this section. The identity $2\sin(x/2)\sin(kx)=\cos((k-1/2)x)-\cos((k+1/2)x)$ makes the sum of the sines telescope: for $x\notin2\pi\mathbb Z$,
\[
\sum_{k=1}^n\sin kx=\frac{\cos(x/2)-\cos((n+1/2)x)}{2\sin(x/2)},\qquad\text{so}\qquad\Bigl|\sum_{k=1}^n\sin kx\Bigr|\le\frac1{|\sin(x/2)|} .
\]
On an interval $[\delta,2\pi-\delta]$ with $0<\delta<\pi$, the bound is at most $1/\sin(\delta/2)$, and Dirichlet's test with $b_k=1/k$ shows that the series converges uniformly there. By periodicity it converges uniformly on every closed set that keeps a positive distance from $2\pi\mathbb Z$, and at the points of $2\pi\mathbb Z$ all terms vanish. So the series converges at every $x\in\mathbb R$, and its sum $g$ is continuous on $\mathbb R\setminus2\pi\mathbb Z$, by \cref{cor:limit-continuous}.

Near $0$ the convergence is not uniform. For $n\ge1$ and $x_n\coloneqq\pi/(4n)$, every $k$ with $n<k\le2n$ has $kx_n\in(\pi/4,\pi/2]$, so $\sin(kx_n)\ge\sin(\pi/4)=\sqrt2/2$, and
\[
\sum_{k=n+1}^{2n}\frac{\sin kx_n}k\ge n\cdot\frac{\sqrt2/2}{2n}=\frac{\sqrt2}4 .
\]
The blocks from $n+1$ to $2n$ do not become uniformly small on any interval $(0,\delta)$, which contains $x_n$ for large $n$, and \eqref{eq:series-uniform-cauchy} fails there. Chapter~9 shows that
\[
g(x)=\frac{\pi-x}2\qquad(0<x<2\pi),
\]
so that $g$ is the sawtooth of \cref{fig:sawtooth}, which jumps by $\pi$ at every multiple of $2\pi$, where it has the value $0$. A uniformly convergent series of continuous functions could not have this sum, and the failure of uniform convergence near $0$ is no accident. The figure also shows that the partial sums overshoot the sawtooth near the jumps, by an amount that does not decrease as $n$ grows; this is the Gibbs phenomenon of Chapter~9, and \cref{sec:uniform-computation} measures it. Abel's series $\sin x-(\sin2x)/2+\cdots$ is $-g(x+\pi)$, which jumps at the odd multiples of $\pi$.

\begin{figure}
\sdfloat{%
\begin{tikzpicture}
\begin{axis}[sdplot, width=0.62\textwidth, height=0.42\textwidth, xmin=0, xmax=6.2832, ymin=-2.05, ymax=2.05, xtick={0,1.5708,3.1416,4.7124,6.2832}, xticklabels={$0$,$\pi/2$,$\pi$,$3\pi/2$,$2\pi$}, ytick={-1.5708,0,1.5708}, yticklabels={$-\pi/2$,$0$,$\pi/2$}, xlabel={$x$}, legend columns=-1, legend style={at={(0.03,0.03)}, anchor=south west}]
\addplot[black!45, line width=1.4pt, domain=0.001:6.2822, samples=2] {(pi-x)/2};
\addlegendentry{$(\pi-x)/2$}
\addplot[black, dashed, line width=0.6pt] coordinates {(0.0000,0.0000) (0.0105,0.0523) (0.0209,0.1046) (0.0314,0.1568) (0.0419,0.2088) (0.0524,0.2605) (0.0628,0.3119) (0.0733,0.3629) (0.0838,0.4135) (0.0942,0.4636) (0.1047,0.5132) (0.1152,0.5621) (0.1257,0.6104) (0.1361,0.6579) (0.1466,0.7047) (0.1571,0.7506) (0.1676,0.7957) (0.1780,0.8398) (0.1885,0.8830) (0.1990,0.9251) (0.2094,0.9662) (0.2199,1.0062) (0.2304,1.0450) (0.2409,1.0826) (0.2513,1.1190) (0.2618,1.1542) (0.2723,1.1881) (0.2827,1.2207) (0.2932,1.2519) (0.3037,1.2818) (0.3142,1.3103) (0.3246,1.3375) (0.3351,1.3632) (0.3456,1.3875) (0.3560,1.4103) (0.3665,1.4318) (0.3770,1.4517) (0.3875,1.4703) (0.3979,1.4873) (0.4084,1.5030) (0.4189,1.5172) (0.4294,1.5299) (0.4398,1.5413) (0.4503,1.5512) (0.4608,1.5597) (0.4712,1.5669) (0.4817,1.5727) (0.4922,1.5772) (0.5027,1.5804) (0.5131,1.5822) (0.5236,1.5829) (0.5341,1.5822) (0.5445,1.5804) (0.5550,1.5775) (0.5655,1.5734) (0.5760,1.5682) (0.5864,1.5620) (0.5969,1.5547) (0.6074,1.5465) (0.6178,1.5373) (0.6283,1.5273) (0.6388,1.5164) (0.6493,1.5047) (0.6597,1.4923) (0.6702,1.4791) (0.6807,1.4653) (0.6912,1.4509) (0.7016,1.4359) (0.7121,1.4204) (0.7226,1.4044) (0.7330,1.3880) (0.7435,1.3713) (0.7540,1.3542) (0.7645,1.3368) (0.7749,1.3192) (0.7854,1.3014) (0.7959,1.2834) (0.8063,1.2654) (0.8168,1.2473) (0.8273,1.2292) (0.8378,1.2112) (0.8482,1.1932) (0.8587,1.1753) (0.8692,1.1575) (0.8796,1.1400) (0.8901,1.1227) (0.9006,1.1056) (0.9111,1.0889) (0.9215,1.0724) (0.9320,1.0563) (0.9425,1.0406) (0.9529,1.0253) (0.9634,1.0104) (0.9739,0.9960) (0.9844,0.9821) (0.9948,0.9686) (1.0053,0.9557) (1.0158,0.9433) (1.0263,0.9314) (1.0367,0.9201) (1.0472,0.9093) (1.0577,0.8991) (1.0681,0.8895) (1.0786,0.8805) (1.0891,0.8721) (1.0996,0.8642) (1.1100,0.8569) (1.1205,0.8502) (1.1310,0.8441) (1.1414,0.8385) (1.1519,0.8335) (1.1624,0.8290) (1.1729,0.8251) (1.1833,0.8217) (1.1938,0.8188) (1.2043,0.8164) (1.2147,0.8145) (1.2252,0.8131) (1.2357,0.8120) (1.2462,0.8114) (1.2566,0.8113) (1.2671,0.8114) (1.2776,0.8120) (1.2881,0.8128) (1.2985,0.8140) (1.3090,0.8155) (1.3195,0.8172) (1.3299,0.8191) (1.3404,0.8213) (1.3509,0.8236) (1.3614,0.8261) (1.3718,0.8286) (1.3823,0.8313) (1.3928,0.8341) (1.4032,0.8369) (1.4137,0.8397) (1.4242,0.8425) (1.4347,0.8452) (1.4451,0.8479) (1.4556,0.8505) (1.4661,0.8530) (1.4765,0.8553) (1.4870,0.8575) (1.4975,0.8595) (1.5080,0.8613) (1.5184,0.8629) (1.5289,0.8642) (1.5394,0.8653) (1.5499,0.8660) (1.5603,0.8665) (1.5708,0.8667) (1.5813,0.8665) (1.5917,0.8660) (1.6022,0.8651) (1.6127,0.8639) (1.6232,0.8623) (1.6336,0.8603) (1.6441,0.8579) (1.6546,0.8552) (1.6650,0.8520) (1.6755,0.8484) (1.6860,0.8445) (1.6965,0.8401) (1.7069,0.8353) (1.7174,0.8301) (1.7279,0.8245) (1.7383,0.8186) (1.7488,0.8122) (1.7593,0.8055) (1.7698,0.7984) (1.7802,0.7909) (1.7907,0.7831) (1.8012,0.7749) (1.8117,0.7664) (1.8221,0.7576) (1.8326,0.7485) (1.8431,0.7391) (1.8535,0.7294) (1.8640,0.7195) (1.8745,0.7094) (1.8850,0.6990) (1.8954,0.6884) (1.9059,0.6777) (1.9164,0.6668) (1.9268,0.6557) (1.9373,0.6446) (1.9478,0.6333) (1.9583,0.6219) (1.9687,0.6105) (1.9792,0.5991) (1.9897,0.5876) (2.0001,0.5761) (2.0106,0.5647) (2.0211,0.5533) (2.0316,0.5420) (2.0420,0.5307) (2.0525,0.5195) (2.0630,0.5085) (2.0735,0.4976) (2.0839,0.4869) (2.0944,0.4763) (2.1049,0.4659) (2.1153,0.4558) (2.1258,0.4458) (2.1363,0.4361) (2.1468,0.4266) (2.1572,0.4174) (2.1677,0.4085) (2.1782,0.3999) (2.1886,0.3915) (2.1991,0.3834) (2.2096,0.3757) (2.2201,0.3683) (2.2305,0.3611) (2.2410,0.3544) (2.2515,0.3479) (2.2619,0.3418) (2.2724,0.3360) (2.2829,0.3305) (2.2934,0.3254) (2.3038,0.3206) (2.3143,0.3161) (2.3248,0.3120) (2.3353,0.3081) (2.3457,0.3046) (2.3562,0.3014) (2.3667,0.2985) (2.3771,0.2958) (2.3876,0.2935) (2.3981,0.2914) (2.4086,0.2896) (2.4190,0.2880) (2.4295,0.2866) (2.4400,0.2855) (2.4504,0.2845) (2.4609,0.2838) (2.4714,0.2832) (2.4819,0.2828) (2.4923,0.2825) (2.5028,0.2824) (2.5133,0.2823) (2.5237,0.2824) (2.5342,0.2825) (2.5447,0.2827) (2.5552,0.2829) (2.5656,0.2831) (2.5761,0.2833) (2.5866,0.2835) (2.5970,0.2837) (2.6075,0.2838) (2.6180,0.2838) (2.6285,0.2838) (2.6389,0.2836) (2.6494,0.2833) (2.6599,0.2829) (2.6704,0.2824) (2.6808,0.2816) (2.6913,0.2807) (2.7018,0.2796) (2.7122,0.2783) (2.7227,0.2768) (2.7332,0.2750) (2.7437,0.2730) (2.7541,0.2707) (2.7646,0.2682) (2.7751,0.2654) (2.7855,0.2623) (2.7960,0.2589) (2.8065,0.2552) (2.8170,0.2513) (2.8274,0.2470) (2.8379,0.2425) (2.8484,0.2376) (2.8588,0.2324) (2.8693,0.2270) (2.8798,0.2212) (2.8903,0.2151) (2.9007,0.2088) (2.9112,0.2021) (2.9217,0.1951) (2.9322,0.1879) (2.9426,0.1804) (2.9531,0.1726) (2.9636,0.1646) (2.9740,0.1563) (2.9845,0.1477) (2.9950,0.1390) (3.0055,0.1300) (3.0159,0.1208) (3.0264,0.1114) (3.0369,0.1019) (3.0473,0.0922) (3.0578,0.0823) (3.0683,0.0723) (3.0788,0.0622) (3.0892,0.0520) (3.0997,0.0417) (3.1102,0.0313) (3.1206,0.0209) (3.1311,0.0105) (3.1416,0.0000) (3.1521,-0.0105) (3.1625,-0.0209) (3.1730,-0.0313) (3.1835,-0.0417) (3.1940,-0.0520) (3.2044,-0.0622) (3.2149,-0.0723) (3.2254,-0.0823) (3.2358,-0.0922) (3.2463,-0.1019) (3.2568,-0.1114) (3.2673,-0.1208) (3.2777,-0.1300) (3.2882,-0.1390) (3.2987,-0.1477) (3.3091,-0.1563) (3.3196,-0.1646) (3.3301,-0.1726) (3.3406,-0.1804) (3.3510,-0.1879) (3.3615,-0.1951) (3.3720,-0.2021) (3.3824,-0.2088) (3.3929,-0.2151) (3.4034,-0.2212) (3.4139,-0.2270) (3.4243,-0.2324) (3.4348,-0.2376) (3.4453,-0.2425) (3.4558,-0.2470) (3.4662,-0.2513) (3.4767,-0.2552) (3.4872,-0.2589) (3.4976,-0.2623) (3.5081,-0.2654) (3.5186,-0.2682) (3.5291,-0.2707) (3.5395,-0.2730) (3.5500,-0.2750) (3.5605,-0.2768) (3.5709,-0.2783) (3.5814,-0.2796) (3.5919,-0.2807) (3.6024,-0.2816) (3.6128,-0.2824) (3.6233,-0.2829) (3.6338,-0.2833) (3.6442,-0.2836) (3.6547,-0.2838) (3.6652,-0.2838) (3.6757,-0.2838) (3.6861,-0.2837) (3.6966,-0.2835) (3.7071,-0.2833) (3.7176,-0.2831) (3.7280,-0.2829) (3.7385,-0.2827) (3.7490,-0.2825) (3.7594,-0.2824) (3.7699,-0.2823) (3.7804,-0.2824) (3.7909,-0.2825) (3.8013,-0.2828) (3.8118,-0.2832) (3.8223,-0.2838) (3.8327,-0.2845) (3.8432,-0.2855) (3.8537,-0.2866) (3.8642,-0.2880) (3.8746,-0.2896) (3.8851,-0.2914) (3.8956,-0.2935) (3.9060,-0.2958) (3.9165,-0.2985) (3.9270,-0.3014) (3.9375,-0.3046) (3.9479,-0.3081) (3.9584,-0.3120) (3.9689,-0.3161) (3.9794,-0.3206) (3.9898,-0.3254) (4.0003,-0.3305) (4.0108,-0.3360) (4.0212,-0.3418) (4.0317,-0.3479) (4.0422,-0.3544) (4.0527,-0.3611) (4.0631,-0.3683) (4.0736,-0.3757) (4.0841,-0.3834) (4.0945,-0.3915) (4.1050,-0.3999) (4.1155,-0.4085) (4.1260,-0.4174) (4.1364,-0.4266) (4.1469,-0.4361) (4.1574,-0.4458) (4.1678,-0.4558) (4.1783,-0.4659) (4.1888,-0.4763) (4.1993,-0.4869) (4.2097,-0.4976) (4.2202,-0.5085) (4.2307,-0.5195) (4.2412,-0.5307) (4.2516,-0.5420) (4.2621,-0.5533) (4.2726,-0.5647) (4.2830,-0.5761) (4.2935,-0.5876) (4.3040,-0.5991) (4.3145,-0.6105) (4.3249,-0.6219) (4.3354,-0.6333) (4.3459,-0.6446) (4.3563,-0.6557) (4.3668,-0.6668) (4.3773,-0.6777) (4.3878,-0.6884) (4.3982,-0.6990) (4.4087,-0.7094) (4.4192,-0.7195) (4.4296,-0.7294) (4.4401,-0.7391) (4.4506,-0.7485) (4.4611,-0.7576) (4.4715,-0.7664) (4.4820,-0.7749) (4.4925,-0.7831) (4.5029,-0.7909) (4.5134,-0.7984) (4.5239,-0.8055) (4.5344,-0.8122) (4.5448,-0.8186) (4.5553,-0.8245) (4.5658,-0.8301) (4.5763,-0.8353) (4.5867,-0.8401) (4.5972,-0.8445) (4.6077,-0.8484) (4.6181,-0.8520) (4.6286,-0.8552) (4.6391,-0.8579) (4.6496,-0.8603) (4.6600,-0.8623) (4.6705,-0.8639) (4.6810,-0.8651) (4.6914,-0.8660) (4.7019,-0.8665) (4.7124,-0.8667) (4.7229,-0.8665) (4.7333,-0.8660) (4.7438,-0.8653) (4.7543,-0.8642) (4.7647,-0.8629) (4.7752,-0.8613) (4.7857,-0.8595) (4.7962,-0.8575) (4.8066,-0.8553) (4.8171,-0.8530) (4.8276,-0.8505) (4.8381,-0.8479) (4.8485,-0.8452) (4.8590,-0.8425) (4.8695,-0.8397) (4.8799,-0.8369) (4.8904,-0.8341) (4.9009,-0.8313) (4.9114,-0.8286) (4.9218,-0.8261) (4.9323,-0.8236) (4.9428,-0.8213) (4.9532,-0.8191) (4.9637,-0.8172) (4.9742,-0.8155) (4.9847,-0.8140) (4.9951,-0.8128) (5.0056,-0.8120) (5.0161,-0.8114) (5.0265,-0.8113) (5.0370,-0.8114) (5.0475,-0.8120) (5.0580,-0.8131) (5.0684,-0.8145) (5.0789,-0.8164) (5.0894,-0.8188) (5.0999,-0.8217) (5.1103,-0.8251) (5.1208,-0.8290) (5.1313,-0.8335) (5.1417,-0.8385) (5.1522,-0.8441) (5.1627,-0.8502) (5.1732,-0.8569) (5.1836,-0.8642) (5.1941,-0.8721) (5.2046,-0.8805) (5.2150,-0.8895) (5.2255,-0.8991) (5.2360,-0.9093) (5.2465,-0.9201) (5.2569,-0.9314) (5.2674,-0.9433) (5.2779,-0.9557) (5.2883,-0.9686) (5.2988,-0.9821) (5.3093,-0.9960) (5.3198,-1.0104) (5.3302,-1.0253) (5.3407,-1.0406) (5.3512,-1.0563) (5.3617,-1.0724) (5.3721,-1.0889) (5.3826,-1.1056) (5.3931,-1.1227) (5.4035,-1.1400) (5.4140,-1.1575) (5.4245,-1.1753) (5.4350,-1.1932) (5.4454,-1.2112) (5.4559,-1.2292) (5.4664,-1.2473) (5.4768,-1.2654) (5.4873,-1.2834) (5.4978,-1.3014) (5.5083,-1.3192) (5.5187,-1.3368) (5.5292,-1.3542) (5.5397,-1.3713) (5.5501,-1.3880) (5.5606,-1.4044) (5.5711,-1.4204) (5.5816,-1.4359) (5.5920,-1.4509) (5.6025,-1.4653) (5.6130,-1.4791) (5.6235,-1.4923) (5.6339,-1.5047) (5.6444,-1.5164) (5.6549,-1.5273) (5.6653,-1.5373) (5.6758,-1.5465) (5.6863,-1.5547) (5.6968,-1.5620) (5.7072,-1.5682) (5.7177,-1.5734) (5.7282,-1.5775) (5.7386,-1.5804) (5.7491,-1.5822) (5.7596,-1.5829) (5.7701,-1.5822) (5.7805,-1.5804) (5.7910,-1.5772) (5.8015,-1.5727) (5.8119,-1.5669) (5.8224,-1.5597) (5.8329,-1.5512) (5.8434,-1.5413) (5.8538,-1.5299) (5.8643,-1.5172) (5.8748,-1.5030) (5.8853,-1.4873) (5.8957,-1.4703) (5.9062,-1.4517) (5.9167,-1.4318) (5.9271,-1.4103) (5.9376,-1.3875) (5.9481,-1.3632) (5.9586,-1.3375) (5.9690,-1.3103) (5.9795,-1.2818) (5.9900,-1.2519) (6.0004,-1.2207) (6.0109,-1.1881) (6.0214,-1.1542) (6.0319,-1.1190) (6.0423,-1.0826) (6.0528,-1.0450) (6.0633,-1.0062) (6.0737,-0.9662) (6.0842,-0.9251) (6.0947,-0.8830) (6.1052,-0.8398) (6.1156,-0.7957) (6.1261,-0.7506) (6.1366,-0.7047) (6.1470,-0.6579) (6.1575,-0.6104) (6.1680,-0.5621) (6.1785,-0.5132) (6.1889,-0.4636) (6.1994,-0.4135) (6.2099,-0.3629) (6.2204,-0.3119) (6.2308,-0.2605) (6.2413,-0.2088) (6.2518,-0.1568) (6.2622,-0.1046) (6.2727,-0.0523) (6.2832,-0.0000)};
\addlegendentry{$n=5$}
\addplot[black, line width=0.6pt] coordinates {(0.0000,0.0000) (0.0105,0.2607) (0.0209,0.5152) (0.0314,0.7574) (0.0419,0.9818) (0.0524,1.1837) (0.0628,1.3592) (0.0733,1.5057) (0.0838,1.6216) (0.0942,1.7065) (0.1047,1.7612) (0.1152,1.7876) (0.1257,1.7887) (0.1361,1.7680) (0.1466,1.7298) (0.1571,1.6787) (0.1676,1.6194) (0.1780,1.5565) (0.1885,1.4943) (0.1990,1.4364) (0.2094,1.3861) (0.2199,1.3455) (0.2304,1.3163) (0.2409,1.2989) (0.2513,1.2933) (0.2618,1.2985) (0.2723,1.3129) (0.2827,1.3346) (0.2932,1.3611) (0.3037,1.3899) (0.3142,1.4186) (0.3246,1.4448) (0.3351,1.4666) (0.3456,1.4823) (0.3560,1.4909) (0.3665,1.4918) (0.3770,1.4850) (0.3875,1.4711) (0.3979,1.4510) (0.4084,1.4261) (0.4189,1.3979) (0.4294,1.3683) (0.4398,1.3389) (0.4503,1.3115) (0.4608,1.2874) (0.4712,1.2679) (0.4817,1.2537) (0.4922,1.2453) (0.5027,1.2425) (0.5131,1.2451) (0.5236,1.2522) (0.5341,1.2628) (0.5445,1.2758) (0.5550,1.2896) (0.5655,1.3031) (0.5760,1.3148) (0.5864,1.3237) (0.5969,1.3289) (0.6074,1.3297) (0.6178,1.3260) (0.6283,1.3177) (0.6388,1.3052) (0.6493,1.2892) (0.6597,1.2704) (0.6702,1.2500) (0.6807,1.2291) (0.6912,1.2086) (0.7016,1.1898) (0.7121,1.1735) (0.7226,1.1604) (0.7330,1.1510) (0.7435,1.1454) (0.7540,1.1436) (0.7645,1.1452) (0.7749,1.1497) (0.7854,1.1564) (0.7959,1.1642) (0.8063,1.1724) (0.8168,1.1799) (0.8273,1.1858) (0.8378,1.1895) (0.8482,1.1904) (0.8587,1.1880) (0.8692,1.1822) (0.8796,1.1733) (0.8901,1.1614) (0.9006,1.1472) (0.9111,1.1313) (0.9215,1.1145) (0.9320,1.0976) (0.9425,1.0815) (0.9529,1.0668) (0.9634,1.0543) (0.9739,1.0444) (0.9844,1.0373) (0.9948,1.0332) (1.0053,1.0319) (1.0158,1.0331) (1.0263,1.0362) (1.0367,1.0407) (1.0472,1.0459) (1.0577,1.0509) (1.0681,1.0552) (1.0786,1.0580) (1.0891,1.0588) (1.0996,1.0571) (1.1100,1.0529) (1.1205,1.0461) (1.1310,1.0367) (1.1414,1.0253) (1.1519,1.0121) (1.1624,0.9979) (1.1729,0.9832) (1.1833,0.9687) (1.1938,0.9552) (1.2043,0.9430) (1.2147,0.9328) (1.2252,0.9248) (1.2357,0.9192) (1.2462,0.9160) (1.2566,0.9150) (1.2671,0.9159) (1.2776,0.9182) (1.2881,0.9213) (1.2985,0.9248) (1.3090,0.9278) (1.3195,0.9300) (1.3299,0.9307) (1.3404,0.9296) (1.3509,0.9264) (1.3614,0.9210) (1.3718,0.9134) (1.3823,0.9038) (1.3928,0.8925) (1.4032,0.8801) (1.4137,0.8670) (1.4242,0.8538) (1.4347,0.8410) (1.4451,0.8291) (1.4556,0.8187) (1.4661,0.8101) (1.4765,0.8035) (1.4870,0.7989) (1.4975,0.7963) (1.5080,0.7956) (1.5184,0.7962) (1.5289,0.7979) (1.5394,0.8001) (1.5499,0.8023) (1.5603,0.8040) (1.5708,0.8046) (1.5813,0.8038) (1.5917,0.8013) (1.6022,0.7970) (1.6127,0.7907) (1.6232,0.7825) (1.6336,0.7727) (1.6441,0.7617) (1.6546,0.7498) (1.6650,0.7375) (1.6755,0.7253) (1.6860,0.7138) (1.6965,0.7033) (1.7069,0.6942) (1.7174,0.6868) (1.7279,0.6811) (1.7383,0.6773) (1.7488,0.6753) (1.7593,0.6746) (1.7698,0.6751) (1.7802,0.6763) (1.7907,0.6778) (1.8012,0.6790) (1.8117,0.6796) (1.8221,0.6791) (1.8326,0.6772) (1.8431,0.6736) (1.8535,0.6683) (1.8640,0.6613) (1.8745,0.6527) (1.8850,0.6428) (1.8954,0.6319) (1.9059,0.6204) (1.9164,0.6088) (1.9268,0.5975) (1.9373,0.5869) (1.9478,0.5775) (1.9583,0.5694) (1.9687,0.5630) (1.9792,0.5582) (1.9897,0.5550) (2.0001,0.5533) (2.0106,0.5528) (2.0211,0.5532) (2.0316,0.5540) (2.0420,0.5549) (2.0525,0.5553) (2.0630,0.5550) (2.0735,0.5535) (2.0839,0.5507) (2.0944,0.5462) (2.1049,0.5402) (2.1153,0.5326) (2.1258,0.5237) (2.1363,0.5136) (2.1468,0.5029) (2.1572,0.4917) (2.1677,0.4806) (2.1782,0.4701) (2.1886,0.4603) (2.1991,0.4517) (2.2096,0.4446) (2.2201,0.4389) (2.2305,0.4348) (2.2410,0.4322) (2.2515,0.4308) (2.2619,0.4304) (2.2724,0.4307) (2.2829,0.4312) (2.2934,0.4315) (2.3038,0.4313) (2.3143,0.4302) (2.3248,0.4279) (2.3353,0.4242) (2.3457,0.4190) (2.3562,0.4124) (2.3667,0.4043) (2.3771,0.3950) (2.3876,0.3849) (2.3981,0.3742) (2.4086,0.3634) (2.4190,0.3528) (2.4295,0.3428) (2.4400,0.3338) (2.4504,0.3260) (2.4609,0.3196) (2.4714,0.3147) (2.4819,0.3112) (2.4923,0.3090) (2.5028,0.3080) (2.5133,0.3077) (2.5237,0.3078) (2.5342,0.3080) (2.5447,0.3079) (2.5552,0.3071) (2.5656,0.3053) (2.5761,0.3023) (2.5866,0.2979) (2.5970,0.2920) (2.6075,0.2847) (2.6180,0.2762) (2.6285,0.2667) (2.6389,0.2565) (2.6494,0.2459) (2.6599,0.2353) (2.6704,0.2252) (2.6808,0.2158) (2.6913,0.2074) (2.7018,0.2003) (2.7122,0.1946) (2.7227,0.1904) (2.7332,0.1874) (2.7437,0.1857) (2.7541,0.1849) (2.7646,0.1847) (2.7751,0.1848) (2.7855,0.1847) (2.7960,0.1842) (2.8065,0.1828) (2.8170,0.1804) (2.8274,0.1767) (2.8379,0.1716) (2.8484,0.1650) (2.8588,0.1572) (2.8693,0.1483) (2.8798,0.1385) (2.8903,0.1282) (2.9007,0.1177) (2.9112,0.1074) (2.9217,0.0977) (2.9322,0.0888) (2.9426,0.0811) (2.9531,0.0747) (2.9636,0.0696) (2.9740,0.0660) (2.9845,0.0636) (2.9950,0.0622) (3.0055,0.0617) (3.0159,0.0616) (3.0264,0.0616) (3.0369,0.0613) (3.0473,0.0603) (3.0578,0.0584) (3.0683,0.0554) (3.0788,0.0510) (3.0892,0.0452) (3.0997,0.0381) (3.1102,0.0298) (3.1206,0.0205) (3.1311,0.0104) (3.1416,0.0000) (3.1521,-0.0104) (3.1625,-0.0205) (3.1730,-0.0298) (3.1835,-0.0381) (3.1940,-0.0452) (3.2044,-0.0510) (3.2149,-0.0554) (3.2254,-0.0584) (3.2358,-0.0603) (3.2463,-0.0613) (3.2568,-0.0616) (3.2673,-0.0616) (3.2777,-0.0617) (3.2882,-0.0622) (3.2987,-0.0636) (3.3091,-0.0660) (3.3196,-0.0696) (3.3301,-0.0747) (3.3406,-0.0811) (3.3510,-0.0888) (3.3615,-0.0977) (3.3720,-0.1074) (3.3824,-0.1177) (3.3929,-0.1282) (3.4034,-0.1385) (3.4139,-0.1483) (3.4243,-0.1572) (3.4348,-0.1650) (3.4453,-0.1716) (3.4558,-0.1767) (3.4662,-0.1804) (3.4767,-0.1828) (3.4872,-0.1842) (3.4976,-0.1847) (3.5081,-0.1848) (3.5186,-0.1847) (3.5291,-0.1849) (3.5395,-0.1857) (3.5500,-0.1874) (3.5605,-0.1904) (3.5709,-0.1946) (3.5814,-0.2003) (3.5919,-0.2074) (3.6024,-0.2158) (3.6128,-0.2252) (3.6233,-0.2353) (3.6338,-0.2459) (3.6442,-0.2565) (3.6547,-0.2667) (3.6652,-0.2762) (3.6757,-0.2847) (3.6861,-0.2920) (3.6966,-0.2979) (3.7071,-0.3023) (3.7176,-0.3053) (3.7280,-0.3071) (3.7385,-0.3079) (3.7490,-0.3080) (3.7594,-0.3078) (3.7699,-0.3077) (3.7804,-0.3080) (3.7909,-0.3090) (3.8013,-0.3112) (3.8118,-0.3147) (3.8223,-0.3196) (3.8327,-0.3260) (3.8432,-0.3338) (3.8537,-0.3428) (3.8642,-0.3528) (3.8746,-0.3634) (3.8851,-0.3742) (3.8956,-0.3849) (3.9060,-0.3950) (3.9165,-0.4043) (3.9270,-0.4124) (3.9375,-0.4190) (3.9479,-0.4242) (3.9584,-0.4279) (3.9689,-0.4302) (3.9794,-0.4313) (3.9898,-0.4315) (4.0003,-0.4312) (4.0108,-0.4307) (4.0212,-0.4304) (4.0317,-0.4308) (4.0422,-0.4322) (4.0527,-0.4348) (4.0631,-0.4389) (4.0736,-0.4446) (4.0841,-0.4517) (4.0945,-0.4603) (4.1050,-0.4701) (4.1155,-0.4806) (4.1260,-0.4917) (4.1364,-0.5029) (4.1469,-0.5136) (4.1574,-0.5237) (4.1678,-0.5326) (4.1783,-0.5402) (4.1888,-0.5462) (4.1993,-0.5507) (4.2097,-0.5535) (4.2202,-0.5550) (4.2307,-0.5553) (4.2412,-0.5549) (4.2516,-0.5540) (4.2621,-0.5532) (4.2726,-0.5528) (4.2830,-0.5533) (4.2935,-0.5550) (4.3040,-0.5582) (4.3145,-0.5630) (4.3249,-0.5694) (4.3354,-0.5775) (4.3459,-0.5869) (4.3563,-0.5975) (4.3668,-0.6088) (4.3773,-0.6204) (4.3878,-0.6319) (4.3982,-0.6428) (4.4087,-0.6527) (4.4192,-0.6613) (4.4296,-0.6683) (4.4401,-0.6736) (4.4506,-0.6772) (4.4611,-0.6791) (4.4715,-0.6796) (4.4820,-0.6790) (4.4925,-0.6778) (4.5029,-0.6763) (4.5134,-0.6751) (4.5239,-0.6746) (4.5344,-0.6753) (4.5448,-0.6773) (4.5553,-0.6811) (4.5658,-0.6868) (4.5763,-0.6942) (4.5867,-0.7033) (4.5972,-0.7138) (4.6077,-0.7253) (4.6181,-0.7375) (4.6286,-0.7498) (4.6391,-0.7617) (4.6496,-0.7727) (4.6600,-0.7825) (4.6705,-0.7907) (4.6810,-0.7970) (4.6914,-0.8013) (4.7019,-0.8038) (4.7124,-0.8046) (4.7229,-0.8040) (4.7333,-0.8023) (4.7438,-0.8001) (4.7543,-0.7979) (4.7647,-0.7962) (4.7752,-0.7956) (4.7857,-0.7963) (4.7962,-0.7989) (4.8066,-0.8035) (4.8171,-0.8101) (4.8276,-0.8187) (4.8381,-0.8291) (4.8485,-0.8410) (4.8590,-0.8538) (4.8695,-0.8670) (4.8799,-0.8801) (4.8904,-0.8925) (4.9009,-0.9038) (4.9114,-0.9134) (4.9218,-0.9210) (4.9323,-0.9264) (4.9428,-0.9296) (4.9532,-0.9307) (4.9637,-0.9300) (4.9742,-0.9278) (4.9847,-0.9248) (4.9951,-0.9213) (5.0056,-0.9182) (5.0161,-0.9159) (5.0265,-0.9150) (5.0370,-0.9160) (5.0475,-0.9192) (5.0580,-0.9248) (5.0684,-0.9328) (5.0789,-0.9430) (5.0894,-0.9552) (5.0999,-0.9687) (5.1103,-0.9832) (5.1208,-0.9979) (5.1313,-1.0121) (5.1417,-1.0253) (5.1522,-1.0367) (5.1627,-1.0461) (5.1732,-1.0529) (5.1836,-1.0571) (5.1941,-1.0588) (5.2046,-1.0580) (5.2150,-1.0552) (5.2255,-1.0509) (5.2360,-1.0459) (5.2465,-1.0407) (5.2569,-1.0362) (5.2674,-1.0331) (5.2779,-1.0319) (5.2883,-1.0332) (5.2988,-1.0373) (5.3093,-1.0444) (5.3198,-1.0543) (5.3302,-1.0668) (5.3407,-1.0815) (5.3512,-1.0976) (5.3617,-1.1145) (5.3721,-1.1313) (5.3826,-1.1472) (5.3931,-1.1614) (5.4035,-1.1733) (5.4140,-1.1822) (5.4245,-1.1880) (5.4350,-1.1904) (5.4454,-1.1895) (5.4559,-1.1858) (5.4664,-1.1799) (5.4768,-1.1724) (5.4873,-1.1642) (5.4978,-1.1564) (5.5083,-1.1497) (5.5187,-1.1452) (5.5292,-1.1436) (5.5397,-1.1454) (5.5501,-1.1510) (5.5606,-1.1604) (5.5711,-1.1735) (5.5816,-1.1898) (5.5920,-1.2086) (5.6025,-1.2291) (5.6130,-1.2500) (5.6235,-1.2704) (5.6339,-1.2892) (5.6444,-1.3052) (5.6549,-1.3177) (5.6653,-1.3260) (5.6758,-1.3297) (5.6863,-1.3289) (5.6968,-1.3237) (5.7072,-1.3148) (5.7177,-1.3031) (5.7282,-1.2896) (5.7386,-1.2758) (5.7491,-1.2628) (5.7596,-1.2522) (5.7701,-1.2451) (5.7805,-1.2425) (5.7910,-1.2453) (5.8015,-1.2537) (5.8119,-1.2679) (5.8224,-1.2874) (5.8329,-1.3115) (5.8434,-1.3389) (5.8538,-1.3683) (5.8643,-1.3979) (5.8748,-1.4261) (5.8853,-1.4510) (5.8957,-1.4711) (5.9062,-1.4850) (5.9167,-1.4918) (5.9271,-1.4909) (5.9376,-1.4823) (5.9481,-1.4666) (5.9586,-1.4448) (5.9690,-1.4186) (5.9795,-1.3899) (5.9900,-1.3611) (6.0004,-1.3346) (6.0109,-1.3129) (6.0214,-1.2985) (6.0319,-1.2933) (6.0423,-1.2989) (6.0528,-1.3163) (6.0633,-1.3455) (6.0737,-1.3861) (6.0842,-1.4364) (6.0947,-1.4943) (6.1052,-1.5565) (6.1156,-1.6194) (6.1261,-1.6787) (6.1366,-1.7298) (6.1470,-1.7680) (6.1575,-1.7887) (6.1680,-1.7876) (6.1785,-1.7612) (6.1889,-1.7065) (6.1994,-1.6216) (6.2099,-1.5057) (6.2204,-1.3592) (6.2308,-1.1837) (6.2413,-0.9818) (6.2518,-0.7574) (6.2622,-0.5152) (6.2727,-0.2607) (6.2832,-0.0000)};
\addlegendentry{$n=25$}
\end{axis}
\end{tikzpicture}}{The partial sums $\sum_{k=1}^n\sin(kx)/k$ for $n=5$ (dashed) and $n=25$ (solid), and the sum $(\pi-x)/2$ of the series on $(0,2\pi)$ (grey). The series converges uniformly on $[\delta,2\pi-\delta]$, but not near $0$ and $2\pi$, where every partial sum vanishes while the sum tends to $\pm\pi/2$; near these points the partial sums overshoot the sum by about $0.28$.\label{fig:sawtooth}}
\end{figure}

\subsection{Continuity, Integration and Differentiation Term by Term}

The theorems of \cref{sec:limit-continuity,sec:limit-integration}, applied to partial sums, give the rules for series.

\begin{corollary}[Series term by term]
\label{cor:series-termwise}
Let $f_k$ be functions on a set $S$ with values in $\mathbb F$.
\begin{enumerate}
\item If $S$ is a metric space, the $f_k$ are continuous and $\sum_kf_k$ converges locally uniformly, then its sum is continuous.
\item If $S=[a,b]$, the $f_k$ are integrable and $\sum_kf_k$ converges uniformly, then $\int_a^b\sum_kf_k=\sum_k\int_a^bf_k$.
\item If $S$ is a bounded interval, the $f_k$ are differentiable, $\sum_kf_k'$ converges uniformly and $\sum_kf_k(x_0)$ converges for some $x_0\in S$, then $\sum_kf_k$ converges uniformly, its sum is differentiable, and $(\sum_kf_k)'=\sum_kf_k'$.
\end{enumerate}
\end{corollary}

\begin{proof}
(i) is \cref{cor:limit-continuous}\,(ii), (ii) is \cref{cor:termwise-integration}, and (iii) is \cref{thm:uniform-derivative}, applied to the partial sums.
\end{proof}

A useful consequence of the M-test and \cref{thm:exchange-limits} concerns series whose terms depend on a further index. It is known as Tannery's theorem.

\begin{corollary}[Limits of series]
\label{cor:tannery}
Let $a_k(n)\in\mathbb F$ for $k,n\in\mathbb N$, with $a_k(n)\to a_k$ as $n\to\infty$ for every $k$, and suppose that $|a_k(n)|\le M_k$ for all $k$ and $n$, where $\sum_kM_k<\infty$. Then $\sum_ka_k$ converges absolutely, and
\[
\lim_{n\to\infty}\sum_{k=0}^\infty a_k(n)=\sum_{k=0}^\infty a_k .
\]
\end{corollary}

\begin{proof}
Regard $a_k$ as a function of $n\in\mathbb N$. By the M-test, $\sum_ka_k(n)$ converges uniformly in $n$, and its partial sums $\sum_{k\le K}a_k(n)$ tend to $\sum_{k\le K}a_k$ as $n\to\infty$. By \cref{thm:exchange-limits}, applied to the partial sums as functions on $\mathbb N$ with the limit $n\to\infty$, the limits $\sum_{k\le K}a_k$ converge as $K\to\infty$, to $\lim_n\sum_ka_k(n)$. Absolute convergence follows from $|a_k|\le M_k$.
\end{proof}

For example, by the binomial theorem, $(1+x/n)^n=\sum_{k=0}^\infty a_k(n)$ with
\[
a_k(n)\coloneqq\binom nk\frac{x^k}{n^k}=\Bigl(1-\frac1n\Bigr)\Bigl(1-\frac2n\Bigr)\cdots\Bigl(1-\frac{k-1}n\Bigr)\frac{x^k}{k!} ,
\]
which vanishes for $k>n$. The product of the factors $1-j/n$ lies between $0$ and $1$ for $k\le n$ and tends to $1$ as $n\to\infty$, so the corollary applies with $M_k=|x|^k/k!$ and gives $\lim_n(1+x/n)^n=\sum_kx^k/k!=e^x$, the formula that \vcref{1}{thm:exp-limit} proved by logarithms.

\subsection{Examples}

The series $f(x)\coloneqq\sum_{k\ge1}x^k/k^2$ converges uniformly on $[-1,1]$, and its sum is continuous there. The series of derivatives, $\sum_{k\ge1}x^{k-1}/k$, converges uniformly on every $[-1,r]$ with $r<1$: on $[-r,r]$ by the M-test with $M_k=r^{k-1}/k$, and on $[-1,0]$ by Leibniz's test as in the example after \cref{thm:m-test}, since for $x\in[-1,0]$ it is an alternating series with decreasing terms. By \cref{cor:series-termwise}\,(iii), $f$ is differentiable on $[-1,1)$, with
\[
f'(x)=\sum_{k=1}^\infty\frac{x^{k-1}}k=-\frac{\ln(1-x)}x\qquad(x\ne0),
\]
by \vcref{1}{thm:log-series}. As $x\to1^-$, $f'(x)\to\infty$, and $f$ is not differentiable at $1$ from the left: by the mean value theorem the difference quotients $(f(1)-f(x))/(1-x)$ are values of $f'$ between $x$ and $1$, and they tend to $\infty$. So a uniformly convergent power series can have a sum that is continuous but not differentiable at an end of its interval of convergence. The sum $f$ is the dilogarithm $\mathrm{Li}_2$, and $f(1)=\sum1/k^2=\pi^2/6$, as Chapter~9 proves.

The next example is the zeta function of \vcref{1}{sec:nonnegative-series}.

\begin{proposition}[The zeta function]
\label{prop:zeta-smooth}
The series $\zeta(s)\coloneqq\sum_{n=1}^\infty n^{-s}$ converges uniformly on every interval $[s_0,\infty)$ with $s_0>1$, but not on $(1,\infty)$. Its sum is infinitely differentiable on $(1,\infty)$, with
\[
\zeta^{(k)}(s)=(-1)^k\sum_{n=2}^\infty\frac{(\ln n)^k}{n^s}\qquad(k\ge1);
\]
it is decreasing and convex, $\zeta(s)\to1$ as $s\to\infty$, and $\zeta(s)\ge1/(s-1)$, so that $\zeta(s)\to\infty$ as $s\to1^+$.
\end{proposition}

\begin{proof}
Let $s_0>1$ and $k\in\mathbb N$. Choose rational numbers $p>0$ and $q>1$ with $p+q\le s_0$. By \vcref{1}{prop:ln-growth}, $(\ln n)^k\le Cn^p$ for all $n\ge1$, with a constant $C$, so that for $s\ge s_0$
\[
\Bigl|\frac{(\ln n)^k}{n^s}\Bigr|\le\frac{(\ln n)^k}{n^{s_0}}\le\frac C{n^{q}},
\]
and $\sum_nn^{-q}<\infty$ (\vcref{1}{cor:p-series}). By the M-test, the series $\sum_n(-\ln n)^kn^{-s}$ converges uniformly on $[s_0,\infty)$ for every $k$; it is the series of the $k$th derivatives of the terms $n^{-s}=e^{-s\ln n}$ (\vcref{1}{thm:power-derivatives}). Applying \cref{cor:series-termwise}\,(iii) on the bounded intervals $[s_0,s_1]$, $k$ times, proves the formula for $\zeta^{(k)}$ on $(1,\infty)$. Then $\zeta'<0$ and $\zeta''>0$, since the terms with $n\ge2$ are of constant sign. On $[2,\infty)$ the partial sums $\sum_{n\le N}n^{-s}$ converge uniformly and tend to $1$ as $s\to\infty$, so $\zeta(s)\to1$ by \cref{thm:exchange-limits}. Comparison with the integral, as in the proof of \vcref{1}{thm:integral-test}, gives $\zeta(s)\ge\int_1^\infty x^{-s}\,dx=1/(s-1)$ for $s>1$. Finally, the convergence is not uniform on $(1,\infty)$: there each partial sum is bounded, by $N$, while $\zeta$ is unbounded, and a uniform limit of bounded functions is bounded (\cref{prop:uniform-sup-norm}).
\end{proof}

The exponential, sine and cosine, defined by power series in \vcref{1}{def:exp,def:sin-cos}, fit the same pattern: their series converge uniformly on every bounded interval, so their sums are continuous and can be differentiated term by term, which gives $\exp'=\exp$, $\sin'=\cos$ and $\cos'=-\sin$. Volume~I proved these facts directly, with the estimates of \vcref{1}{thm:termwise-diff}; the present theorems give the same results without estimates specific to power series, and \cref{sec:power-series-revisited} returns to this.

\subsection{Infinite Products}

Products of infinitely many factors close to $1$ converge uniformly under a condition of M-test type. The partial products are compared with exponentials through the inequality $1+t\le e^t$ for real $t$ (\vcref{1}{prop:ln-inequalities}).

\begin{proposition}[Uniform convergence of products]
\label{prop:uniform-products}
Let $u_k\colon S\to\mathbb C$ be functions with $|u_k(s)|\le M_k$ for all $s$, where $\sum_kM_k<\infty$. Then the partial products $P_n\coloneqq\prod_{k=0}^n(1+u_k)$ converge uniformly on $S$ to a function $P$ with $|P(s)|\le\exp(\sum_kM_k)$. Moreover $P(s)=0$ if and only if $1+u_k(s)=0$ for some $k$.
\end{proposition}

\begin{proof}
For complex numbers $w_1,\ldots,w_n$,
\begin{equation}
\Bigl|\prod_{k=1}^n(1+w_k)-1\Bigr|\le\prod_{k=1}^n(1+|w_k|)-1\le\exp\Bigl(\sum_{k=1}^n|w_k|\Bigr)-1 .
\label{eq:product-bound}
\end{equation}
The second inequality follows from $1+t\le e^t$. The first holds for $n=1$, and if it holds for $n-1$, then with $Q\coloneqq\prod_{k<n}(1+w_k)$ and $Q'\coloneqq\prod_{k<n}(1+|w_k|)$,
\[
|Q(1+w_n)-1|=|(Q-1)(1+w_n)+w_n|\le(Q'-1)(1+|w_n|)+|w_n|=Q'(1+|w_n|)-1 .
\]
Let $C\coloneqq\exp(\sum_kM_k)$. By \eqref{eq:product-bound}, $|P_m|\le C$, and for $n>m$,
\[
|P_n-P_m|=|P_m|\,\Bigl|\prod_{k=m+1}^n(1+u_k)-1\Bigr|\le C\Bigl(\exp\Bigl(\sum_{k>m}M_k\Bigr)-1\Bigr)
\]
at every point, which tends to $0$ as $m\to\infty$. So $(P_n)$ is uniformly Cauchy and converges uniformly to some $P$, with $|P|\le C$. If no factor vanishes at $s$, choose $m$ with $\sum_{k>m}M_k<\ln2$. Then, by \eqref{eq:product-bound}, $|\prod_{k=m+1}^n(1+u_k(s))|\ge1-(\exp(\sum_{k>m}M_k)-1)$ for all $n>m$, and the right side, $\eta\coloneqq2-\exp(\sum_{k>m}M_k)$, is positive, and $|P(s)|\ge|P_m(s)|\,\eta>0$. If some factor vanishes at $s$, so do all $P_n(s)$ from then on, and $P(s)=0$.
\end{proof}

For example, for $|x|\le r<1$ the product $\prod_{k\ge0}(1+x^{2^k})$ converges uniformly, with $M_k=r^{2^k}$; multiplying out $\prod_{k=0}^n(1+x^{2^k})$ gives every power $x^j$ with $0\le j<2^{n+1}$ exactly once, by the uniqueness of binary expansions, so the product is $\sum_jx^j=1/(1-x)$. The products that matter most, such as Euler's product $\sin\pi x=\pi x\prod_{k\ge1}(1-x^2/k^2)$, which converges uniformly on bounded sets by the proposition, are the subject of Chapter~15, where they are studied for complex arguments.

\section{Double Sequences and Double Series}
\label{sec:double-series}

\subsection{Double Limits}

A double sequence $(a_{mn})$ in a metric space has the iterated limits of \cref{sec:limit-continuity} and a third notion of limit, in which both indices grow together. It \term{converges} to $L$, written $\lim_{m,n\to\infty}a_{mn}=L$, if for every $\varepsilon>0$ there is $N$ with $d(a_{mn},L)<\varepsilon$ for all $m,n\ge N$. This is the limit of the function $(m,n)\mapsto a_{mn}$ on $\mathbb N^2$ as both arguments tend to infinity, and it is unique when it exists.

\begin{proposition}[Double and iterated limits]
\label{prop:double-limits}
Let $(a_{mn})$ be a double sequence in a metric space $Y$.
\begin{enumerate}
\item If $\lim_{m,n}a_{mn}=L$ and $b_m\coloneqq\lim_na_{mn}$ exists for every $m$, then $b_m\to L$; that is, $\lim_m\lim_na_{mn}=L$.
\item If $b_m\coloneqq\lim_na_{mn}$ exists for every $m$, and $c_n\coloneqq\lim_ma_{mn}$ exists uniformly in $n$, and $Y$ is complete, then the double limit exists and equals both iterated limits.
\end{enumerate}
\end{proposition}

\begin{proof}
(i) Given $\varepsilon>0$, choose $N$ as in the definition. For $m\ge N$, letting $n\to\infty$ in $d(a_{mn},L)<\varepsilon$ gives $d(b_m,L)\le\varepsilon$. (ii) By \cref{cor:iterated-limits}, both iterated limits exist and are equal to some $L$, and $c_n\to L$. Given $\varepsilon>0$, choose $N$ with $d(a_{mn},c_n)<\varepsilon$ for $m\ge N$ and all $n$, and with $d(c_n,L)<\varepsilon$ for $n\ge N$. For $m,n\ge N$, $d(a_{mn},L)<2\varepsilon$.
\end{proof}

The three notions are independent. For $a_{mn}=m/(m+n)$ the iterated limits are $0$ and $1$, and there is no double limit, by (i). For $a_{mn}=(-1)^m(1/m+1/n)$, with $m,n\ge1$, the double limit is $0$, since $|a_{mn}|\le1/m+1/n$; the iterated limit $\lim_m\lim_na_{mn}=\lim_m(-1)^m/m=0$ exists, but $\lim_ma_{mn}$ exists for no $n$. For $a_{mn}=mn/(m^2+n^2)$ both iterated limits are $0$, while $a_{nn}=1/2$, so there is no double limit. Part (ii) is the statement that the double limit exists when the convergence in one index is uniform in the other.

\subsection{Sums of Nonnegative Families}

A double series $\sum_{m,n}a_{mn}$ has no preferred order of summation. It can be summed by rows, $\sum_m\sum_na_{mn}$, by columns, $\sum_n\sum_ma_{mn}$, by squares or by diagonals, or along any enumeration of $\mathbb N^2$. The clean way to compare these is to define a sum that does not depend on any order, first for nonnegative terms. Let $I$ be a countable set, and $(a_i)_{i\in I}$ a family of numbers $a_i\ge0$. Its \term{sum} is
\[
\sum_{i\in I}a_i\coloneqq\sup\Bigl\{\sum_{i\in F}a_i: F\subseteq I\text{ finite}\Bigr\}\in[0,\infty] ,
\]
a supremum in the extended real line of \cref{def:extended-reals} (\cref{prop:extended-sup}). For $I=\mathbb N$ this is the sum of the series $\sum_ia_i$, which by \vcref{1}{prop:nonnegative-series} is the supremum of the partial sums, the value $\infty$ meaning that the series diverges: every finite set $F\subseteq\mathbb N$ lies in some $\{0,\ldots,n\}$, and the partial sums are sums over finite sets. A series in $[0,\infty]$, whose terms may themselves be $\infty$, is summed in the same way, as the supremum of its partial sums. For two increasing sequences $(x_N)$ and $(y_N)$ in $[0,\infty]$, $\sup_N(x_N+y_N)=\sup_Nx_N+\sup_Ny_N$, since $x_N+y_M\le x_K+y_K$ for $K\ge M,N$; so limits of finite sums of increasing sequences are the sums of the limits, with the convention $x+\infty=\infty$.

\begin{theorem}[Rearranging nonnegative families]
\label{thm:nonnegative-families}
Let $(a_i)_{i\in I}$ be a family of nonnegative numbers indexed by a countable set $I$.
\begin{enumerate}
\item For every bijection $\sigma\colon\mathbb N\to I$, $\sum_{k=0}^\infty a_{\sigma(k)}=\sum_{i\in I}a_i$.
\item For every partition $I=\bigcup_{j\in J}I_j$ of $I$ into disjoint sets, with $J$ countable,
\[
\sum_{i\in I}a_i=\sum_{j\in J}\Bigl(\sum_{i\in I_j}a_i\Bigr) ,
\]
where the outer sum is a sum of numbers in $[0,\infty]$, defined as above.
\item In particular, for a double sequence $a_{mn}\ge0$,
\[
\sum_{m=0}^\infty\sum_{n=0}^\infty a_{mn}=\sum_{n=0}^\infty\sum_{m=0}^\infty a_{mn}=\sum_{(m,n)\in\mathbb N^2}a_{mn}\in[0,\infty] .
\]
\end{enumerate}
\end{theorem}

\begin{proof}
Write $S\coloneqq\sum_{i\in I}a_i$. (i) Each partial sum $\sum_{k\le n}a_{\sigma(k)}$ is a sum over the finite set $\sigma(\{0,\ldots,n\})$, so it is at most $S$. Conversely, a finite set $F\subseteq I$ lies in $\sigma(\{0,\ldots,n\})$ for some $n$, since $\sigma$ is surjective, and so $\sum_Fa_i$ is at most a partial sum. The suprema agree.

(ii) We allow the outer sum over $J$ to be an ordinary sum in $[0,\infty]$, as defined above, also when $J$ is infinite and enumerated as $j_0,j_1,\ldots$; by (i), applied to numbers in $[0,\infty]$, where the proof is the same, its value does not depend on the enumeration. Let $T\coloneqq\sum_j\sum_{I_j}a_i$. A finite set $F\subseteq I$ meets only finitely many $I_j$, and
\[
\sum_{i\in F}a_i=\sum_j\sum_{i\in F\cap I_j}a_i\le\sum_j\sum_{i\in I_j}a_i=T ,
\]
the middle sum having finitely many nonzero terms; so $S\le T$. Conversely, let $j_0,\ldots,j_r$ be finitely many indices, and $F_q\subseteq I_{j_q}$ finite sets. The union of the disjoint sets $F_q$ is finite, so $\sum_{q\le r}\sum_{F_q}a_i\le S$. Since $\sup_F(x_F+c)=\sup_Fx_F+c$ for every family $(x_F)$ in $[0,\infty]$ and every $c\in[0,\infty]$, taking the supremum over each $F_q$ in turn gives $\sum_{q\le r}\sum_{I_{j_q}}a_i\le S$, and the supremum over $r$ gives $T\le S$.

(iii) Apply (ii) to the partition of $\mathbb N^2$ into the rows $\{m\}\times\mathbb N$ and into the columns $\mathbb N\times\{n\}$, together with (i) for the sums over a single row or column, which are ordinary series.
\end{proof}

Part (iii) is the theorem of Tonelli for series, which the note on the Lebesgue integral and the Fubini--Tonelli theorem in \vcref{1}{sec:fubini} described as a consequence of measure theory, sums being integrals with respect to counting measure. Here it is proved directly. No convergence hypothesis is needed: for nonnegative terms, every order of summation gives the same value in $[0,\infty]$, finite or not.

\subsection{Absolutely Summable Families}

A family $(a_i)_{i\in I}$ of complex numbers, indexed by a countable set, is \term{absolutely summable} if $\sum_{i\in I}|a_i|<\infty$. By \cref{thm:nonnegative-families}\,(i) this holds if and only if the series $\sum_ka_{\sigma(k)}$ converges absolutely for one, and then for every, bijection $\sigma\colon\mathbb N\to I$.

\begin{theorem}[Rearranging absolutely summable families]
\label{thm:summable-families}
Let $(a_i)_{i\in I}$ be an absolutely summable family of complex numbers.
\begin{enumerate}
\item The sum $\sum_{i\in I}a_i\coloneqq\sum_{k=0}^\infty a_{\sigma(k)}$ is the same for every bijection $\sigma\colon\mathbb N\to I$, and $|\sum_{i\in I}a_i|\le\sum_{i\in I}|a_i|$.
\item For every partition $I=\bigcup_{j\in J}I_j$ into disjoint sets, with $J$ countable, every subfamily $(a_i)_{i\in I_j}$ is absolutely summable, the family of the sums $s_j\coloneqq\sum_{i\in I_j}a_i$ is absolutely summable, and $\sum_{j\in J}s_j=\sum_{i\in I}a_i$.
\item In particular, for a double sequence with $\sum_{m,n}|a_{mn}|<\infty$, every row and every column converges absolutely, and
\[
\sum_{m=0}^\infty\sum_{n=0}^\infty a_{mn}=\sum_{n=0}^\infty\sum_{m=0}^\infty a_{mn}=\sum_{(m,n)\in\mathbb N^2}a_{mn} .
\]
\end{enumerate}
For finite sets $I$ or $I_j$, the sums are finite sums.
\end{theorem}

\begin{proof}
(i) Two bijections $\sigma,\tau\colon\mathbb N\to I$ give series that are rearrangements of each other, $a_{\tau(k)}=a_{\sigma(\pi(k))}$ with the bijection $\pi\coloneqq\sigma^{-1}\circ\tau$ of $\mathbb N$. For real terms they have the same sum by \vcref{1}{thm:rearrangement-absolute}, and for complex terms the real and imaginary parts are rearranged in the same way. The inequality holds for the partial sums, hence for the limit.

(ii) Suppose first that the $a_i$ are real, and let $a_i^+\coloneqq\max(a_i,0)$ and $a_i^-\coloneqq\max(-a_i,0)$, so that $a_i=a_i^+-a_i^-$ and $0\le a_i^\pm\le|a_i|$. By \cref{thm:nonnegative-families}\,(ii), applied to $|a_i|$, every subfamily is absolutely summable and $\sum_j\sum_{I_j}|a_i|=\sum_I|a_i|<\infty$; since $|s_j|\le\sum_{I_j}|a_i|$ by (i), the family $(s_j)$ is absolutely summable. By the linearity of convergent series, $s_j=\sum_{I_j}a_i^+-\sum_{I_j}a_i^-$, and $\sum_Ia_i=\sum_Ia_i^+-\sum_Ia_i^-$, where all sums of the nonnegative families $a^\pm$ are finite. By \cref{thm:nonnegative-families}\,(ii) for $a^+$ and for $a^-$,
\[
\sum_js_j=\sum_j\sum_{I_j}a_i^+-\sum_j\sum_{I_j}a_i^-=\sum_Ia_i^+-\sum_Ia_i^-=\sum_Ia_i .
\]
For complex terms, apply this to the real and the imaginary parts, which are absolutely summable because $|\operatorname{Re}a_i|,|\operatorname{Im}a_i|\le|a_i|$. (iii) follows from (ii) as in \cref{thm:nonnegative-families}.
\end{proof}

This is the theorem of Fubini for series. Without absolute summability the two iterated sums can both exist and differ, as for the double sequence with the rows
\[
\begin{bmatrix}
1 & -1 & 0 & 0 & \cdots\\
0 & 1 & -1 & 0 & \cdots\\
0 & 0 & 1 & -1 & \cdots\\
\vdots & & & \ddots & \ddots
\end{bmatrix},
\]
that is, $a_{mm}=1$, $a_{m,m+1}=-1$, and $a_{mn}=0$ otherwise. Every row sums to $0$, so $\sum_m\sum_na_{mn}=0$. The column $0$ sums to $1$, every other column to $1-1=0$, so $\sum_n\sum_ma_{mn}=1$. Here $\sum_{m,n}|a_{mn}|=\infty$. The example is the discrete counterpart of the integrand of \vcref{1}{sec:fubini} whose two iterated integrals differ.

Several earlier results are special cases of \cref{thm:summable-families}. With $I=\mathbb N$, part (i) is the rearrangement theorem \vcref{1}{thm:rearrangement-absolute} for complex terms. With the partition of $\mathbb N$ into consecutive blocks, part (ii) is the grouping of terms, which for absolutely convergent series is always allowed. And with $a_{mn}\coloneqq x_my_n$ for absolutely convergent series $\sum x_m$ and $\sum y_n$, the family is absolutely summable, since $\sum_m\sum_n|x_m||y_n|=\sum_m|x_m|\sum_n|y_n|$ by \cref{thm:nonnegative-families}\,(iii); the partition into rows gives the product $(\sum x_m)(\sum y_n)$, and the partition into the finite diagonals $\{(m,n): m+n=k\}$ gives the Cauchy product $\sum_k\sum_{m+n=k}x_my_n$. Thus Cauchy's theorem on products, \vcref{1}{thm:cauchy-product}, follows, and any other arrangement of the products $x_my_n$ gives the same sum as well.

\subsection{Applications}

\Cref{thm:nonnegative-families} sums the values of the zeta function. For $n\ge2$, $\zeta(n)-1=\sum_{k\ge2}k^{-n}$, and summing first over $n$, by geometric series,
\[
\sum_{n=2}^\infty\bigl(\zeta(n)-1\bigr)=\sum_{k=2}^\infty\sum_{n=2}^\infty\frac1{k^n}=\sum_{k=2}^\infty\frac{1/k^2}{1-1/k}=\sum_{k=2}^\infty\frac1{k(k-1)}=1 ,
\]
the last sum telescoping, since $1/(k(k-1))=1/(k-1)-1/k$. In the same way, $\sum_{n\ge1}(\zeta(2n)-1)=\sum_{k\ge2}1/(k^2-1)=3/4$, because $1/(k^2-1)=(1/(k-1)-1/(k+1))/2$ and the telescoping sum is $(1+1/2)/2$. The family $(-1)^nk^{-n}$, $n,k\ge2$, is absolutely summable, by the first computation, and \cref{thm:summable-families} gives
\[
\sum_{n=2}^\infty(-1)^n\bigl(\zeta(n)-1\bigr)=\sum_{k=2}^\infty\frac{1/k^2}{1+1/k}=\sum_{k=2}^\infty\frac1{k(k+1)}=\frac12 .
\]

A classical rearrangement is that of the \term{Lambert series} $\sum_{n\ge1}x^n/(1-x^n)$. For $|x|<1$, expand each term as a geometric series, $x^n/(1-x^n)=\sum_{k\ge1}x^{nk}$. The family $(x^{nk})_{n,k\ge1}$ is absolutely summable, since $\sum_n\sum_k|x|^{nk}=\sum_n|x|^n/(1-|x|^n)\le\sum_n|x|^n/(1-|x|)<\infty$. Partition the index pairs by the value of the product $m=nk$: the pairs with $nk=m$ correspond to the divisors $n$ of $m$, and there are $d(m)$ of them, the number of divisors of $m$. By \cref{thm:summable-families}\,(ii),
\[
\sum_{n=1}^\infty\frac{x^n}{1-x^n}=\sum_{m=1}^\infty d(m)\,x^m\qquad(|x|<1),
\]
a power series that generates the divisor function. Such identities between series and products of functions are the starting point of the analytic theory of numbers~\cite{hardywright2008}.

\section{Power Series Revisited}
\label{sec:power-series-revisited}

\Vcref{1}{sec:power-series} proved the main properties of power series without the notion of uniform convergence, by direct estimates. This section shows how they follow from the general theorems, and then proves three results that Volume~I stated without proof: that the sum of a power series is analytic at every point of its interval of convergence, Abel's theorem on the behaviour at an end of the interval, and Tauber's partial converse of it. The coefficients may be complex; the variable is real, as in Volume~I, and the note on complex power series in \vcref{1}{sec:power-series} applies to everything below. The complex theory of Chapters~10 and~12 extends the results to discs in the complex plane.

\subsection{Uniform Convergence of Power Series}

\begin{proposition}
\label{prop:power-series-uniform}
Let $\sum_ka_k(x-c)^k$ have the radius of convergence $R>0$. For every $r$ with $0<r<R$, the series converges absolutely and uniformly on $[c-r,c+r]$; hence it converges locally uniformly, and uniformly on compact sets, in the interval $(c-R,c+R)$.
\end{proposition}

\begin{proof}
Choose $r<r_0<R$. By \vcref{1}{thm:radius} the terms $a_kr_0^k$ are bounded, say by $C$, and by \vcref{1}{lem:abel}, $|a_k(x-c)^k|\le C(r/r_0)^k$ for $|x-c|\le r$. The M-test with $M_k=C(r/r_0)^k$ applies. Every point of $(c-R,c+R)$ lies in the interior of some $[c-r,c+r]$ with $r<R$, which gives local uniformity, and \cref{prop:locally-uniform} gives uniformity on compact sets.
\end{proof}

The convergence need not be uniform on the whole interval of convergence: the geometric series has the partial sums $(1-x^{n+1})/(1-x)$, which are bounded on $(-1,1)$, while the sum $1/(1-x)$ is not. With the proposition, the theorems of \vcref{1}{sec:power-series} become special cases of this chapter. The sum $f$ of the series is continuous on $(c-R,c+R)$ by \cref{cor:series-termwise}\,(i). The series of derivatives $\sum_kka_k(x-c)^{k-1}$ has the same radius $R$, by \vcref{1}{lem:derived-radius}, so it converges uniformly on every $[c-r,c+r]$, and \cref{cor:series-termwise}\,(iii) shows that $f$ is differentiable there with $f'(x)=\sum_kka_k(x-c)^{k-1}$, which is \vcref{1}{thm:termwise-diff}. Integrating term by term over $[c,x]$, by \cref{cor:series-termwise}\,(ii), gives \vcref{1}{cor:termwise-int}. The direct proofs of Volume~I have the advantage of using nothing about uniform convergence; the present ones show that nothing beyond uniform convergence on smaller intervals is needed.

\subsection{Sums of Power Series Are Analytic}

A function is \term{analytic} at a point $b$ if it is the sum of a power series with centre $b$ and a positive radius on some interval around $b$, as in \vcref{1}{def:taylor-series} and \vcref{2}{def:analytic}; it is analytic on an open interval if it is analytic at every point of it. A power series is analytic at its centre by definition. The paragraph after \vcref{1}{thm:taylor-convergence} asserted that it is analytic at every point of its interval of convergence, by a rearrangement of a double series. With \cref{thm:summable-families} the proof is short.

\begin{theorem}[Changing the centre]
\label{thm:change-centre}
Let $f(x)=\sum_{k=0}^\infty a_k(x-c)^k$ have the radius of convergence $R>0$, with $a_k\in\mathbb F$, and let $|b-c|<R$. Then for $|x-b|<R-|b-c|$,
\[
f(x)=\sum_{j=0}^\infty b_j(x-b)^j,\qquad b_j=\sum_{k=j}^\infty\binom kja_k(b-c)^{k-j}=\frac{f^{(j)}(b)}{j!},
\]
where the series defining $b_j$ converge absolutely, and the radius of convergence of the series with centre $b$ is at least $R-|b-c|$. In particular, $f$ is analytic on $(c-R,c+R)$.
\end{theorem}

\begin{proof}
Let $|x-b|<R-|b-c|$, so that $\rho\coloneqq|x-b|+|b-c|<R$. By the binomial theorem,
\[
a_k(x-c)^k=a_k\bigl((x-b)+(b-c)\bigr)^k=\sum_{j=0}^k\binom kja_k(b-c)^{k-j}(x-b)^j .
\]
Let $t_{kj}$ be the term with index $j$ of the last sum, and consider the family of these terms, indexed by the pairs $(k,j)$ with $0\le j\le k$. Summing the absolute values first over $j$, by the binomial theorem again,
\[
\sum_{k=0}^\infty\sum_{j=0}^k|t_{kj}|=\sum_{k=0}^\infty|a_k|\bigl(|x-b|+|b-c|\bigr)^k=\sum_{k=0}^\infty|a_k|\rho^k<\infty ,
\]
because the series converges absolutely at distance $\rho<R$ from its centre (\vcref{1}{thm:radius}). So the family is absolutely summable, by \cref{thm:nonnegative-families}, and by \cref{thm:summable-families} its sum can be computed by the partition into the sets with fixed $k$, which gives $f(x)$, or into the sets with fixed $j$, which gives $\sum_jb_j(x-b)^j$; the series defining $b_j$ converges absolutely as a sum over one of the parts. Thus the power series $\sum_jb_j(x-b)^j$ converges at every $x$ with $|x-b|<R-|b-c|$, and its radius is at least $R-|b-c|$. Its sum is $f$ on the interval $|x-b|<R-|b-c|$, so by \vcref{1}{thm:termwise-diff}, applied to this series at its centre, $b_j=f^{(j)}(b)/j!$.
\end{proof}

The bound $R-|b-c|$ for the new radius can be strict. The geometric series $\sum x^k$, with $c=0$ and $R=1$, has the sum $1/(1-x)$, and for $|b|<1$
\[
\frac1{1-x}=\frac1{(1-b)-(x-b)}=\sum_{j=0}^\infty\frac{(x-b)^j}{(1-b)^{j+1}}\qquad(|x-b|<1-b),
\]
with the radius $1-b$, which for $b<0$ exceeds $1-|b|$; the new interval of convergence reaches beyond the old one, to the left. The radius at $b$ is the distance from $b$ to the singularity at $1$. For $1/(1+x^2)$, whose series at $0$ has radius $1$, the radius at $b$ turns out to be $\sqrt{1+b^2}$, the distance from $b$ to the complex points $\pm i$, as the note on complex power series in \vcref{1}{sec:power-series} explained; Chapter~12 proves the general statement. Extending a function beyond the interval of convergence of one power series by the power series at other centres is the process of \term{analytic continuation}, which Chapter~12 studies in the complex plane.

\subsection{The Identity Theorem}

An analytic function is determined by very little information: by its derivatives at one point, or by its values on a set with an accumulation point. This is the property that makes analytic functions rigid, in contrast with the smooth function $e^{-1/x^2}$ of \vcref{1}{prop:flat-function}, which vanishes with all its derivatives at $0$ without vanishing anywhere else.

\begin{theorem}[Identity theorem]
\label{thm:identity}
Let $I$ be an open interval and $f\colon I\to\mathbb F$ analytic on $I$. The following are equivalent:
\begin{enumerate}
\item the set of zeros of $f$ has an accumulation point in $I$;
\item there is a point $a\in I$ with $f^{(k)}(a)=0$ for all $k\in\mathbb N$;
\item $f=0$ on $I$.
\end{enumerate}
\end{theorem}

\begin{proof}
An analytic function is infinitely differentiable, by \vcref{1}{thm:termwise-diff} applied near each point, and near each point it is the sum of its Taylor series. (iii) implies (i), since every point of an open interval is an accumulation point of it. (i)$\Rightarrow$(ii): let $a\in I$ be an accumulation point of the zeros, and $x_j\to a$ zeros of $f$ with $x_j\ne a$ (\cref{prop:closure-sequences}). Near $a$, $f$ is the sum of a power series with centre $a$, which vanishes at the points $x_j$ that lie close enough to $a$; by \vcref{1}{cor:power-series-unique}, all its coefficients vanish, and these are $f^{(k)}(a)/k!$. (ii)$\Rightarrow$(iii): let $U$ be the set of points $x\in I$ with $f^{(k)}(x)=0$ for all $k$. It is nonempty by (ii). It is open: if $x\in U$, the Taylor series of $f$ at $x$ is the zero series, so $f=0$ on an interval around $x$, and all derivatives vanish there. It is closed in $I$, as the intersection of the closed sets $\{x\in I: f^{(k)}(x)=0\}$, which are closed because the derivatives are continuous (\cref{thm:continuity-characterisations}). Since $I$ is connected (\cref{thm:connected-real}), $U=I$, and $f=0$ on $I$.
\end{proof}

\begin{corollary}
\label{cor:identity}
Let $I$ be an open interval.
\begin{enumerate}
\item If $f$ is analytic on $I$ and not identically zero, its zeros are isolated, and every closed bounded subinterval of $I$ contains only finitely many of them.
\item If $f$ and $g$ are analytic on $I$ and agree on a set that has an accumulation point in $I$, or if $f^{(k)}(a)=g^{(k)}(a)$ for all $k$ at one point $a\in I$, then $f=g$ on $I$.
\end{enumerate}
\end{corollary}

\begin{proof}
(i) By the theorem, the zeros have no accumulation point in $I$, so each zero is isolated, and infinitely many zeros in $[\alpha,\beta]\subseteq I$ would have an accumulation point there, by the Bolzano--Weierstrass theorem (\vcref{1}{thm:bolzano-weierstrass}). (ii) Apply the theorem to $f-g$, which is analytic, since the difference of two power series with the same centre is a power series whose radius is at least the smaller radius.
\end{proof}

The accumulation point must lie in $I$: the theorem says nothing about zeros that accumulate only at an end of $I$. The connectedness of $I$ is essential as well, since on a union of two disjoint intervals a function may vanish on one and not on the other. The proof is the first appearance of a pattern that the complex analysis of Part~II uses constantly: a property that holds on a nonempty set that is open and closed holds on a whole connected set. By \cref{thm:change-centre} the corollary applies to sums of power series on their intervals of convergence; for instance, if the sums of two power series with centre $0$ agree on $(0,\delta)$, the series are equal.

\subsection{Abel's Theorem}

At an end of the interval of convergence a power series may converge or diverge, and \vcref{1}{sec:power-series} gave examples of all possibilities. When it converges, the sum is continuous there from the inside. This is the theorem of Abel that the note on complex power series in \vcref{1}{sec:power-series} stated.

\begin{theorem}[Abel's theorem]
\label{thm:abel}
Let $a_k\in\mathbb C$ and $R\in(0,\infty)$, and suppose that $\sum_ka_kR^k$ converges. Then $\sum_ka_kx^k$ converges uniformly on $[0,R]$, its sum $f$ is continuous on $[0,R]$, and
\[
\lim_{x\to R^-}\sum_{k=0}^\infty a_kx^k=\sum_{k=0}^\infty a_kR^k .
\]
The same holds at $-R$, on $[-R,0]$, if $\sum_ka_k(-R)^k$ converges.
\end{theorem}

\begin{proof}
For $x\in[0,R]$ write $a_kx^k=(a_kR^k)(x/R)^k$. The constant functions $a_kR^k$ form a series that converges, uniformly on $[0,R]$ since it does not depend on $x$, and for each $x\in[0,R]$ the sequence $(x/R)^k$ is decreasing and bounded by $1$. By Abel's test, \cref{thm:abel-dirichlet}\,(ii), $\sum_ka_kx^k$ converges uniformly on $[0,R]$, and its sum is continuous by \cref{cor:series-termwise}\,(i), since the partial sums are polynomials. The statement at $-R$ follows by applying this to the coefficients $(-1)^ka_k$.
\end{proof}

So the sum of a power series is continuous on the whole set where the series converges, the interval of convergence together with whichever of its ends belong to it. Here are three applications. In each, the value of the series at the end is computed from a formula valid inside the interval, without an estimate of the remainder at the end.
\begin{enumerate}
\item For $|x|<1$, $\ln(1+x)=\sum_{k\ge1}(-1)^{k+1}x^k/k$, by integrating the geometric series term by term (\vcref{1}{cor:termwise-int,thm:ln-lambda}). At $x=1$ the series converges by Leibniz's test, so by Abel's theorem and the continuity of the logarithm,
\[
\ln2=\lim_{x\to1^-}\ln(1+x)=1-\frac12+\frac13-\frac14+\cdots .
\]
\item In the same way, $\arctan x=\sum_{k\ge0}(-1)^kx^{2k+1}/(2k+1)$ for $|x|<1$, and the series converges at $x=1$, so $\pi/4=\arctan1=1-1/3+1/5-\cdots$.
\item By the binomial series, \vcref{1}{thm:binomial-series} with $r=-1/2$, $(1+x)^{-1/2}=\sum_{k\ge0}(-1)^kc_kx^k$ for $|x|<1$, where, as computed in \vcref{1}{sec:expansions},
\[
c_k\coloneqq\binom{2k}k\frac1{4^k}=\prod_{j=1}^k\frac{2j-1}{2j} .
\]
The numbers $c_k$ decrease, since $c_{k+1}/c_k=(2k+1)/(2k+2)$, and they tend to $0$: since $(2j-1)/(2j)\le2j/(2j+1)$,
\[
c_k^2=\prod_{j=1}^k\Bigl(\frac{2j-1}{2j}\Bigr)^2\le\prod_{j=1}^k\frac{2j-1}{2j}\cdot\frac{2j}{2j+1}=\frac1{2k+1} .
\]
So the series converges at $x=1$ by Leibniz's test, and Abel's theorem gives $\sum_{k\ge0}(-1)^kc_k=1/\sqrt2$.
\end{enumerate}
Volume~I proved (i) and (ii) differently, in \vcref{1}{thm:log-series,thm:arctan-series}, by integrating the finite geometric sum with its exact remainder. That argument gives more, an explicit bound for the error of each partial sum at $x=1$, which Abel's theorem does not provide; Abel's theorem gives less, but it needs nothing beyond the convergence of the series at the end.

The converse of Abel's theorem is false. The geometric series $\sum_k(-1)^kx^k$ has the sum $1/(1+x)$, which tends to $1/2$ as $x\to1^-$, while $\sum_k(-1)^k$ diverges. The existence of $\lim_{x\to R^-}f(x)$ does not imply convergence at $R$, and the second half of this section studies what it does imply. For complex $z$, Abel's theorem extends to $z\to R$ within any sector $\{z: |R-z|\le K(R-|z|)\}$ with $K\ge1$, a so-called Stolz angle, but not to arbitrary approach from inside the disc~\cite{knopp1990}.

\subsection{Products of Series}

The Cauchy product of two convergent series need not converge, as Cauchy's example in \vcref{1}{sec:absolute} shows. When it does converge, its sum is the right one.

\begin{theorem}[Abel's theorem on products]
\label{thm:abel-product}
Let $\sum_ka_k=A$, $\sum_kb_k=B$ and their Cauchy product $\sum_nc_n$, $c_n=\sum_{k=0}^na_kb_{n-k}$, all converge, with complex terms, and let $\sum_nc_n=C$. Then $C=AB$.
\end{theorem}

\begin{proof}
The terms of the three series are bounded, since they tend to $0$, so the three power series $\sum a_kx^k$, $\sum b_kx^k$ and $\sum c_nx^n$ have radii at least $1$ and converge absolutely for $0\le x<1$. Their sums $A(x)$, $B(x)$ and $C(x)$ satisfy $C(x)=A(x)B(x)$, because the Cauchy product of the absolutely convergent series $\sum a_kx^k$ and $\sum b_kx^k$ is $\sum c_nx^n$ (\vcref{1}{lem:complex-series}). As $x\to1^-$, Abel's theorem gives $A(x)\to A$, $B(x)\to B$ and $C(x)\to C$; hence $C=AB$.
\end{proof}

The theorem is in Abel's paper of 1826~\cite{abel1826}. For an example, square the alternating harmonic series $\sum_{k\ge0}(-1)^k/(k+1)=\ln2$. By partial fractions, $1/((k+1)(n-k+1))=(1/(k+1)+1/(n-k+1))/(n+2)$, so the Cauchy product has the terms
\[
c_n=(-1)^n\sum_{k=0}^n\frac1{(k+1)(n-k+1)}=(-1)^n\frac{2H_{n+1}}{n+2},\qquad H_m\coloneqq1+\frac12+\cdots+\frac1m .
\]
The numbers $H_{n+1}/(n+2)$ decrease, since $(n+3)H_{n+1}-(n+2)H_{n+2}=H_{n+1}-1\ge0$, and they tend to $0$, since $H_{n+1}\le1+\ln(n+1)$ by \vcref{1}{thm:integral-test}\,(iii); by Leibniz's test the product series converges, and the theorem gives
\[
\sum_{n=0}^\infty(-1)^n\frac{2H_{n+1}}{n+2}=(\ln2)^2 .
\]
Neither factor converges absolutely, so \vcref{1}{thm:cauchy-product} and Mertens's theorem do not apply.

\subsection{Abel and Ces\`aro Summation}

The note on sums of divergent series in \vcref{1}{sec:asymptotic-expansions} described two methods that assign sums to some divergent series. A series $\sum_ka_k$ with partial sums $s_n$ is \term{Ces\`aro summable} to $s$ if the means $\sigma_n\coloneqq(s_0+\cdots+s_n)/(n+1)$ tend to $s$, and \term{Abel summable} to $s$ if $\sum_ka_kx^k$ converges for $0\le x<1$ and its sum tends to $s$ as $x\to1^-$. A convergent series is Ces\`aro summable to its sum, by \cref{cor:cesaro}, and Abel summable to its sum, by Abel's theorem. The two methods are not independent.

\begin{theorem}[Ces\`aro summability implies Abel summability]
\label{thm:cesaro-abel}
If $\sum_ka_k$ is Ces\`aro summable to $s$, then it is Abel summable to $s$.
\end{theorem}

\begin{proof}
Let $|\sigma_n|\le M$ for all $n$. Then $|s_n|=|(n+1)\sigma_n-n\sigma_{n-1}|\le(2n+1)M$ and $|a_n|=|s_n-s_{n-1}|\le4nM+M$, with $s_{-1}\coloneqq0$ and $\sigma_{-1}\coloneqq0$. These bounds grow at most linearly in $n$, so the power series $\sum a_nx^n$, $\sum s_nx^n$ and $\sum(n+1)\sigma_nx^n$ converge absolutely for $|x|<1$, by comparison with $\sum(n+1)|x|^n$. Their sums are related by
\[
\sum_{n=0}^\infty a_nx^n=(1-x)\sum_{n=0}^\infty s_nx^n=(1-x)^2\sum_{n=0}^\infty(n+1)\sigma_nx^n\qquad(|x|<1),
\]
because the Cauchy product of $\sum a_nx^n$ with $\sum x^n=1/(1-x)$ is $\sum s_nx^n$, and that of $\sum s_nx^n$ with $\sum x^n$ is $\sum(n+1)\sigma_nx^n$ (\vcref{1}{thm:cauchy-product,lem:complex-series}). Since also $(1-x)^2\sum_n(n+1)x^n=1$, by \vcref{1}{sec:power-series}, for $0\le x<1$
\[
f(x)-s\coloneqq\sum_na_nx^n-s=(1-x)^2\sum_{n=0}^\infty(n+1)(\sigma_n-s)x^n .
\]
Given $\varepsilon>0$, choose $N$ with $|\sigma_n-s|<\varepsilon$ for $n>N$. Then
\[
|f(x)-s|\le(1-x)^2\sum_{n=0}^N(n+1)|\sigma_n-s|+\varepsilon(1-x)^2\sum_{n=0}^\infty(n+1)x^n\le(1-x)^2C_N+\varepsilon ,
\]
with $C_N\coloneqq\sum_{n\le N}(n+1)|\sigma_n-s|$, and the right side is less than $2\varepsilon$ for $x$ close to $1$.
\end{proof}

The theorem is due to Frobenius~\cite{frobenius1880}. The series $1-1+1-\cdots$ has the Ces\`aro means $\sigma_n$ equal to $1/2$ or $(n+2)/(2(n+1))$, which tend to $1/2$, and the Abel sum $\lim1/(1+x)=1/2$. The converse of the theorem fails: $\sum_k(-1)^k(k+1)=1-2+3-\cdots$ is Abel summable to $1/4$, since $\sum_k(-1)^k(k+1)x^k=1/(1+x)^2$ by differentiating the geometric series, but its partial sums $1,-1,2,-2,3,\ldots$ have the means $\sigma_{2m+1}=0$ and $\sigma_{2m}=(m+1)/(2m+1)\to1/2$, which do not converge. Abel's method thus sums strictly more series than Ces\`aro's. Chapter~9 uses Ces\`aro means of Fourier series, in Fej\'er's theorem, and Abel means, through the Poisson kernel, and Chapter~14 returns to the second.

\subsection{Tauber's Theorem}

A series that is Abel summable converges if its terms are small enough. Results of this kind, which recover convergence from summability under a condition on the terms, are called \term{Tauberian theorems}, after the first of them.

\begin{theorem}[Tauber]
\label{thm:tauber}
Let $\sum_ka_kx^k$ converge for $0\le x<1$, with a sum $f(x)$ that tends to $s$ as $x\to1^-$, and suppose that $ka_k\to0$. Then $\sum_ka_k$ converges, and its sum is $s$.
\end{theorem}

\begin{proof}
For $0<x<1$ and $n\ge1$,
\[
s_n-f(x)=\sum_{k=1}^na_k(1-x^k)-\sum_{k>n}a_kx^k ,
\]
where $s_n$ is the partial sum. Since $1-x^k=(1-x)(1+x+\cdots+x^{k-1})\le k(1-x)$, the first sum is at most $(1-x)\sum_{k\le n}k|a_k|$ in absolute value. With $\varepsilon_n\coloneqq\sup_{k>n}k|a_k|$, which tends to $0$, the second is at most
\[
\sum_{k>n}\frac{k|a_k|}kx^k\le\frac{\varepsilon_n}n\sum_{k>n}x^k\le\frac{\varepsilon_n}{n(1-x)} .
\]
Take $x=x_n\coloneqq1-1/n$ for $n\ge2$. Then
\[
|s_n-f(x_n)|\le\frac1n\sum_{k=1}^nk|a_k|+\varepsilon_n .
\]
The first term on the right is the $n$th Ces\`aro mean of the sequence $(k|a_k|)_{k\ge1}$, which tends to $0$, so it tends to $0$ by \cref{cor:cesaro}. Hence $s_n-f(x_n)\to0$, and since $x_n\to1^-$, $f(x_n)\to s$ and $s_n\to s$.
\end{proof}

Tauber proved the theorem in 1897~\cite{tauber1897}. The condition $ka_k\to0$ excludes series like $1-1+1-\cdots$, whose terms do not tend to zero at all, but it is far from necessary for convergence, and the right condition was found later, as the following note describes.

\begin{note}{Tauberian theorems}
Littlewood proved in 1911 that the condition of \cref{thm:tauber} can be weakened to $|ka_k|\le C$ for a constant $C$, so that, for instance, an Abel summable series with $|a_k|\le1/k$ converges~\cite{littlewood1911}. His proof is much harder than Tauber's. For series with nonnegative terms, Hardy and Littlewood proved that $\sum_ka_kx^k\sim1/(1-x)$ as $x\to1^-$ implies $a_0+\cdots+a_n\sim n$, a statement that relates the growth of a power series near the boundary to the growth of its partial sums and that is used in the analytic theory of numbers; Karamata found a short proof of it in 1930, and Littlewood's theorem can be derived from it. Hardy's book presents these results and the whole theory~\cite{hardy1949divergent}. In the same spirit, a Ces\`aro summable series with $ka_k$ bounded converges, a theorem of Hardy~\cite{hardy1949divergent}. Wiener's general Tauberian theorem of 1932, which belongs to Fourier analysis, contains many of these results as special cases; Korevaar's book surveys the theory and its applications~\cite{korevaar2004}.
\end{note}

\section{Continuous Nowhere Differentiable Functions}
\label{sec:nowhere-differentiable}

The functions $\sqrt{x^2+1/n}$ showed that a uniform limit of smooth functions can fail to be differentiable at a point. Uniform limits can do much worse: they can fail to be differentiable at every point. The note on continuous nowhere differentiable functions in \vcref{1}{sec:derivative} described the example that Weierstrass presented in 1872. This section proves first that a simpler function, built from the distance to the nearest integer, has no derivative anywhere, and then that Weierstrass's function, under his original conditions, has none either.

\subsection{The Takagi Function}

Let $\varphi(x)\coloneqq\operatorname{dist}(x,\mathbb Z)=\min_{m\in\mathbb Z}|x-m|$ be the distance from $x$ to the nearest integer. It is continuous, periodic with the period $1$, takes values in $[0,1/2]$, and it is linear with slope $1$ on each interval $[m,m+1/2]$ and with slope $-1$ on each interval $[m+1/2,m+1]$, $m\in\mathbb Z$. The \term{Takagi function} is
\begin{equation}
T(x)\coloneqq\sum_{n=0}^\infty\frac{\varphi(2^nx)}{2^n}\qquad(x\in\mathbb R).
\label{eq:takagi}
\end{equation}
The $n$th term is a sawtooth of height $2^{-n-1}$ with $2^n$ teeth in each interval of length $1$. By the M-test with $M_n=2^{-n-1}$, the series converges uniformly on $\mathbb R$, so $T$ is continuous, periodic with the period $1$, and $0\le T\le1$. Its partial sums $T_N\coloneqq\sum_{n\le N}2^{-n}\varphi(2^n\cdot)$ are piecewise linear, with corners at the multiples of $2^{-N-1}$, and $|T-T_N|\le2^{-N-1}$. \Cref{fig:nowhere-differentiable}A shows $T$ with its first partial sums. The slopes of $T_N$ are sums of $N+1$ terms $\pm1$, so they grow without bound as $N$ increases, while the teeth added at each step become flatter in height.

The proof that $T$ is nowhere differentiable uses the slopes of chords over short intervals that contain the point, rather than the difference quotients from the point itself.

\begin{lemma}[Chords around a point]
\label{lem:chords}
Let $f$ be defined on an interval $I$ and differentiable at $x\in I$. If $u_m\le x\le v_m$ are points of $I$ with $u_m<v_m$ and $v_m-u_m\to0$, then
\[
\frac{f(v_m)-f(u_m)}{v_m-u_m}\to f'(x)\qquad(m\to\infty).
\]
\end{lemma}

\begin{proof}
If $u_m<x<v_m$, then with $\lambda\coloneqq(v_m-x)/(v_m-u_m)\in(0,1)$,
\[
\frac{f(v_m)-f(u_m)}{v_m-u_m}=\lambda\,\frac{f(v_m)-f(x)}{v_m-x}+(1-\lambda)\,\frac{f(x)-f(u_m)}{x-u_m} ,
\]
a convex combination of two difference quotients at $x$, as one checks by multiplying out. Its distance from $f'(x)$ is at most the larger of the distances of the two quotients from $f'(x)$, and these tend to $0$ since $v_m-x\to0$ and $x-u_m\to0$. If $u_m=x$ or $v_m=x$, the chord slope is itself one of the two quotients.
\end{proof}

The lemma is false for chords that do not contain the point: for $f(x)=x^2\sin(1/x)$, with $f(0)=0$, which is differentiable at $0$ with $f'(0)=0$ (\vcref{1}{prop:x2sin}), the chords between the points $1/(2\pi m+\pi/2)$ and $1/(2\pi m-\pi/2)$, both on the same side of $0$, have slopes that do not tend to $0$.

\begin{theorem}[Takagi]
\label{thm:takagi}
The Takagi function is continuous on $\mathbb R$ and differentiable at no point.
\end{theorem}

\begin{proof}
Continuity was shown above. Fix $x\in\mathbb R$, and for $m\ge1$ let $u_m\coloneqq k2^{-m}$ and $v_m\coloneqq(k+1)2^{-m}$, where $k\coloneqq\lfloor2^mx\rfloor$, so that $u_m\le x<v_m$ and $v_m-u_m=2^{-m}$. Consider the terms of \eqref{eq:takagi} at $u_m$ and $v_m$.

For $n\ge m$, the numbers $2^nu_m=2^{n-m}k$ and $2^nv_m=2^{n-m}(k+1)$ are integers, so $\varphi(2^nu_m)=\varphi(2^nv_m)=0$. For $n<m$, the points $2^nu_m$ and $2^nv_m$ are consecutive multiples of $2^{n-m}$, and $2^{n-m}\le1/2$; so the interval between them contains no multiple of $1/2$ in its interior, and on it $\varphi$ is linear with a slope $\varepsilon_n\in\{1,-1\}$. Hence
\[
\frac{\varphi(2^nv_m)-\varphi(2^nu_m)}{2^n(v_m-u_m)}=\varepsilon_n\qquad(n<m),
\]
and the chord slope of $T$ is
\[
D_m\coloneqq\frac{T(v_m)-T(u_m)}{v_m-u_m}=\sum_{n=0}^{m-1}\frac{\varphi(2^nv_m)-\varphi(2^nu_m)}{2^n(v_m-u_m)}=\sum_{n=0}^{m-1}\varepsilon_n .
\]
The sign $\varepsilon_n$ is determined by which half of an interval $[j/2^n,(j+1)/2^n]$ the point $x$ lies in, and it does not depend on $m>n$, since the intervals $[u_m,v_m]$ are nested. So $D_{m+1}-D_m=\varepsilon_m=\pm1$ for every $m$, and the sequence $(D_m)$ is not a Cauchy sequence. If $T$ were differentiable at $x$, then $D_m\to T'(x)$ by \cref{lem:chords}, which is impossible.
\end{proof}

The proof shows more than nondifferentiability: the chord slopes $D_m$ are integers with the parity of $m$, and they perform a walk with steps $\pm1$ that is determined by the binary digits of $x$, the step $\varepsilon_n$ being $+1$ if the $(n+1)$st binary digit of $x$ after the point is $0$ and $-1$ if it is $1$. For $x=1/3=0.0101\ldots$ in binary, the chord slopes are $1,0,1,0,\ldots$, bounded but not convergent; for most $x$ they wander off, and \cref{sec:uniform-computation} follows such a walk. At $x=1/3$, $\varphi(2^n/3)=1/3$ for every $n$, so $T(1/3)=2/3$, which is the largest value of $T$~\cite{allaart2012}. Takagi gave the function in 1903~\cite{takagi1903}, and van der Waerden rediscovered it in 1930 with the base $10$ in place of $2$~\cite{vanderwaerden1930}.

\subsection{Weierstrass's Function}

For $0<a<1$ and $b>0$, the series
\begin{equation}
W(x)\coloneqq\sum_{n=0}^\infty a^n\cos(b^n\pi x)
\label{eq:weierstrass}
\end{equation}
converges uniformly on $\mathbb R$ by the M-test with $M_n=a^n$, and its sum is continuous; this proves the continuity claimed in the note in \vcref{1}{sec:derivative}. The series of derivatives, $-\pi\sum_n(ab)^n\sin(b^n\pi x)$, converges uniformly if $ab<1$, by the M-test with $M_n=\pi(ab)^n$, and then $W$ is continuously differentiable by \cref{cor:series-termwise}\,(iii). If $ab\ge1$, the terms of the series of derivatives do not tend to $0$, and the theorems of this chapter say nothing. In fact differentiability then fails everywhere.

\begin{theorem}[Weierstrass]
\label{thm:weierstrass-function}
Let $0<a<1$, let $b$ be an odd integer, and suppose that $ab>1+3\pi/2$. Then the function $W$ of \eqref{eq:weierstrass} is continuous on $\mathbb R$ and differentiable at no point. More precisely, at every point $x$ there are points $y_m<x<z_m$ with $y_m,z_m\to x$ such that the difference quotients $(W(y_m)-W(x))/(y_m-x)$ and $(W(z_m)-W(x))/(z_m-x)$ have opposite signs and absolute values that tend to $\infty$.
\end{theorem}

\begin{proof}
Fix $x$ and $m\ge1$. Let $\alpha_m$ be the integer nearest to $b^mx$, in the sense that $b^mx=\alpha_m+\xi_m$ with $-1/2\le\xi_m<1/2$, and let
\[
y_m\coloneqq\frac{\alpha_m-1}{b^m},\qquad z_m\coloneqq\frac{\alpha_m+1}{b^m} .
\]
Then $x-y_m=(1+\xi_m)/b^m\in[1/(2b^m),3/(2b^m))$ and $z_m-x=(1-\xi_m)/b^m\in(1/(2b^m),3/(2b^m)]$, so $y_m<x<z_m$ and both points tend to $x$, since $b\ge ab>1$. We split the difference quotient at $y_m$ into the terms with $n<m$ and those with $n\ge m$:
\[
\frac{W(y_m)-W(x)}{y_m-x}=P_m+Q_m,\qquad P_m\coloneqq\sum_{n=0}^{m-1}a^n\frac{\cos(b^n\pi y_m)-\cos(b^n\pi x)}{y_m-x},
\]
and $Q_m$ is the same sum over $n\ge m$.

For $n<m$, the mean value theorem and $|\sin|\le1$ give $|\cos(b^n\pi y)-\cos(b^n\pi x)|\le b^n\pi|y-x|$, so
\[
|P_m|\le\pi\sum_{n=0}^{m-1}(ab)^n=\pi\frac{(ab)^m-1}{ab-1}<\frac{\pi}{ab-1}(ab)^m .
\]
For $n\ge m$, $b^{n-m}$ is an odd integer, since $b$ is odd. Hence $b^n\pi y_m=b^{n-m}(\alpha_m-1)\pi$ is an integer multiple of $\pi$ with the parity of $\alpha_m-1$, and $\cos(b^n\pi y_m)=(-1)^{\alpha_m-1}=-(-1)^{\alpha_m}$. Similarly $b^n\pi x=b^{n-m}\alpha_m\pi+b^{n-m}\xi_m\pi$, and by the addition theorem, since $\sin(b^{n-m}\alpha_m\pi)=0$ and $\cos(b^{n-m}\alpha_m\pi)=(-1)^{\alpha_m}$,
\[
\cos(b^n\pi x)=(-1)^{\alpha_m}\cos(b^{n-m}\pi\xi_m) .
\]
With $y_m-x=-(1+\xi_m)/b^m$, therefore,
\[
Q_m=(-1)^{\alpha_m}b^m\sum_{n=m}^\infty a^n\frac{1+\cos(b^{n-m}\pi\xi_m)}{1+\xi_m} .
\]
All terms of the last sum are nonnegative, and the term with $n=m$ is at least $a^m/(3/2)$, because $\cos(\pi\xi_m)\ge0$ for $|\xi_m|\le1/2$ and $0<1+\xi_m<3/2$. So $Q_m=2(-1)^{\alpha_m}\eta_m(ab)^m/3$ with $\eta_m\ge1$, and
\[
\frac{W(y_m)-W(x)}{y_m-x}=(-1)^{\alpha_m}(ab)^m\Bigl(\frac23\eta_m+\theta_m\frac\pi{ab-1}\Bigr),\qquad|\theta_m|<1 .
\]
The same computation at $z_m$, where $\cos(b^n\pi z_m)=(-1)^{\alpha_m+1}$ and $z_m-x=(1-\xi_m)/b^m$ with $0<1-\xi_m\le3/2$, gives
\[
\frac{W(z_m)-W(x)}{z_m-x}=-(-1)^{\alpha_m}(ab)^m\Bigl(\frac23\eta_m'+\theta_m'\frac\pi{ab-1}\Bigr),\qquad\eta_m'\ge1,\ |\theta_m'|<1 .
\]
The hypothesis $ab>1+3\pi/2$ says exactly that $\pi/(ab-1)<2/3$. So both brackets are at least $2/3-\pi/(ab-1)>0$, the two quotients have opposite signs, and their absolute values are at least $(2/3-\pi/(ab-1))(ab)^m$, which tends to $\infty$. If $W$ were differentiable at $x$, both quotients would tend to $W'(x)$, which is impossible.
\end{proof}

The proof follows Weierstrass's own, as published by du Bois-Reymond~\cite{duboisreymond1875}. The condition $ab>1+3\pi/2$ is an artefact of the method, which controls the first $m$ terms by a crude bound. Hardy proved in 1916 that $W$ is nowhere differentiable whenever $0<a<1$, $b>1$ and $ab\ge1$, without the assumption that $b$ is an integer~\cite{hardy1916}; his proof is considerably harder. \Cref{fig:nowhere-differentiable}B shows $W$ for $a=1/2$ and $b=3$, a case of Hardy's theorem but not of \cref{thm:weierstrass-function}. Together with the condition $ab<1$ for continuous differentiability, Hardy's theorem shows that the transition is sharp: the function is either continuously differentiable or differentiable nowhere.

\begin{figure}
\sdfloat{%
\sdpanel{A}{\begin{tikzpicture}
\begin{axis}[sdplot, width=0.47\textwidth, height=0.36\textwidth, xmin=0, xmax=1.02, ymin=0, ymax=0.98, xtick={0,0.25,0.5,0.75,1}, xticklabels={$0$,$1/4$,$1/2$,$3/4$,$1$}, ytick={0,0.5,0.6667}, yticklabels={$0$,$1/2$,$2/3$}, xlabel={$x$}, legend columns=2, legend style={at={(0.97,0.99)}, anchor=north east}]
\addplot[black, densely dotted, line width=0.7pt] coordinates {(0,0) (0.5,0.5) (1,0)};
\addlegendentry{$T_0$}
\addplot[black, dashed, line width=0.6pt] coordinates {(0,0) (0.25,0.5) (0.5,0.5) (0.75,0.5) (1,0)};
\addlegendentry{$T_1$}
\addplot[black!45, line width=0.6pt] coordinates {(0,0) (0.125,0.375) (0.25,0.5) (0.375,0.625) (0.5,0.5) (0.625,0.625) (0.75,0.5) (0.875,0.375) (1,0)};
\addlegendentry{$T_2$}
\addplot[black, line width=0.5pt] coordinates {(0.0000,0.0000) (0.0020,0.0176) (0.0039,0.0312) (0.0059,0.0449) (0.0078,0.0547) (0.0098,0.0684) (0.0117,0.0781) (0.0137,0.0879) (0.0156,0.0938) (0.0176,0.1074) (0.0195,0.1172) (0.0215,0.1270) (0.0234,0.1328) (0.0254,0.1426) (0.0273,0.1484) (0.0293,0.1543) (0.0312,0.1562) (0.0332,0.1699) (0.0352,0.1797) (0.0371,0.1895) (0.0391,0.1953) (0.0410,0.2051) (0.0430,0.2109) (0.0449,0.2168) (0.0469,0.2188) (0.0488,0.2285) (0.0508,0.2344) (0.0527,0.2402) (0.0547,0.2422) (0.0566,0.2480) (0.0586,0.2500) (0.0605,0.2520) (0.0625,0.2500) (0.0645,0.2637) (0.0664,0.2734) (0.0684,0.2832) (0.0703,0.2891) (0.0723,0.2988) (0.0742,0.3047) (0.0762,0.3105) (0.0781,0.3125) (0.0801,0.3223) (0.0820,0.3281) (0.0840,0.3340) (0.0859,0.3359) (0.0879,0.3418) (0.0898,0.3438) (0.0918,0.3457) (0.0938,0.3438) (0.0957,0.3535) (0.0977,0.3594) (0.0996,0.3652) (0.1016,0.3672) (0.1035,0.3730) (0.1055,0.3750) (0.1074,0.3770) (0.1094,0.3750) (0.1113,0.3809) (0.1133,0.3828) (0.1152,0.3848) (0.1172,0.3828) (0.1191,0.3848) (0.1211,0.3828) (0.1230,0.3809) (0.1250,0.3750) (0.1270,0.3887) (0.1289,0.3984) (0.1309,0.4082) (0.1328,0.4141) (0.1348,0.4238) (0.1367,0.4297) (0.1387,0.4355) (0.1406,0.4375) (0.1426,0.4473) (0.1445,0.4531) (0.1465,0.4590) (0.1484,0.4609) (0.1504,0.4668) (0.1523,0.4688) (0.1543,0.4707) (0.1562,0.4688) (0.1582,0.4785) (0.1602,0.4844) (0.1621,0.4902) (0.1641,0.4922) (0.1660,0.4980) (0.1680,0.5000) (0.1699,0.5020) (0.1719,0.5000) (0.1738,0.5059) (0.1758,0.5078) (0.1777,0.5098) (0.1797,0.5078) (0.1816,0.5098) (0.1836,0.5078) (0.1855,0.5059) (0.1875,0.5000) (0.1895,0.5098) (0.1914,0.5156) (0.1934,0.5215) (0.1953,0.5234) (0.1973,0.5293) (0.1992,0.5312) (0.2012,0.5332) (0.2031,0.5312) (0.2051,0.5371) (0.2070,0.5391) (0.2090,0.5410) (0.2109,0.5391) (0.2129,0.5410) (0.2148,0.5391) (0.2168,0.5371) (0.2188,0.5312) (0.2207,0.5371) (0.2227,0.5391) (0.2246,0.5410) (0.2266,0.5391) (0.2285,0.5410) (0.2305,0.5391) (0.2324,0.5371) (0.2344,0.5312) (0.2363,0.5332) (0.2383,0.5312) (0.2402,0.5293) (0.2422,0.5234) (0.2441,0.5215) (0.2461,0.5156) (0.2480,0.5098) (0.2500,0.5000) (0.2520,0.5137) (0.2539,0.5234) (0.2559,0.5332) (0.2578,0.5391) (0.2598,0.5488) (0.2617,0.5547) (0.2637,0.5605) (0.2656,0.5625) (0.2676,0.5723) (0.2695,0.5781) (0.2715,0.5840) (0.2734,0.5859) (0.2754,0.5918) (0.2773,0.5938) (0.2793,0.5957) (0.2812,0.5938) (0.2832,0.6035) (0.2852,0.6094) (0.2871,0.6152) (0.2891,0.6172) (0.2910,0.6230) (0.2930,0.6250) (0.2949,0.6270) (0.2969,0.6250) (0.2988,0.6309) (0.3008,0.6328) (0.3027,0.6348) (0.3047,0.6328) (0.3066,0.6348) (0.3086,0.6328) (0.3105,0.6309) (0.3125,0.6250) (0.3145,0.6348) (0.3164,0.6406) (0.3184,0.6465) (0.3203,0.6484) (0.3223,0.6543) (0.3242,0.6562) (0.3262,0.6582) (0.3281,0.6562) (0.3301,0.6621) (0.3320,0.6641) (0.3340,0.6660) (0.3359,0.6641) (0.3379,0.6660) (0.3398,0.6641) (0.3418,0.6621) (0.3438,0.6562) (0.3457,0.6621) (0.3477,0.6641) (0.3496,0.6660) (0.3516,0.6641) (0.3535,0.6660) (0.3555,0.6641) (0.3574,0.6621) (0.3594,0.6562) (0.3613,0.6582) (0.3633,0.6562) (0.3652,0.6543) (0.3672,0.6484) (0.3691,0.6465) (0.3711,0.6406) (0.3730,0.6348) (0.3750,0.6250) (0.3770,0.6348) (0.3789,0.6406) (0.3809,0.6465) (0.3828,0.6484) (0.3848,0.6543) (0.3867,0.6562) (0.3887,0.6582) (0.3906,0.6562) (0.3926,0.6621) (0.3945,0.6641) (0.3965,0.6660) (0.3984,0.6641) (0.4004,0.6660) (0.4023,0.6641) (0.4043,0.6621) (0.4062,0.6562) (0.4082,0.6621) (0.4102,0.6641) (0.4121,0.6660) (0.4141,0.6641) (0.4160,0.6660) (0.4180,0.6641) (0.4199,0.6621) (0.4219,0.6562) (0.4238,0.6582) (0.4258,0.6562) (0.4277,0.6543) (0.4297,0.6484) (0.4316,0.6465) (0.4336,0.6406) (0.4355,0.6348) (0.4375,0.6250) (0.4395,0.6309) (0.4414,0.6328) (0.4434,0.6348) (0.4453,0.6328) (0.4473,0.6348) (0.4492,0.6328) (0.4512,0.6309) (0.4531,0.6250) (0.4551,0.6270) (0.4570,0.6250) (0.4590,0.6230) (0.4609,0.6172) (0.4629,0.6152) (0.4648,0.6094) (0.4668,0.6035) (0.4688,0.5938) (0.4707,0.5957) (0.4727,0.5938) (0.4746,0.5918) (0.4766,0.5859) (0.4785,0.5840) (0.4805,0.5781) (0.4824,0.5723) (0.4844,0.5625) (0.4863,0.5605) (0.4883,0.5547) (0.4902,0.5488) (0.4922,0.5391) (0.4941,0.5332) (0.4961,0.5234) (0.4980,0.5137) (0.5000,0.5000) (0.5020,0.5137) (0.5039,0.5234) (0.5059,0.5332) (0.5078,0.5391) (0.5098,0.5488) (0.5117,0.5547) (0.5137,0.5605) (0.5156,0.5625) (0.5176,0.5723) (0.5195,0.5781) (0.5215,0.5840) (0.5234,0.5859) (0.5254,0.5918) (0.5273,0.5938) (0.5293,0.5957) (0.5312,0.5938) (0.5332,0.6035) (0.5352,0.6094) (0.5371,0.6152) (0.5391,0.6172) (0.5410,0.6230) (0.5430,0.6250) (0.5449,0.6270) (0.5469,0.6250) (0.5488,0.6309) (0.5508,0.6328) (0.5527,0.6348) (0.5547,0.6328) (0.5566,0.6348) (0.5586,0.6328) (0.5605,0.6309) (0.5625,0.6250) (0.5645,0.6348) (0.5664,0.6406) (0.5684,0.6465) (0.5703,0.6484) (0.5723,0.6543) (0.5742,0.6562) (0.5762,0.6582) (0.5781,0.6562) (0.5801,0.6621) (0.5820,0.6641) (0.5840,0.6660) (0.5859,0.6641) (0.5879,0.6660) (0.5898,0.6641) (0.5918,0.6621) (0.5938,0.6562) (0.5957,0.6621) (0.5977,0.6641) (0.5996,0.6660) (0.6016,0.6641) (0.6035,0.6660) (0.6055,0.6641) (0.6074,0.6621) (0.6094,0.6562) (0.6113,0.6582) (0.6133,0.6562) (0.6152,0.6543) (0.6172,0.6484) (0.6191,0.6465) (0.6211,0.6406) (0.6230,0.6348) (0.6250,0.6250) (0.6270,0.6348) (0.6289,0.6406) (0.6309,0.6465) (0.6328,0.6484) (0.6348,0.6543) (0.6367,0.6562) (0.6387,0.6582) (0.6406,0.6562) (0.6426,0.6621) (0.6445,0.6641) (0.6465,0.6660) (0.6484,0.6641) (0.6504,0.6660) (0.6523,0.6641) (0.6543,0.6621) (0.6562,0.6562) (0.6582,0.6621) (0.6602,0.6641) (0.6621,0.6660) (0.6641,0.6641) (0.6660,0.6660) (0.6680,0.6641) (0.6699,0.6621) (0.6719,0.6562) (0.6738,0.6582) (0.6758,0.6562) (0.6777,0.6543) (0.6797,0.6484) (0.6816,0.6465) (0.6836,0.6406) (0.6855,0.6348) (0.6875,0.6250) (0.6895,0.6309) (0.6914,0.6328) (0.6934,0.6348) (0.6953,0.6328) (0.6973,0.6348) (0.6992,0.6328) (0.7012,0.6309) (0.7031,0.6250) (0.7051,0.6270) (0.7070,0.6250) (0.7090,0.6230) (0.7109,0.6172) (0.7129,0.6152) (0.7148,0.6094) (0.7168,0.6035) (0.7188,0.5938) (0.7207,0.5957) (0.7227,0.5938) (0.7246,0.5918) (0.7266,0.5859) (0.7285,0.5840) (0.7305,0.5781) (0.7324,0.5723) (0.7344,0.5625) (0.7363,0.5605) (0.7383,0.5547) (0.7402,0.5488) (0.7422,0.5391) (0.7441,0.5332) (0.7461,0.5234) (0.7480,0.5137) (0.7500,0.5000) (0.7520,0.5098) (0.7539,0.5156) (0.7559,0.5215) (0.7578,0.5234) (0.7598,0.5293) (0.7617,0.5312) (0.7637,0.5332) (0.7656,0.5312) (0.7676,0.5371) (0.7695,0.5391) (0.7715,0.5410) (0.7734,0.5391) (0.7754,0.5410) (0.7773,0.5391) (0.7793,0.5371) (0.7812,0.5312) (0.7832,0.5371) (0.7852,0.5391) (0.7871,0.5410) (0.7891,0.5391) (0.7910,0.5410) (0.7930,0.5391) (0.7949,0.5371) (0.7969,0.5312) (0.7988,0.5332) (0.8008,0.5312) (0.8027,0.5293) (0.8047,0.5234) (0.8066,0.5215) (0.8086,0.5156) (0.8105,0.5098) (0.8125,0.5000) (0.8145,0.5059) (0.8164,0.5078) (0.8184,0.5098) (0.8203,0.5078) (0.8223,0.5098) (0.8242,0.5078) (0.8262,0.5059) (0.8281,0.5000) (0.8301,0.5020) (0.8320,0.5000) (0.8340,0.4980) (0.8359,0.4922) (0.8379,0.4902) (0.8398,0.4844) (0.8418,0.4785) (0.8438,0.4688) (0.8457,0.4707) (0.8477,0.4688) (0.8496,0.4668) (0.8516,0.4609) (0.8535,0.4590) (0.8555,0.4531) (0.8574,0.4473) (0.8594,0.4375) (0.8613,0.4355) (0.8633,0.4297) (0.8652,0.4238) (0.8672,0.4141) (0.8691,0.4082) (0.8711,0.3984) (0.8730,0.3887) (0.8750,0.3750) (0.8770,0.3809) (0.8789,0.3828) (0.8809,0.3848) (0.8828,0.3828) (0.8848,0.3848) (0.8867,0.3828) (0.8887,0.3809) (0.8906,0.3750) (0.8926,0.3770) (0.8945,0.3750) (0.8965,0.3730) (0.8984,0.3672) (0.9004,0.3652) (0.9023,0.3594) (0.9043,0.3535) (0.9062,0.3438) (0.9082,0.3457) (0.9102,0.3438) (0.9121,0.3418) (0.9141,0.3359) (0.9160,0.3340) (0.9180,0.3281) (0.9199,0.3223) (0.9219,0.3125) (0.9238,0.3105) (0.9258,0.3047) (0.9277,0.2988) (0.9297,0.2891) (0.9316,0.2832) (0.9336,0.2734) (0.9355,0.2637) (0.9375,0.2500) (0.9395,0.2520) (0.9414,0.2500) (0.9434,0.2480) (0.9453,0.2422) (0.9473,0.2402) (0.9492,0.2344) (0.9512,0.2285) (0.9531,0.2188) (0.9551,0.2168) (0.9570,0.2109) (0.9590,0.2051) (0.9609,0.1953) (0.9629,0.1895) (0.9648,0.1797) (0.9668,0.1699) (0.9688,0.1562) (0.9707,0.1543) (0.9727,0.1484) (0.9746,0.1426) (0.9766,0.1328) (0.9785,0.1270) (0.9805,0.1172) (0.9824,0.1074) (0.9844,0.0938) (0.9863,0.0879) (0.9883,0.0781) (0.9902,0.0684) (0.9922,0.0547) (0.9941,0.0449) (0.9961,0.0312) (0.9980,0.0176) (1.0000,0.0000)};
\addlegendentry{$T$}
\end{axis}
\end{tikzpicture}}\hfill
\sdpanel{B}{\begin{tikzpicture}
\begin{axis}[sdplot, width=0.47\textwidth, height=0.36\textwidth, xmin=0, xmax=2.04, ymin=-2.1, ymax=2.1, xtick={0,0.5,1,1.5,2}, ytick={-2,-1,0,1,2}, xlabel={$x$}]
\addplot[black, line width=0.35pt] coordinates {(0.0000,1.9922) (0.0010,1.9555) (0.0020,1.9317) (0.0029,1.9132) (0.0039,1.8546) (0.0049,1.8543) (0.0059,1.8569) (0.0068,1.8300) (0.0078,1.8466) (0.0088,1.8204) (0.0098,1.7672) (0.0107,1.7573) (0.0117,1.7124) (0.0127,1.6915) (0.0137,1.7239) (0.0146,1.7165) (0.0156,1.7343) (0.0166,1.7568) (0.0176,1.7154) (0.0186,1.7058) (0.0195,1.6925) (0.0205,1.6535) (0.0215,1.6803) (0.0225,1.6940) (0.0234,1.6885) (0.0244,1.7196) (0.0254,1.6907) (0.0264,1.6455) (0.0273,1.6304) (0.0283,1.5694) (0.0293,1.5429) (0.0303,1.5514) (0.0312,1.5166) (0.0322,1.5163) (0.0332,1.5047) (0.0342,1.4398) (0.0352,1.4188) (0.0361,1.3855) (0.0371,1.3465) (0.0381,1.3798) (0.0391,1.3925) (0.0400,1.4063) (0.0410,1.4542) (0.0420,1.4395) (0.0430,1.4295) (0.0439,1.4425) (0.0449,1.4127) (0.0459,1.4324) (0.0469,1.4707) (0.0479,1.4686) (0.0488,1.5049) (0.0498,1.5087) (0.0508,1.4659) (0.0518,1.4636) (0.0527,1.4299) (0.0537,1.3994) (0.0547,1.4285) (0.0557,1.4202) (0.0566,1.4221) (0.0576,1.4423) (0.0586,1.3957) (0.0596,1.3697) (0.0605,1.3561) (0.0615,1.3098) (0.0625,1.3286) (0.0635,1.3554) (0.0645,1.3573) (0.0654,1.4054) (0.0664,1.4112) (0.0674,1.3897) (0.0684,1.4090) (0.0693,1.3890) (0.0703,1.3883) (0.0713,1.4343) (0.0723,1.4337) (0.0732,1.4517) (0.0742,1.4672) (0.0752,1.4130) (0.0762,1.3857) (0.0771,1.3505) (0.0781,1.2877) (0.0791,1.2900) (0.0801,1.2795) (0.0811,1.2526) (0.0820,1.2657) (0.0830,1.2233) (0.0840,1.1703) (0.0850,1.1531) (0.0859,1.0956) (0.0869,1.0784) (0.0879,1.1003) (0.0889,1.0823) (0.0898,1.1012) (0.0908,1.1088) (0.0918,1.0622) (0.0928,1.0570) (0.0938,1.0350) (0.0947,1.0025) (0.0957,1.0358) (0.0967,1.0414) (0.0977,1.0415) (0.0986,1.0688) (0.0996,1.0273) (0.1006,0.9860) (0.1016,0.9637) (0.1025,0.8966) (0.1035,0.8791) (0.1045,0.8810) (0.1055,0.8460) (0.1064,0.8539) (0.1074,0.8347) (0.1084,0.7762) (0.1094,0.7654) (0.1104,0.7303) (0.1113,0.7060) (0.1123,0.7475) (0.1133,0.7570) (0.1143,0.7812) (0.1152,0.8260) (0.1162,0.8051) (0.1172,0.8049) (0.1182,0.8146) (0.1191,0.7887) (0.1201,0.8238) (0.1211,0.8614) (0.1221,0.8691) (0.1230,0.9176) (0.1240,0.9186) (0.1250,0.8887) (0.1260,0.8966) (0.1270,0.8635) (0.1279,0.8501) (0.1289,0.8848) (0.1299,0.8760) (0.1309,0.8907) (0.1318,0.9082) (0.1328,0.8629) (0.1338,0.8520) (0.1348,0.8401) (0.1357,0.8079) (0.1367,0.8472) (0.1377,0.8785) (0.1387,0.8974) (0.1396,0.9585) (0.1406,0.9641) (0.1416,0.9585) (0.1426,0.9854) (0.1436,0.9675) (0.1445,0.9848) (0.1455,1.0340) (0.1465,1.0365) (0.1475,1.0688) (0.1484,1.0824) (0.1494,1.0359) (0.1504,1.0257) (0.1514,0.9944) (0.1523,0.9505) (0.1533,0.9715) (0.1543,0.9651) (0.1553,0.9556) (0.1562,0.9765) (0.1572,0.9328) (0.1582,0.8950) (0.1592,0.8813) (0.1602,0.8282) (0.1611,0.8298) (0.1621,0.8544) (0.1631,0.8449) (0.1641,0.8800) (0.1650,0.8873) (0.1660,0.8538) (0.1670,0.8648) (0.1680,0.8463) (0.1689,0.8331) (0.1699,0.8787) (0.1709,0.8846) (0.1719,0.8983) (0.1729,0.9251) (0.1738,0.8799) (0.1748,0.8505) (0.1758,0.8268) (0.1768,0.7655) (0.1777,0.7649) (0.1787,0.7675) (0.1797,0.7433) (0.1807,0.7642) (0.1816,0.7426) (0.1826,0.6970) (0.1836,0.6959) (0.1846,0.6605) (0.1855,0.6519) (0.1865,0.6970) (0.1875,0.7033) (0.1885,0.7371) (0.1895,0.7752) (0.1904,0.7512) (0.1914,0.7608) (0.1924,0.7665) (0.1934,0.7483) (0.1943,0.7968) (0.1953,0.8313) (0.1963,0.8480) (0.1973,0.9011) (0.1982,0.8932) (0.1992,0.8702) (0.2002,0.8765) (0.2012,0.8363) (0.2021,0.8315) (0.2031,0.8598) (0.2041,0.8445) (0.2051,0.8638) (0.2061,0.8695) (0.2070,0.8220) (0.2080,0.8177) (0.2090,0.7990) (0.2100,0.7747) (0.2109,0.8210) (0.2119,0.8442) (0.2129,0.8681) (0.2139,0.9237) (0.2148,0.9146) (0.2158,0.9101) (0.2168,0.9259) (0.2178,0.8977) (0.2188,0.9186) (0.2197,0.9551) (0.2207,0.9506) (0.2217,0.9832) (0.2227,0.9808) (0.2236,0.9316) (0.2246,0.9216) (0.2256,0.8781) (0.2266,0.8380) (0.2275,0.8555) (0.2285,0.8341) (0.2295,0.8230) (0.2305,0.8283) (0.2314,0.7662) (0.2324,0.7254) (0.2334,0.6954) (0.2344,0.6334) (0.2354,0.6370) (0.2363,0.6472) (0.2373,0.6336) (0.2383,0.6662) (0.2393,0.6556) (0.2402,0.6194) (0.2412,0.6240) (0.2422,0.5893) (0.2432,0.5760) (0.2441,0.6090) (0.2451,0.5961) (0.2461,0.6038) (0.2471,0.6086) (0.2480,0.5455) (0.2490,0.5114) (0.2500,0.4696) (0.2510,0.4025) (0.2520,0.4021) (0.2529,0.3889) (0.2539,0.3617) (0.2549,0.3752) (0.2559,0.3335) (0.2568,0.2839) (0.2578,0.2705) (0.2588,0.2178) (0.2598,0.2077) (0.2607,0.2366) (0.2617,0.2268) (0.2627,0.2557) (0.2637,0.2727) (0.2646,0.2374) (0.2656,0.2445) (0.2666,0.2344) (0.2676,0.2155) (0.2686,0.2624) (0.2695,0.2807) (0.2705,0.2948) (0.2715,0.3349) (0.2725,0.3056) (0.2734,0.2776) (0.2744,0.2673) (0.2754,0.2120) (0.2764,0.2067) (0.2773,0.2189) (0.2783,0.1940) (0.2793,0.2115) (0.2803,0.1997) (0.2812,0.1486) (0.2822,0.1441) (0.2832,0.1130) (0.2842,0.0929) (0.2852,0.1362) (0.2861,0.1453) (0.2871,0.1688) (0.2881,0.2101) (0.2891,0.1843) (0.2900,0.1787) (0.2910,0.1804) (0.2920,0.1458) (0.2930,0.1714) (0.2939,0.1967) (0.2949,0.1915) (0.2959,0.2257) (0.2969,0.2099) (0.2979,0.1632) (0.2988,0.1526) (0.2998,0.0992) (0.3008,0.0659) (0.3018,0.0787) (0.3027,0.0469) (0.3037,0.0389) (0.3047,0.0319) (0.3057,-0.0379) (0.3066,-0.0726) (0.3076,-0.1095) (0.3086,-0.1657) (0.3096,-0.1502) (0.3105,-0.1438) (0.3115,-0.1483) (0.3125,-0.1109) (0.3135,-0.1294) (0.3145,-0.1570) (0.3154,-0.1522) (0.3164,-0.1915) (0.3174,-0.1934) (0.3184,-0.1634) (0.3193,-0.1787) (0.3203,-0.1621) (0.3213,-0.1642) (0.3223,-0.2240) (0.3232,-0.2454) (0.3242,-0.2874) (0.3252,-0.3391) (0.3262,-0.3247) (0.3271,-0.3371) (0.3281,-0.3497) (0.3291,-0.3314) (0.3301,-0.3764) (0.3311,-0.4127) (0.3320,-0.4243) (0.3330,-0.4733) (0.3340,-0.4650) (0.3350,-0.4336) (0.3359,-0.4342) (0.3369,-0.3885) (0.3379,-0.3707) (0.3389,-0.3913) (0.3398,-0.3665) (0.3408,-0.3711) (0.3418,-0.3682) (0.3428,-0.3067) (0.3438,-0.2850) (0.3447,-0.2540) (0.3457,-0.2107) (0.3467,-0.2391) (0.3477,-0.2503) (0.3486,-0.2569) (0.3496,-0.3004) (0.3506,-0.2827) (0.3516,-0.2634) (0.3525,-0.2704) (0.3535,-0.2329) (0.3545,-0.2397) (0.3555,-0.2699) (0.3564,-0.2570) (0.3574,-0.2799) (0.3584,-0.2756) (0.3594,-0.2198) (0.3604,-0.2041) (0.3613,-0.1611) (0.3623,-0.1162) (0.3633,-0.1340) (0.3643,-0.1185) (0.3652,-0.1093) (0.3662,-0.1240) (0.3672,-0.0730) (0.3682,-0.0384) (0.3691,-0.0216) (0.3701,0.0298) (0.3711,0.0182) (0.3721,-0.0079) (0.3730,-0.0067) (0.3740,-0.0529) (0.3750,-0.0632) (0.3760,-0.0422) (0.3770,-0.0651) (0.3779,-0.0530) (0.3789,-0.0562) (0.3799,-0.1119) (0.3809,-0.1238) (0.3818,-0.1513) (0.3828,-0.1829) (0.3838,-0.1440) (0.3848,-0.1290) (0.3857,-0.1116) (0.3867,-0.0630) (0.3877,-0.0789) (0.3887,-0.0881) (0.3896,-0.0770) (0.3906,-0.1086) (0.3916,-0.0889) (0.3926,-0.0537) (0.3936,-0.0575) (0.3945,-0.0224) (0.3955,-0.0223) (0.3965,-0.0662) (0.3975,-0.0693) (0.3984,-0.1054) (0.3994,-0.1351) (0.4004,-0.1064) (0.4014,-0.1158) (0.4023,-0.1120) (0.4033,-0.0915) (0.4043,-0.1368) (0.4053,-0.1579) (0.4062,-0.1679) (0.4072,-0.2075) (0.4082,-0.1794) (0.4092,-0.1448) (0.4102,-0.1318) (0.4111,-0.0714) (0.4121,-0.0545) (0.4131,-0.0608) (0.4141,-0.0260) (0.4150,-0.0305) (0.4160,-0.0117) (0.4170,0.0529) (0.4180,0.0717) (0.4189,0.1122) (0.4199,0.1487) (0.4209,0.1179) (0.4219,0.1169) (0.4229,0.1068) (0.4238,0.0722) (0.4248,0.1039) (0.4258,0.1210) (0.4268,0.1243) (0.4277,0.1670) (0.4287,0.1524) (0.4297,0.1303) (0.4307,0.1422) (0.4316,0.1134) (0.4326,0.1273) (0.4336,0.1774) (0.4346,0.1877) (0.4355,0.2359) (0.4365,0.2693) (0.4375,0.2498) (0.4385,0.2717) (0.4395,0.2735) (0.4404,0.2663) (0.4414,0.3229) (0.4424,0.3483) (0.4434,0.3689) (0.4443,0.4130) (0.4453,0.3858) (0.4463,0.3599) (0.4473,0.3491) (0.4482,0.2927) (0.4492,0.2861) (0.4502,0.2946) (0.4512,0.2657) (0.4521,0.2788) (0.4531,0.2605) (0.4541,0.2035) (0.4551,0.1924) (0.4561,0.1535) (0.4570,0.1265) (0.4580,0.1619) (0.4590,0.1623) (0.4600,0.1782) (0.4609,0.2108) (0.4619,0.1766) (0.4629,0.1642) (0.4639,0.1586) (0.4648,0.1180) (0.4658,0.1390) (0.4668,0.1591) (0.4678,0.1507) (0.4688,0.1823) (0.4697,0.1637) (0.4707,0.1167) (0.4717,0.1063) (0.4727,0.0537) (0.4736,0.0237) (0.4746,0.0398) (0.4756,0.0124) (0.4766,0.0110) (0.4775,0.0104) (0.4785,-0.0510) (0.4795,-0.0753) (0.4805,-0.1019) (0.4814,-0.1454) (0.4824,-0.1160) (0.4834,-0.0961) (0.4844,-0.0851) (0.4854,-0.0320) (0.4863,-0.0351) (0.4873,-0.0453) (0.4883,-0.0236) (0.4893,-0.0455) (0.4902,-0.0285) (0.4912,0.0191) (0.4922,0.0219) (0.4932,0.0574) (0.4941,0.0726) (0.4951,0.0309) (0.4961,0.0276) (0.4971,0.0021) (0.4980,-0.0323) (0.4990,-0.0017) (0.5000,0.0000) (0.5010,0.0017) (0.5020,0.0323) (0.5029,-0.0021) (0.5039,-0.0276) (0.5049,-0.0309) (0.5059,-0.0726) (0.5068,-0.0574) (0.5078,-0.0219) (0.5088,-0.0191) (0.5098,0.0285) (0.5107,0.0455) (0.5117,0.0236) (0.5127,0.0453) (0.5137,0.0351) (0.5146,0.0320) (0.5156,0.0851) (0.5166,0.0961) (0.5176,0.1160) (0.5186,0.1454) (0.5195,0.1019) (0.5205,0.0753) (0.5215,0.0510) (0.5225,-0.0104) (0.5234,-0.0110) (0.5244,-0.0124) (0.5254,-0.0398) (0.5264,-0.0237) (0.5273,-0.0537) (0.5283,-0.1063) (0.5293,-0.1167) (0.5303,-0.1637) (0.5312,-0.1823) (0.5322,-0.1507) (0.5332,-0.1591) (0.5342,-0.1390) (0.5352,-0.1180) (0.5361,-0.1586) (0.5371,-0.1642) (0.5381,-0.1766) (0.5391,-0.2108) (0.5400,-0.1782) (0.5410,-0.1623) (0.5420,-0.1619) (0.5430,-0.1265) (0.5439,-0.1535) (0.5449,-0.1924) (0.5459,-0.2035) (0.5469,-0.2605) (0.5479,-0.2788) (0.5488,-0.2657) (0.5498,-0.2946) (0.5508,-0.2861) (0.5518,-0.2927) (0.5527,-0.3491) (0.5537,-0.3599) (0.5547,-0.3858) (0.5557,-0.4130) (0.5566,-0.3689) (0.5576,-0.3483) (0.5586,-0.3229) (0.5596,-0.2663) (0.5605,-0.2735) (0.5615,-0.2717) (0.5625,-0.2498) (0.5635,-0.2693) (0.5645,-0.2359) (0.5654,-0.1877) (0.5664,-0.1774) (0.5674,-0.1273) (0.5684,-0.1134) (0.5693,-0.1422) (0.5703,-0.1303) (0.5713,-0.1524) (0.5723,-0.1670) (0.5732,-0.1243) (0.5742,-0.1210) (0.5752,-0.1039) (0.5762,-0.0722) (0.5771,-0.1068) (0.5781,-0.1169) (0.5791,-0.1179) (0.5801,-0.1487) (0.5811,-0.1122) (0.5820,-0.0717) (0.5830,-0.0529) (0.5840,0.0117) (0.5850,0.0305) (0.5859,0.0260) (0.5869,0.0608) (0.5879,0.0545) (0.5889,0.0714) (0.5898,0.1318) (0.5908,0.1448) (0.5918,0.1794) (0.5928,0.2075) (0.5938,0.1679) (0.5947,0.1579) (0.5957,0.1368) (0.5967,0.0915) (0.5977,0.1120) (0.5986,0.1158) (0.5996,0.1064) (0.6006,0.1351) (0.6016,0.1054) (0.6025,0.0693) (0.6035,0.0662) (0.6045,0.0223) (0.6055,0.0224) (0.6064,0.0575) (0.6074,0.0537) (0.6084,0.0889) (0.6094,0.1086) (0.6104,0.0770) (0.6113,0.0881) (0.6123,0.0789) (0.6133,0.0630) (0.6143,0.1116) (0.6152,0.1290) (0.6162,0.1440) (0.6172,0.1829) (0.6182,0.1513) (0.6191,0.1238) (0.6201,0.1119) (0.6211,0.0562) (0.6221,0.0530) (0.6230,0.0651) (0.6240,0.0422) (0.6250,0.0632) (0.6260,0.0529) (0.6270,0.0067) (0.6279,0.0079) (0.6289,-0.0182) (0.6299,-0.0298) (0.6309,0.0216) (0.6318,0.0384) (0.6328,0.0730) (0.6338,0.1240) (0.6348,0.1093) (0.6357,0.1185) (0.6367,0.1340) (0.6377,0.1162) (0.6387,0.1611) (0.6396,0.2041) (0.6406,0.2198) (0.6416,0.2756) (0.6426,0.2799) (0.6436,0.2570) (0.6445,0.2699) (0.6455,0.2397) (0.6465,0.2329) (0.6475,0.2704) (0.6484,0.2634) (0.6494,0.2827) (0.6504,0.3004) (0.6514,0.2569) (0.6523,0.2503) (0.6533,0.2391) (0.6543,0.2107) (0.6553,0.2540) (0.6562,0.2850) (0.6572,0.3067) (0.6582,0.3682) (0.6592,0.3711) (0.6602,0.3665) (0.6611,0.3913) (0.6621,0.3707) (0.6631,0.3885) (0.6641,0.4342) (0.6650,0.4336) (0.6660,0.4650) (0.6670,0.4733) (0.6680,0.4243) (0.6689,0.4127) (0.6699,0.3764) (0.6709,0.3314) (0.6719,0.3497) (0.6729,0.3371) (0.6738,0.3247) (0.6748,0.3391) (0.6758,0.2874) (0.6768,0.2454) (0.6777,0.2240) (0.6787,0.1642) (0.6797,0.1621) (0.6807,0.1787) (0.6816,0.1634) (0.6826,0.1934) (0.6836,0.1915) (0.6846,0.1522) (0.6855,0.1570) (0.6865,0.1294) (0.6875,0.1109) (0.6885,0.1483) (0.6895,0.1438) (0.6904,0.1502) (0.6914,0.1657) (0.6924,0.1095) (0.6934,0.0726) (0.6943,0.0379) (0.6953,-0.0319) (0.6963,-0.0389) (0.6973,-0.0469) (0.6982,-0.0787) (0.6992,-0.0659) (0.7002,-0.0992) (0.7012,-0.1526) (0.7021,-0.1632) (0.7031,-0.2099) (0.7041,-0.2257) (0.7051,-0.1915) (0.7061,-0.1967) (0.7070,-0.1714) (0.7080,-0.1458) (0.7090,-0.1804) (0.7100,-0.1787) (0.7109,-0.1843) (0.7119,-0.2101) (0.7129,-0.1688) (0.7139,-0.1453) (0.7148,-0.1362) (0.7158,-0.0929) (0.7168,-0.1130) (0.7178,-0.1441) (0.7188,-0.1486) (0.7197,-0.1997) (0.7207,-0.2115) (0.7217,-0.1940) (0.7227,-0.2189) (0.7236,-0.2067) (0.7246,-0.2120) (0.7256,-0.2673) (0.7266,-0.2776) (0.7275,-0.3056) (0.7285,-0.3349) (0.7295,-0.2948) (0.7305,-0.2807) (0.7314,-0.2624) (0.7324,-0.2155) (0.7334,-0.2344) (0.7344,-0.2445) (0.7354,-0.2374) (0.7363,-0.2727) (0.7373,-0.2557) (0.7383,-0.2268) (0.7393,-0.2366) (0.7402,-0.2077) (0.7412,-0.2178) (0.7422,-0.2705) (0.7432,-0.2839) (0.7441,-0.3335) (0.7451,-0.3752) (0.7461,-0.3617) (0.7471,-0.3889) (0.7480,-0.4021) (0.7490,-0.4025) (0.7500,-0.4696) (0.7510,-0.5114) (0.7520,-0.5455) (0.7529,-0.6086) (0.7539,-0.6038) (0.7549,-0.5961) (0.7559,-0.6090) (0.7568,-0.5760) (0.7578,-0.5893) (0.7588,-0.6240) (0.7598,-0.6194) (0.7607,-0.6556) (0.7617,-0.6662) (0.7627,-0.6336) (0.7637,-0.6472) (0.7646,-0.6370) (0.7656,-0.6334) (0.7666,-0.6954) (0.7676,-0.7254) (0.7686,-0.7662) (0.7695,-0.8283) (0.7705,-0.8230) (0.7715,-0.8341) (0.7725,-0.8555) (0.7734,-0.8380) (0.7744,-0.8781) (0.7754,-0.9216) (0.7764,-0.9316) (0.7773,-0.9808) (0.7783,-0.9832) (0.7793,-0.9506) (0.7803,-0.9551) (0.7812,-0.9186) (0.7822,-0.8977) (0.7832,-0.9259) (0.7842,-0.9101) (0.7852,-0.9146) (0.7861,-0.9237) (0.7871,-0.8681) (0.7881,-0.8442) (0.7891,-0.8210) (0.7900,-0.7747) (0.7910,-0.7990) (0.7920,-0.8177) (0.7930,-0.8220) (0.7939,-0.8695) (0.7949,-0.8638) (0.7959,-0.8445) (0.7969,-0.8598) (0.7979,-0.8315) (0.7988,-0.8363) (0.7998,-0.8765) (0.8008,-0.8702) (0.8018,-0.8932) (0.8027,-0.9011) (0.8037,-0.8480) (0.8047,-0.8313) (0.8057,-0.7968) (0.8066,-0.7483) (0.8076,-0.7665) (0.8086,-0.7608) (0.8096,-0.7512) (0.8105,-0.7752) (0.8115,-0.7371) (0.8125,-0.7033) (0.8135,-0.6970) (0.8145,-0.6519) (0.8154,-0.6605) (0.8164,-0.6959) (0.8174,-0.6970) (0.8184,-0.7426) (0.8193,-0.7642) (0.8203,-0.7433) (0.8213,-0.7675) (0.8223,-0.7649) (0.8232,-0.7655) (0.8242,-0.8268) (0.8252,-0.8505) (0.8262,-0.8799) (0.8271,-0.9251) (0.8281,-0.8983) (0.8291,-0.8846) (0.8301,-0.8787) (0.8311,-0.8331) (0.8320,-0.8463) (0.8330,-0.8648) (0.8340,-0.8538) (0.8350,-0.8873) (0.8359,-0.8800) (0.8369,-0.8449) (0.8379,-0.8544) (0.8389,-0.8298) (0.8398,-0.8282) (0.8408,-0.8813) (0.8418,-0.8950) (0.8428,-0.9328) (0.8438,-0.9765) (0.8447,-0.9556) (0.8457,-0.9651) (0.8467,-0.9715) (0.8477,-0.9505) (0.8486,-0.9944) (0.8496,-1.0257) (0.8506,-1.0359) (0.8516,-1.0824) (0.8525,-1.0688) (0.8535,-1.0365) (0.8545,-1.0340) (0.8555,-0.9848) (0.8564,-0.9675) (0.8574,-0.9854) (0.8584,-0.9585) (0.8594,-0.9641) (0.8604,-0.9585) (0.8613,-0.8974) (0.8623,-0.8785) (0.8633,-0.8472) (0.8643,-0.8079) (0.8652,-0.8401) (0.8662,-0.8520) (0.8672,-0.8629) (0.8682,-0.9082) (0.8691,-0.8907) (0.8701,-0.8760) (0.8711,-0.8848) (0.8721,-0.8501) (0.8730,-0.8635) (0.8740,-0.8966) (0.8750,-0.8887) (0.8760,-0.9186) (0.8770,-0.9176) (0.8779,-0.8691) (0.8789,-0.8614) (0.8799,-0.8238) (0.8809,-0.7887) (0.8818,-0.8146) (0.8828,-0.8049) (0.8838,-0.8051) (0.8848,-0.8260) (0.8857,-0.7812) (0.8867,-0.7570) (0.8877,-0.7475) (0.8887,-0.7060) (0.8896,-0.7303) (0.8906,-0.7654) (0.8916,-0.7762) (0.8926,-0.8347) (0.8936,-0.8539) (0.8945,-0.8460) (0.8955,-0.8810) (0.8965,-0.8791) (0.8975,-0.8966) (0.8984,-0.9637) (0.8994,-0.9860) (0.9004,-1.0273) (0.9014,-1.0688) (0.9023,-1.0415) (0.9033,-1.0414) (0.9043,-1.0358) (0.9053,-1.0025) (0.9062,-1.0350) (0.9072,-1.0570) (0.9082,-1.0622) (0.9092,-1.1088) (0.9102,-1.1012) (0.9111,-1.0823) (0.9121,-1.1003) (0.9131,-1.0784) (0.9141,-1.0956) (0.9150,-1.1531) (0.9160,-1.1703) (0.9170,-1.2233) (0.9180,-1.2657) (0.9189,-1.2526) (0.9199,-1.2795) (0.9209,-1.2900) (0.9219,-1.2877) (0.9229,-1.3505) (0.9238,-1.3857) (0.9248,-1.4130) (0.9258,-1.4672) (0.9268,-1.4517) (0.9277,-1.4337) (0.9287,-1.4343) (0.9297,-1.3883) (0.9307,-1.3890) (0.9316,-1.4090) (0.9326,-1.3897) (0.9336,-1.4112) (0.9346,-1.4054) (0.9355,-1.3573) (0.9365,-1.3554) (0.9375,-1.3286) (0.9385,-1.3098) (0.9395,-1.3561) (0.9404,-1.3697) (0.9414,-1.3957) (0.9424,-1.4423) (0.9434,-1.4221) (0.9443,-1.4202) (0.9453,-1.4285) (0.9463,-1.3994) (0.9473,-1.4299) (0.9482,-1.4636) (0.9492,-1.4659) (0.9502,-1.5087) (0.9512,-1.5049) (0.9521,-1.4686) (0.9531,-1.4707) (0.9541,-1.4324) (0.9551,-1.4127) (0.9561,-1.4425) (0.9570,-1.4295) (0.9580,-1.4395) (0.9590,-1.4542) (0.9600,-1.4063) (0.9609,-1.3925) (0.9619,-1.3798) (0.9629,-1.3465) (0.9639,-1.3855) (0.9648,-1.4188) (0.9658,-1.4398) (0.9668,-1.5047) (0.9678,-1.5163) (0.9688,-1.5166) (0.9697,-1.5514) (0.9707,-1.5429) (0.9717,-1.5694) (0.9727,-1.6304) (0.9736,-1.6455) (0.9746,-1.6907) (0.9756,-1.7196) (0.9766,-1.6885) (0.9775,-1.6940) (0.9785,-1.6803) (0.9795,-1.6535) (0.9805,-1.6925) (0.9814,-1.7058) (0.9824,-1.7154) (0.9834,-1.7568) (0.9844,-1.7343) (0.9854,-1.7165) (0.9863,-1.7239) (0.9873,-1.6915) (0.9883,-1.7124) (0.9893,-1.7573) (0.9902,-1.7672) (0.9912,-1.8204) (0.9922,-1.8466) (0.9932,-1.8300) (0.9941,-1.8569) (0.9951,-1.8543) (0.9961,-1.8546) (0.9971,-1.9132) (0.9980,-1.9317) (0.9990,-1.9555) (1.0000,-1.9922) (1.0010,-1.9555) (1.0020,-1.9317) (1.0029,-1.9132) (1.0039,-1.8546) (1.0049,-1.8543) (1.0059,-1.8569) (1.0068,-1.8300) (1.0078,-1.8466) (1.0088,-1.8204) (1.0098,-1.7672) (1.0107,-1.7573) (1.0117,-1.7124) (1.0127,-1.6915) (1.0137,-1.7239) (1.0146,-1.7165) (1.0156,-1.7343) (1.0166,-1.7568) (1.0176,-1.7154) (1.0186,-1.7058) (1.0195,-1.6925) (1.0205,-1.6535) (1.0215,-1.6803) (1.0225,-1.6940) (1.0234,-1.6885) (1.0244,-1.7196) (1.0254,-1.6907) (1.0264,-1.6455) (1.0273,-1.6304) (1.0283,-1.5694) (1.0293,-1.5429) (1.0303,-1.5514) (1.0312,-1.5166) (1.0322,-1.5163) (1.0332,-1.5047) (1.0342,-1.4398) (1.0352,-1.4188) (1.0361,-1.3855) (1.0371,-1.3465) (1.0381,-1.3798) (1.0391,-1.3925) (1.0400,-1.4063) (1.0410,-1.4542) (1.0420,-1.4395) (1.0430,-1.4295) (1.0439,-1.4425) (1.0449,-1.4127) (1.0459,-1.4324) (1.0469,-1.4707) (1.0479,-1.4686) (1.0488,-1.5049) (1.0498,-1.5087) (1.0508,-1.4659) (1.0518,-1.4636) (1.0527,-1.4299) (1.0537,-1.3994) (1.0547,-1.4285) (1.0557,-1.4202) (1.0566,-1.4221) (1.0576,-1.4423) (1.0586,-1.3957) (1.0596,-1.3697) (1.0605,-1.3561) (1.0615,-1.3098) (1.0625,-1.3286) (1.0635,-1.3554) (1.0645,-1.3573) (1.0654,-1.4054) (1.0664,-1.4112) (1.0674,-1.3897) (1.0684,-1.4090) (1.0693,-1.3890) (1.0703,-1.3883) (1.0713,-1.4343) (1.0723,-1.4337) (1.0732,-1.4517) (1.0742,-1.4672) (1.0752,-1.4130) (1.0762,-1.3857) (1.0771,-1.3505) (1.0781,-1.2877) (1.0791,-1.2900) (1.0801,-1.2795) (1.0811,-1.2526) (1.0820,-1.2657) (1.0830,-1.2233) (1.0840,-1.1703) (1.0850,-1.1531) (1.0859,-1.0956) (1.0869,-1.0784) (1.0879,-1.1003) (1.0889,-1.0823) (1.0898,-1.1012) (1.0908,-1.1088) (1.0918,-1.0622) (1.0928,-1.0570) (1.0938,-1.0350) (1.0947,-1.0025) (1.0957,-1.0358) (1.0967,-1.0414) (1.0977,-1.0415) (1.0986,-1.0688) (1.0996,-1.0273) (1.1006,-0.9860) (1.1016,-0.9637) (1.1025,-0.8966) (1.1035,-0.8791) (1.1045,-0.8810) (1.1055,-0.8460) (1.1064,-0.8539) (1.1074,-0.8347) (1.1084,-0.7762) (1.1094,-0.7654) (1.1104,-0.7303) (1.1113,-0.7060) (1.1123,-0.7475) (1.1133,-0.7570) (1.1143,-0.7812) (1.1152,-0.8260) (1.1162,-0.8051) (1.1172,-0.8049) (1.1182,-0.8146) (1.1191,-0.7887) (1.1201,-0.8238) (1.1211,-0.8614) (1.1221,-0.8691) (1.1230,-0.9176) (1.1240,-0.9186) (1.1250,-0.8887) (1.1260,-0.8966) (1.1270,-0.8635) (1.1279,-0.8501) (1.1289,-0.8848) (1.1299,-0.8760) (1.1309,-0.8907) (1.1318,-0.9082) (1.1328,-0.8629) (1.1338,-0.8520) (1.1348,-0.8401) (1.1357,-0.8079) (1.1367,-0.8472) (1.1377,-0.8785) (1.1387,-0.8974) (1.1396,-0.9585) (1.1406,-0.9641) (1.1416,-0.9585) (1.1426,-0.9854) (1.1436,-0.9675) (1.1445,-0.9848) (1.1455,-1.0340) (1.1465,-1.0365) (1.1475,-1.0688) (1.1484,-1.0824) (1.1494,-1.0359) (1.1504,-1.0257) (1.1514,-0.9944) (1.1523,-0.9505) (1.1533,-0.9715) (1.1543,-0.9651) (1.1553,-0.9556) (1.1562,-0.9765) (1.1572,-0.9328) (1.1582,-0.8950) (1.1592,-0.8813) (1.1602,-0.8282) (1.1611,-0.8298) (1.1621,-0.8544) (1.1631,-0.8449) (1.1641,-0.8800) (1.1650,-0.8873) (1.1660,-0.8538) (1.1670,-0.8648) (1.1680,-0.8463) (1.1689,-0.8331) (1.1699,-0.8787) (1.1709,-0.8846) (1.1719,-0.8983) (1.1729,-0.9251) (1.1738,-0.8799) (1.1748,-0.8505) (1.1758,-0.8268) (1.1768,-0.7655) (1.1777,-0.7649) (1.1787,-0.7675) (1.1797,-0.7433) (1.1807,-0.7642) (1.1816,-0.7426) (1.1826,-0.6970) (1.1836,-0.6959) (1.1846,-0.6605) (1.1855,-0.6519) (1.1865,-0.6970) (1.1875,-0.7033) (1.1885,-0.7371) (1.1895,-0.7752) (1.1904,-0.7512) (1.1914,-0.7608) (1.1924,-0.7665) (1.1934,-0.7483) (1.1943,-0.7968) (1.1953,-0.8313) (1.1963,-0.8480) (1.1973,-0.9011) (1.1982,-0.8932) (1.1992,-0.8702) (1.2002,-0.8765) (1.2012,-0.8363) (1.2021,-0.8315) (1.2031,-0.8598) (1.2041,-0.8445) (1.2051,-0.8638) (1.2061,-0.8695) (1.2070,-0.8220) (1.2080,-0.8177) (1.2090,-0.7990) (1.2100,-0.7747) (1.2109,-0.8210) (1.2119,-0.8442) (1.2129,-0.8681) (1.2139,-0.9237) (1.2148,-0.9146) (1.2158,-0.9101) (1.2168,-0.9259) (1.2178,-0.8977) (1.2188,-0.9186) (1.2197,-0.9551) (1.2207,-0.9506) (1.2217,-0.9832) (1.2227,-0.9808) (1.2236,-0.9316) (1.2246,-0.9216) (1.2256,-0.8781) (1.2266,-0.8380) (1.2275,-0.8555) (1.2285,-0.8341) (1.2295,-0.8230) (1.2305,-0.8283) (1.2314,-0.7662) (1.2324,-0.7254) (1.2334,-0.6954) (1.2344,-0.6334) (1.2354,-0.6370) (1.2363,-0.6472) (1.2373,-0.6336) (1.2383,-0.6662) (1.2393,-0.6556) (1.2402,-0.6194) (1.2412,-0.6240) (1.2422,-0.5893) (1.2432,-0.5760) (1.2441,-0.6090) (1.2451,-0.5961) (1.2461,-0.6038) (1.2471,-0.6086) (1.2480,-0.5455) (1.2490,-0.5114) (1.2500,-0.4696) (1.2510,-0.4025) (1.2520,-0.4021) (1.2529,-0.3889) (1.2539,-0.3617) (1.2549,-0.3752) (1.2559,-0.3335) (1.2568,-0.2839) (1.2578,-0.2705) (1.2588,-0.2178) (1.2598,-0.2077) (1.2607,-0.2366) (1.2617,-0.2268) (1.2627,-0.2557) (1.2637,-0.2727) (1.2646,-0.2374) (1.2656,-0.2445) (1.2666,-0.2344) (1.2676,-0.2155) (1.2686,-0.2624) (1.2695,-0.2807) (1.2705,-0.2948) (1.2715,-0.3349) (1.2725,-0.3056) (1.2734,-0.2776) (1.2744,-0.2673) (1.2754,-0.2120) (1.2764,-0.2067) (1.2773,-0.2189) (1.2783,-0.1940) (1.2793,-0.2115) (1.2803,-0.1997) (1.2812,-0.1486) (1.2822,-0.1441) (1.2832,-0.1130) (1.2842,-0.0929) (1.2852,-0.1362) (1.2861,-0.1453) (1.2871,-0.1688) (1.2881,-0.2101) (1.2891,-0.1843) (1.2900,-0.1787) (1.2910,-0.1804) (1.2920,-0.1458) (1.2930,-0.1714) (1.2939,-0.1967) (1.2949,-0.1915) (1.2959,-0.2257) (1.2969,-0.2099) (1.2979,-0.1632) (1.2988,-0.1526) (1.2998,-0.0992) (1.3008,-0.0659) (1.3018,-0.0787) (1.3027,-0.0469) (1.3037,-0.0389) (1.3047,-0.0319) (1.3057,0.0379) (1.3066,0.0726) (1.3076,0.1095) (1.3086,0.1657) (1.3096,0.1502) (1.3105,0.1438) (1.3115,0.1483) (1.3125,0.1109) (1.3135,0.1294) (1.3145,0.1570) (1.3154,0.1522) (1.3164,0.1915) (1.3174,0.1934) (1.3184,0.1634) (1.3193,0.1787) (1.3203,0.1621) (1.3213,0.1642) (1.3223,0.2240) (1.3232,0.2454) (1.3242,0.2874) (1.3252,0.3391) (1.3262,0.3247) (1.3271,0.3371) (1.3281,0.3497) (1.3291,0.3314) (1.3301,0.3764) (1.3311,0.4127) (1.3320,0.4243) (1.3330,0.4733) (1.3340,0.4650) (1.3350,0.4336) (1.3359,0.4342) (1.3369,0.3885) (1.3379,0.3707) (1.3389,0.3913) (1.3398,0.3665) (1.3408,0.3711) (1.3418,0.3682) (1.3428,0.3067) (1.3438,0.2850) (1.3447,0.2540) (1.3457,0.2107) (1.3467,0.2391) (1.3477,0.2503) (1.3486,0.2569) (1.3496,0.3004) (1.3506,0.2827) (1.3516,0.2634) (1.3525,0.2704) (1.3535,0.2329) (1.3545,0.2397) (1.3555,0.2699) (1.3564,0.2570) (1.3574,0.2799) (1.3584,0.2756) (1.3594,0.2198) (1.3604,0.2041) (1.3613,0.1611) (1.3623,0.1162) (1.3633,0.1340) (1.3643,0.1185) (1.3652,0.1093) (1.3662,0.1240) (1.3672,0.0730) (1.3682,0.0384) (1.3691,0.0216) (1.3701,-0.0298) (1.3711,-0.0182) (1.3721,0.0079) (1.3730,0.0067) (1.3740,0.0529) (1.3750,0.0632) (1.3760,0.0422) (1.3770,0.0651) (1.3779,0.0530) (1.3789,0.0562) (1.3799,0.1119) (1.3809,0.1238) (1.3818,0.1513) (1.3828,0.1829) (1.3838,0.1440) (1.3848,0.1290) (1.3857,0.1116) (1.3867,0.0630) (1.3877,0.0789) (1.3887,0.0881) (1.3896,0.0770) (1.3906,0.1086) (1.3916,0.0889) (1.3926,0.0537) (1.3936,0.0575) (1.3945,0.0224) (1.3955,0.0223) (1.3965,0.0662) (1.3975,0.0693) (1.3984,0.1054) (1.3994,0.1351) (1.4004,0.1064) (1.4014,0.1158) (1.4023,0.1120) (1.4033,0.0915) (1.4043,0.1368) (1.4053,0.1579) (1.4062,0.1679) (1.4072,0.2075) (1.4082,0.1794) (1.4092,0.1448) (1.4102,0.1318) (1.4111,0.0714) (1.4121,0.0545) (1.4131,0.0608) (1.4141,0.0260) (1.4150,0.0305) (1.4160,0.0117) (1.4170,-0.0529) (1.4180,-0.0717) (1.4189,-0.1122) (1.4199,-0.1487) (1.4209,-0.1179) (1.4219,-0.1169) (1.4229,-0.1068) (1.4238,-0.0722) (1.4248,-0.1039) (1.4258,-0.1210) (1.4268,-0.1243) (1.4277,-0.1670) (1.4287,-0.1524) (1.4297,-0.1303) (1.4307,-0.1422) (1.4316,-0.1134) (1.4326,-0.1273) (1.4336,-0.1774) (1.4346,-0.1877) (1.4355,-0.2359) (1.4365,-0.2693) (1.4375,-0.2498) (1.4385,-0.2717) (1.4395,-0.2735) (1.4404,-0.2663) (1.4414,-0.3229) (1.4424,-0.3483) (1.4434,-0.3689) (1.4443,-0.4130) (1.4453,-0.3858) (1.4463,-0.3599) (1.4473,-0.3491) (1.4482,-0.2927) (1.4492,-0.2861) (1.4502,-0.2946) (1.4512,-0.2657) (1.4521,-0.2788) (1.4531,-0.2605) (1.4541,-0.2035) (1.4551,-0.1924) (1.4561,-0.1535) (1.4570,-0.1265) (1.4580,-0.1619) (1.4590,-0.1623) (1.4600,-0.1782) (1.4609,-0.2108) (1.4619,-0.1766) (1.4629,-0.1642) (1.4639,-0.1586) (1.4648,-0.1180) (1.4658,-0.1390) (1.4668,-0.1591) (1.4678,-0.1507) (1.4688,-0.1823) (1.4697,-0.1637) (1.4707,-0.1167) (1.4717,-0.1063) (1.4727,-0.0537) (1.4736,-0.0237) (1.4746,-0.0398) (1.4756,-0.0124) (1.4766,-0.0110) (1.4775,-0.0104) (1.4785,0.0510) (1.4795,0.0753) (1.4805,0.1019) (1.4814,0.1454) (1.4824,0.1160) (1.4834,0.0961) (1.4844,0.0851) (1.4854,0.0320) (1.4863,0.0351) (1.4873,0.0453) (1.4883,0.0236) (1.4893,0.0455) (1.4902,0.0285) (1.4912,-0.0191) (1.4922,-0.0219) (1.4932,-0.0574) (1.4941,-0.0726) (1.4951,-0.0309) (1.4961,-0.0276) (1.4971,-0.0021) (1.4980,0.0323) (1.4990,0.0017) (1.5000,-0.0000) (1.5010,-0.0017) (1.5020,-0.0323) (1.5029,0.0021) (1.5039,0.0276) (1.5049,0.0309) (1.5059,0.0726) (1.5068,0.0574) (1.5078,0.0219) (1.5088,0.0191) (1.5098,-0.0285) (1.5107,-0.0455) (1.5117,-0.0236) (1.5127,-0.0453) (1.5137,-0.0351) (1.5146,-0.0320) (1.5156,-0.0851) (1.5166,-0.0961) (1.5176,-0.1160) (1.5186,-0.1454) (1.5195,-0.1019) (1.5205,-0.0753) (1.5215,-0.0510) (1.5225,0.0104) (1.5234,0.0110) (1.5244,0.0124) (1.5254,0.0398) (1.5264,0.0237) (1.5273,0.0537) (1.5283,0.1063) (1.5293,0.1167) (1.5303,0.1637) (1.5312,0.1823) (1.5322,0.1507) (1.5332,0.1591) (1.5342,0.1390) (1.5352,0.1180) (1.5361,0.1586) (1.5371,0.1642) (1.5381,0.1766) (1.5391,0.2108) (1.5400,0.1782) (1.5410,0.1623) (1.5420,0.1619) (1.5430,0.1265) (1.5439,0.1535) (1.5449,0.1924) (1.5459,0.2035) (1.5469,0.2605) (1.5479,0.2788) (1.5488,0.2657) (1.5498,0.2946) (1.5508,0.2861) (1.5518,0.2927) (1.5527,0.3491) (1.5537,0.3599) (1.5547,0.3858) (1.5557,0.4130) (1.5566,0.3689) (1.5576,0.3483) (1.5586,0.3229) (1.5596,0.2663) (1.5605,0.2735) (1.5615,0.2717) (1.5625,0.2498) (1.5635,0.2693) (1.5645,0.2359) (1.5654,0.1877) (1.5664,0.1774) (1.5674,0.1273) (1.5684,0.1134) (1.5693,0.1422) (1.5703,0.1303) (1.5713,0.1524) (1.5723,0.1670) (1.5732,0.1243) (1.5742,0.1210) (1.5752,0.1039) (1.5762,0.0722) (1.5771,0.1068) (1.5781,0.1169) (1.5791,0.1179) (1.5801,0.1487) (1.5811,0.1122) (1.5820,0.0717) (1.5830,0.0529) (1.5840,-0.0117) (1.5850,-0.0305) (1.5859,-0.0260) (1.5869,-0.0608) (1.5879,-0.0545) (1.5889,-0.0714) (1.5898,-0.1318) (1.5908,-0.1448) (1.5918,-0.1794) (1.5928,-0.2075) (1.5938,-0.1679) (1.5947,-0.1579) (1.5957,-0.1368) (1.5967,-0.0915) (1.5977,-0.1120) (1.5986,-0.1158) (1.5996,-0.1064) (1.6006,-0.1351) (1.6016,-0.1054) (1.6025,-0.0693) (1.6035,-0.0662) (1.6045,-0.0223) (1.6055,-0.0224) (1.6064,-0.0575) (1.6074,-0.0537) (1.6084,-0.0889) (1.6094,-0.1086) (1.6104,-0.0770) (1.6113,-0.0881) (1.6123,-0.0789) (1.6133,-0.0630) (1.6143,-0.1116) (1.6152,-0.1290) (1.6162,-0.1440) (1.6172,-0.1829) (1.6182,-0.1513) (1.6191,-0.1238) (1.6201,-0.1119) (1.6211,-0.0562) (1.6221,-0.0530) (1.6230,-0.0651) (1.6240,-0.0422) (1.6250,-0.0632) (1.6260,-0.0529) (1.6270,-0.0067) (1.6279,-0.0079) (1.6289,0.0182) (1.6299,0.0298) (1.6309,-0.0216) (1.6318,-0.0384) (1.6328,-0.0730) (1.6338,-0.1240) (1.6348,-0.1093) (1.6357,-0.1185) (1.6367,-0.1340) (1.6377,-0.1162) (1.6387,-0.1611) (1.6396,-0.2041) (1.6406,-0.2198) (1.6416,-0.2756) (1.6426,-0.2799) (1.6436,-0.2570) (1.6445,-0.2699) (1.6455,-0.2397) (1.6465,-0.2329) (1.6475,-0.2704) (1.6484,-0.2634) (1.6494,-0.2827) (1.6504,-0.3004) (1.6514,-0.2569) (1.6523,-0.2503) (1.6533,-0.2391) (1.6543,-0.2107) (1.6553,-0.2540) (1.6562,-0.2850) (1.6572,-0.3067) (1.6582,-0.3682) (1.6592,-0.3711) (1.6602,-0.3665) (1.6611,-0.3913) (1.6621,-0.3707) (1.6631,-0.3885) (1.6641,-0.4342) (1.6650,-0.4336) (1.6660,-0.4650) (1.6670,-0.4733) (1.6680,-0.4243) (1.6689,-0.4127) (1.6699,-0.3764) (1.6709,-0.3314) (1.6719,-0.3497) (1.6729,-0.3371) (1.6738,-0.3247) (1.6748,-0.3391) (1.6758,-0.2874) (1.6768,-0.2454) (1.6777,-0.2240) (1.6787,-0.1642) (1.6797,-0.1621) (1.6807,-0.1787) (1.6816,-0.1634) (1.6826,-0.1934) (1.6836,-0.1915) (1.6846,-0.1522) (1.6855,-0.1570) (1.6865,-0.1294) (1.6875,-0.1109) (1.6885,-0.1483) (1.6895,-0.1438) (1.6904,-0.1502) (1.6914,-0.1657) (1.6924,-0.1095) (1.6934,-0.0726) (1.6943,-0.0379) (1.6953,0.0319) (1.6963,0.0389) (1.6973,0.0469) (1.6982,0.0787) (1.6992,0.0659) (1.7002,0.0992) (1.7012,0.1526) (1.7021,0.1632) (1.7031,0.2099) (1.7041,0.2257) (1.7051,0.1915) (1.7061,0.1967) (1.7070,0.1714) (1.7080,0.1458) (1.7090,0.1804) (1.7100,0.1787) (1.7109,0.1843) (1.7119,0.2101) (1.7129,0.1688) (1.7139,0.1453) (1.7148,0.1362) (1.7158,0.0929) (1.7168,0.1130) (1.7178,0.1441) (1.7188,0.1486) (1.7197,0.1997) (1.7207,0.2115) (1.7217,0.1940) (1.7227,0.2189) (1.7236,0.2067) (1.7246,0.2120) (1.7256,0.2673) (1.7266,0.2776) (1.7275,0.3056) (1.7285,0.3349) (1.7295,0.2948) (1.7305,0.2807) (1.7314,0.2624) (1.7324,0.2155) (1.7334,0.2344) (1.7344,0.2445) (1.7354,0.2374) (1.7363,0.2727) (1.7373,0.2557) (1.7383,0.2268) (1.7393,0.2366) (1.7402,0.2077) (1.7412,0.2178) (1.7422,0.2705) (1.7432,0.2839) (1.7441,0.3335) (1.7451,0.3752) (1.7461,0.3617) (1.7471,0.3889) (1.7480,0.4021) (1.7490,0.4025) (1.7500,0.4696) (1.7510,0.5114) (1.7520,0.5455) (1.7529,0.6086) (1.7539,0.6038) (1.7549,0.5961) (1.7559,0.6090) (1.7568,0.5760) (1.7578,0.5893) (1.7588,0.6240) (1.7598,0.6194) (1.7607,0.6556) (1.7617,0.6662) (1.7627,0.6336) (1.7637,0.6472) (1.7646,0.6370) (1.7656,0.6334) (1.7666,0.6954) (1.7676,0.7254) (1.7686,0.7662) (1.7695,0.8283) (1.7705,0.8230) (1.7715,0.8341) (1.7725,0.8555) (1.7734,0.8380) (1.7744,0.8781) (1.7754,0.9216) (1.7764,0.9316) (1.7773,0.9808) (1.7783,0.9832) (1.7793,0.9506) (1.7803,0.9551) (1.7812,0.9186) (1.7822,0.8977) (1.7832,0.9259) (1.7842,0.9101) (1.7852,0.9146) (1.7861,0.9237) (1.7871,0.8681) (1.7881,0.8442) (1.7891,0.8210) (1.7900,0.7747) (1.7910,0.7990) (1.7920,0.8177) (1.7930,0.8220) (1.7939,0.8695) (1.7949,0.8638) (1.7959,0.8445) (1.7969,0.8598) (1.7979,0.8315) (1.7988,0.8363) (1.7998,0.8765) (1.8008,0.8702) (1.8018,0.8932) (1.8027,0.9011) (1.8037,0.8480) (1.8047,0.8313) (1.8057,0.7968) (1.8066,0.7483) (1.8076,0.7665) (1.8086,0.7608) (1.8096,0.7512) (1.8105,0.7752) (1.8115,0.7371) (1.8125,0.7033) (1.8135,0.6970) (1.8145,0.6519) (1.8154,0.6605) (1.8164,0.6959) (1.8174,0.6970) (1.8184,0.7426) (1.8193,0.7642) (1.8203,0.7433) (1.8213,0.7675) (1.8223,0.7649) (1.8232,0.7655) (1.8242,0.8268) (1.8252,0.8505) (1.8262,0.8799) (1.8271,0.9251) (1.8281,0.8983) (1.8291,0.8846) (1.8301,0.8787) (1.8311,0.8331) (1.8320,0.8463) (1.8330,0.8648) (1.8340,0.8538) (1.8350,0.8873) (1.8359,0.8800) (1.8369,0.8449) (1.8379,0.8544) (1.8389,0.8298) (1.8398,0.8282) (1.8408,0.8813) (1.8418,0.8950) (1.8428,0.9328) (1.8438,0.9765) (1.8447,0.9556) (1.8457,0.9651) (1.8467,0.9715) (1.8477,0.9505) (1.8486,0.9944) (1.8496,1.0257) (1.8506,1.0359) (1.8516,1.0824) (1.8525,1.0688) (1.8535,1.0365) (1.8545,1.0340) (1.8555,0.9848) (1.8564,0.9675) (1.8574,0.9854) (1.8584,0.9585) (1.8594,0.9641) (1.8604,0.9585) (1.8613,0.8974) (1.8623,0.8785) (1.8633,0.8472) (1.8643,0.8079) (1.8652,0.8401) (1.8662,0.8520) (1.8672,0.8629) (1.8682,0.9082) (1.8691,0.8907) (1.8701,0.8760) (1.8711,0.8848) (1.8721,0.8501) (1.8730,0.8635) (1.8740,0.8966) (1.8750,0.8887) (1.8760,0.9186) (1.8770,0.9176) (1.8779,0.8691) (1.8789,0.8614) (1.8799,0.8238) (1.8809,0.7887) (1.8818,0.8146) (1.8828,0.8049) (1.8838,0.8051) (1.8848,0.8260) (1.8857,0.7812) (1.8867,0.7570) (1.8877,0.7475) (1.8887,0.7060) (1.8896,0.7303) (1.8906,0.7654) (1.8916,0.7762) (1.8926,0.8347) (1.8936,0.8539) (1.8945,0.8460) (1.8955,0.8810) (1.8965,0.8791) (1.8975,0.8966) (1.8984,0.9637) (1.8994,0.9860) (1.9004,1.0273) (1.9014,1.0688) (1.9023,1.0415) (1.9033,1.0414) (1.9043,1.0358) (1.9053,1.0025) (1.9062,1.0350) (1.9072,1.0570) (1.9082,1.0622) (1.9092,1.1088) (1.9102,1.1012) (1.9111,1.0823) (1.9121,1.1003) (1.9131,1.0784) (1.9141,1.0956) (1.9150,1.1531) (1.9160,1.1703) (1.9170,1.2233) (1.9180,1.2657) (1.9189,1.2526) (1.9199,1.2795) (1.9209,1.2900) (1.9219,1.2877) (1.9229,1.3505) (1.9238,1.3857) (1.9248,1.4130) (1.9258,1.4672) (1.9268,1.4517) (1.9277,1.4337) (1.9287,1.4343) (1.9297,1.3883) (1.9307,1.3890) (1.9316,1.4090) (1.9326,1.3897) (1.9336,1.4112) (1.9346,1.4054) (1.9355,1.3573) (1.9365,1.3554) (1.9375,1.3286) (1.9385,1.3098) (1.9395,1.3561) (1.9404,1.3697) (1.9414,1.3957) (1.9424,1.4423) (1.9434,1.4221) (1.9443,1.4202) (1.9453,1.4285) (1.9463,1.3994) (1.9473,1.4299) (1.9482,1.4636) (1.9492,1.4659) (1.9502,1.5087) (1.9512,1.5049) (1.9521,1.4686) (1.9531,1.4707) (1.9541,1.4324) (1.9551,1.4127) (1.9561,1.4425) (1.9570,1.4295) (1.9580,1.4395) (1.9590,1.4542) (1.9600,1.4063) (1.9609,1.3925) (1.9619,1.3798) (1.9629,1.3465) (1.9639,1.3855) (1.9648,1.4188) (1.9658,1.4398) (1.9668,1.5047) (1.9678,1.5163) (1.9688,1.5166) (1.9697,1.5514) (1.9707,1.5429) (1.9717,1.5694) (1.9727,1.6304) (1.9736,1.6455) (1.9746,1.6907) (1.9756,1.7196) (1.9766,1.6885) (1.9775,1.6940) (1.9785,1.6803) (1.9795,1.6535) (1.9805,1.6925) (1.9814,1.7058) (1.9824,1.7154) (1.9834,1.7568) (1.9844,1.7343) (1.9854,1.7165) (1.9863,1.7239) (1.9873,1.6915) (1.9883,1.7124) (1.9893,1.7573) (1.9902,1.7672) (1.9912,1.8204) (1.9922,1.8466) (1.9932,1.8300) (1.9941,1.8569) (1.9951,1.8543) (1.9961,1.8546) (1.9971,1.9132) (1.9980,1.9317) (1.9990,1.9555) (2.0000,1.9922)};
\end{axis}
\end{tikzpicture}}}{Two continuous functions that are differentiable nowhere. (A) The Takagi function $T$ on $[0,1]$ (solid), plotted at the $513$ points $j/512$, where it equals its partial sum $T_8$, together with the partial sums $T_0=\varphi$, $T_1$ and $T_2$; its maximum $2/3$ is attained, among other points, at $1/3$ and $2/3$. (B) Weierstrass's function $\sum a^n\cos(b^n\pi x)$ for $a=1/2$ and $b=3$, approximated by its first eight terms at $2049$ points of $[0,2]$; the omitted terms change it by at most $2^{-7}$. Its period is $2$.\label{fig:nowhere-differentiable}}
\end{figure}

The two functions are extreme forms of the phenomenon in \cref{tab:interchange-failures}: uniform convergence of the partial sums controls their values but not their slopes, and here the slopes of the partial sums grow without bound everywhere. Such functions are not rare. Chapter~5 proves, with Baire's theorem of \cref{sec:baire}, Banach's theorem that in the complete metric space $C[0,1]$ the functions that are nowhere differentiable form a residual set: in the sense of category, a typical continuous function is nowhere differentiable~\cite{banach1931}.

\subsection{A Curve That Fills a Square}

Uniform limits also produce continuous curves with surprising images. A continuous map of $[0,1]$ onto the square $[0,1]^2$ was found by Peano in 1890~\cite{peano1890curve}, and Hilbert gave a geometric construction the following year~\cite{hilbert1891}. The following construction, due to Schoenberg~\cite{schoenberg1938}, uses series of the kind studied in this chapter. Let $p\colon\mathbb R\to[0,1]$ be the even function with the period $2$ given on $[0,1]$ by
\[
p(t)\coloneqq\begin{cases}0,&0\le t\le1/3,\\3t-1,&1/3\le t\le2/3,\\1,&2/3\le t\le1 ,\end{cases}
\]
which is continuous, and let
\[
\Phi(t)\coloneqq\sum_{n=0}^\infty\frac{p(3^{2n}t)}{2^{n+1}},\qquad\Psi(t)\coloneqq\sum_{n=0}^\infty\frac{p(3^{2n+1}t)}{2^{n+1}}\qquad(0\le t\le1).
\]

\begin{proposition}[A space-filling curve]
\label{prop:space-filling}
The map $t\mapsto(\Phi(t),\Psi(t))$ is continuous from $[0,1]$ onto $[0,1]^2$.
\end{proposition}

\begin{proof}
Both series converge uniformly by the M-test with $M_n=2^{-n-1}$, so $\Phi$ and $\Psi$ are continuous, with values in $[0,\sum_n2^{-n-1}]=[0,1]$. Let $(x,y)\in[0,1]^2$, and choose binary expansions $x=\sum_{n\ge0}a_{2n}2^{-n-1}$ and $y=\sum_{n\ge0}a_{2n+1}2^{-n-1}$ with digits $a_k\in\{0,1\}$; this defines $a_k$ for every $k\in\mathbb N$. Let
\[
t\coloneqq\sum_{k=0}^\infty\frac{2a_k}{3^{k+1}}\in[0,1] ,
\]
a point of the Cantor set of \cref{def:cantor-set}. We show that $p(3^kt)=a_k$ for every $k$. Indeed,
\[
3^kt=\sum_{j<k}2a_j3^{k-j-1}+r_k,\qquad r_k\coloneqq\frac{2a_k}3+\sum_{j>k}\frac{2a_j}{3^{j-k+1}} ,
\]
where the first sum is an even integer and $0\le\sum_{j>k}2a_j3^{-(j-k+1)}\le2\sum_{i\ge2}3^{-i}=1/3$. Since $p$ has the period $2$, $p(3^kt)=p(r_k)$. If $a_k=0$, then $0\le r_k\le1/3$ and $p(r_k)=0$; if $a_k=1$, then $2/3\le r_k\le1$ and $p(r_k)=1$. Hence $\Phi(t)=\sum_na_{2n}2^{-n-1}=x$ and $\Psi(t)=\sum_na_{2n+1}2^{-n-1}=y$.
\end{proof}

The curve cannot be injective. If it were, it would be a homeomorphism $h$ of $[0,1]$ onto the square, by \cref{thm:compact-homeomorphism}, and $h$ would map the disconnected set $[0,1]\setminus\{1/2\}$ onto the square without the point $h(1/2)$, which is path-connected, since two of its points $q$ and $q'$ are joined by the segment $[q,q']$ or, if this contains $h(1/2)$, by two segments through a point of the square off the line through $q$ and $q'$; so it is connected by \cref{prop:path-connected}\,(i). This is impossible, since the inverse of $h$ is continuous and maps connected sets to connected sets (\cref{thm:connected-image}). The coordinate functions $\Phi$ and $\Psi$ are themselves nowhere differentiable, a fact that we do not prove here; Sagan's book describes these and many other space-filling curves~\cite{sagan1994}. The existence of such curves shows that the dimension of a set is not preserved by continuous maps, while \cref{thm:compact-homeomorphism} and Brouwer's invariance of dimension, mentioned in the Notes of \cref{ch:compactness}, show that it is preserved by homeomorphisms.

\section{Uniform Errors in Computation}
\label{sec:uniform-computation}

A numerical method that replaces a function by a partial sum, a polynomial or an iterate is judged by its error over a whole interval, and that error is the distance in the supremum norm. A computer cannot evaluate a supremum over infinitely many points; it evaluates the error on a finite grid, and the maximum over the grid is a lower bound for the supremum, accurate when the grid resolves the features of the error. This section measures uniform errors in four examples from the chapter, compares them with the bounds of the theorems, and shows where a grid fails to see the supremum. The computations were made on an Apple M4 Pro processor with Python~3.15 and its binary64 floating-point numbers, using the array operations and the random number generator of the version of the stardust library of 25 November 2026; reference values were computed with the ball arithmetic of python-flint (version 0.9.0), with error bounds far below the digits shown.

\subsection{Suprema of Partial Sums}

The two series $\sum_{k\ge1}\cos(kx)/k^2$ and $\sum_{k\ge1}\sin(kx)/k$ have the sums
\[
f(x)=\frac{\pi^2}6-\frac{\pi x}2+\frac{x^2}4\quad(0\le x\le2\pi),\qquad g(x)=\frac{\pi-x}2\quad(0<x<2\pi),
\]
as Chapter~9 proves. The first converges uniformly by the M-test, with the remainder bound $\sum_{k>n}1/k^2$, and the second converges uniformly away from the multiples of $2\pi$ only, by \cref{sec:function-series}. We evaluated the partial sums $C_n$ and $S_n$ of the two series, adding the terms one by one, on the grid of the $2^{16}+1$ points $j\pi/2^{16}$, $j=0,\ldots,2^{16}$, of $[0,\pi]$; both sums are symmetric about $\pi$ up to sign, so this interval suffices. \Cref{tab:sup-errors} shows the results.

\begin{table}
\sdfloat{%
\begin{sdtable}{colspec={rllllll}, row{1}={font=\small\bfseries}}
\toprule
$n$ & $\sum_{k>n}k^{-2}$ & $\max|C_n-f|$ & $\max_{[0.1,\pi]}|S_n-g|$ & $g-S_n$ at $1/n$ & $\max(S_n-g)$ \\
\midrule
$10$ & $9.517\times10^{-2}$ & $9.517\times10^{-2}$ & $5.83\times10^{-1}$ & $0.5829$ & $0.2823$ \\
$10^2$ & $9.950\times10^{-3}$ & $9.950\times10^{-3}$ & $8.45\times10^{-2}$ & $0.6205$ & $0.2812$ \\
$10^3$ & $9.995\times10^{-4}$ & $9.995\times10^{-4}$ & $9.95\times10^{-3}$ & $0.6243$ & $0.2811$ \\
$10^4$ & $1.000\times10^{-4}$ & $1.000\times10^{-4}$ & $9.95\times10^{-4}$ & $0.6247$ & $0.2811$ \\
\bottomrule
\end{sdtable}}{Uniform errors of the partial sums $C_n$ of $\sum\cos(kx)/k^2$ and $S_n$ of $\sum\sin(kx)/k$. The second and third columns are the M-test bound and the largest error of $C_n$ on the grid of $[0,\pi]$; the fourth is the largest error of $S_n$ on the grid points of $[0.1,\pi]$; the fifth is the error of $S_n$ at the single point $1/n$; the last is the largest amount by which $S_n$ exceeds $g$, on a grid of $4000$ points of $(0,2\pi/n]$.\label{tab:sup-errors}}
\end{table}

For the cosine series, the largest error on the grid agrees with the M-test bound in all the digits shown, and in fact to at least nine digits. This is no coincidence: at $x=0$ all terms $\cos(kx)/k^2$ attain their bounds $M_k=1/k^2$ simultaneously, so the remainder there is exactly $\sum_{k>n}1/k^2$, and the grid contains $x=0$. The M-test bound is sharp for this series, and it decreases like $1/n$, by \vcref{1}{cor:p-series}.

For the sine series, the fourth column shows uniform convergence on $[0.1,\pi]$, with errors close to $1/(0.1\,n)$ for large $n$; the estimate in the proof of \cref{thm:abel-dirichlet}, with $M=1/\sin(0.05)$ and $b_k=1/k$, bounds the remainder by $4/((n+1)\sin0.05)$, about $80/n$, which is correct in its order but eight times too large. The fifth column shows the failure of uniform convergence near $0$ by the sequential test, \cref{prop:sequential-test}: the error at the point $1/n$ does not tend to $0$. Its limit can be computed. The partial sum $S_n(1/n)=(1/n)\sum_{k\le n}\sin(k/n)/(k/n)$ is a Riemann sum of the continuous function $\sin t/t$ on $[0,1]$, so it tends to $\int_0^1\sin t/t\,dt$ (\vcref{1}{thm:riemann-sums}), and $g(1/n)\to\pi/2$; the limit of the error is $\pi/2-\int_0^1\sin t/t\,dt=0.62471$, which the column approaches. The last column is the overshoot of \cref{fig:sawtooth}: in the same way, $S_n(\pi/n)\to\int_0^\pi\sin t/t\,dt$ while $g(\pi/n)\to\pi/2$, and the largest excess converges to $\int_0^\pi\sin t/t\,dt-\pi/2=0.28114$, the constant of the Gibbs phenomenon of Chapter~9. On the whole interval $(0,\pi]$, the supremum of $|S_n-g|$ is at least $\pi/2$ for every $n$, since $S_n(x)\to0$ and $g(x)\to\pi/2$ as $x\to0^+$.

The fixed grid of $[0,\pi]$ does not show this. For $n=10^4$ the largest error of $S_n$ on that grid was $1.097$, attained at the first grid point after $0$, instead of a value close to $\pi/2$, because the grid spacing $\pi/2^{16}=4.8\times10^{-5}$ is no longer small compared with the width $1/n$ of the region where the error is large; and the excess $S_n-g$ on that grid reached only $0.274$, short of $0.2811$. A grid can only resolve what is wider than its spacing, and the region of nonuniform convergence shrinks as $n$ grows. The sequential test with points adapted to $n$, as in the fifth column, is the reliable way to detect nonuniform convergence numerically.

\subsection{Square Roots by Polynomials}

The polynomials $p_n$ of \cref{prop:sqrt-polynomials} satisfy $0\le\sqrt x-p_n(x)\le2/(n+1)$. We evaluated them by the recursion on the grid $x=t^2$ with $t=j/2^{20}$, $j=0,\ldots,2^{20}$, which is fine near $0$, where the error is largest. For $n=16$, $64$, $256$ and $1024$, the largest error on the grid was $0.0344$, $0.00868$, $0.00217$ and $0.000544$, while the bound is $0.118$, $0.0308$, $0.00778$ and $0.00195$; the ratio of error to bound decreased from $0.29$ to $0.28$. The products of $n$ with the largest error were $0.551$, $0.555$, $0.557$ and $0.557$, so the error behaves like $0.557/n$, and it was attained near $\sqrt x\approx1.28/n$: the slow, algebraic rate predicted by the proposition is the true rate, and the region of largest error moves into the point $0$, where the square root has an infinite slope. No error on the grid was negative, in agreement with $p_n\le\sqrt x$, although the iteration was carried out in floating-point arithmetic. As an approximation method the iteration is wasteful: the degree $2^{n-1}$ of $p_n$ grows exponentially while the error decreases like $1/n$, and far better polynomials of the same degree exist. Their construction, by best approximation in the supremum norm, is the subject of a note in Chapter~5 and of Volume~XII.

\subsection{Abel's Theorem at the End of the Interval}

By Abel's theorem the series $\sum_{k\ge1}(-1)^{k+1}x^k/k$ converges uniformly on $[0,1]$ to $\ln(1+x)$. We computed the remainders $R_n(x)\coloneqq\ln(1+x)-\sum_{k\le n}(-1)^{k+1}x^k/k$ on the grid of the $2^{16}+1$ points $j/2^{16}$ of $[0,1]$, with the logarithm function of the library. For $n=10$, $100$ and $1000$, the largest of $|R_n|$ over the grid was $4.75\times10^{-2}$, $4.975\times10^{-3}$ and $5.00\times10^{-4}$, attained in each case at $x=1$; these are close to $1/(2n)$. On the part $[0,0.9]$ of the grid, the largest remainders were $1.57\times10^{-2}$ and $1.25\times10^{-7}$ for $n=10$ and $100$, below the bounds $0.9^{n+1}/(n+1)$ of \vcref{1}{thm:log-series}, and for $n=1000$ the computed value $2.0\times10^{-15}$ consisted of rounding errors only. Uniform convergence on $[0,1]$ thus takes place at the slow rate of the end point, while on every $[0,r]$ with $r<1$ the convergence is geometric. Abel's theorem guarantees convergence at the end, not speed: to compute $\ln2$ to within $10^{-10}$ from the alternating harmonic series takes about $5\times10^9$ terms.

The partial sums at $x=1$ alternate about $\ln2$, and averaging neighbours cancels much of the error. Let $s_n$ be the $n$th partial sum, $s_n^{(0)}\coloneqq s_n$, and $s_n^{(j+1)}\coloneqq(s_n^{(j)}+s_{n+1}^{(j)})/2$. From the $n$ partial sums $s_1,\ldots,s_n$, $n-1$ averaging steps give the single value
\[
E_n\coloneqq s_1^{(n-1)}=\frac1{2^{n-1}}\sum_{j=0}^{n-1}\binom{n-1}js_{j+1} ,
\]
a weighted mean of the partial sums with binomial weights, which is the form of Euler's transformation of an alternating series that Knopp describes~\cite{knopp1990}. \Cref{tab:ln2-averaging} compares $s_n$, the mean $(s_{n-1}+s_n)/2$ of the last two partial sums, and $E_n$.

\begin{table}
\sdfloat{%
\begin{sdtable}{colspec={rlll}, row{1}={font=\small\bfseries}}
\toprule
$n$ & $s_n-\ln2$ & $(s_{n-1}+s_n)/2-\ln2$ & $E_n-\ln2$ \\
\midrule
$5$ & $9.02\times10^{-2}$ & $-9.81\times10^{-3}$ & $1.64\times10^{-3}$ \\
$10$ & $-4.75\times10^{-2}$ & $2.49\times10^{-3}$ & $1.53\times10^{-5}$ \\
$20$ & $-2.44\times10^{-2}$ & $6.24\times10^{-4}$ & $4.17\times10^{-9}$ \\
$30$ & $-1.64\times10^{-2}$ & $2.78\times10^{-4}$ & $1.89\times10^{-12}$ \\
\bottomrule
\end{sdtable}}{Errors of three values computed from the first $n$ partial sums of $1-1/2+1/3-\cdots=\ln2$: the last partial sum, the mean of the last two, and the binomial mean $E_n$ obtained by $n-1$ steps of averaging neighbours. The values were computed in binary64 arithmetic and agree with exact rational computations in the digits shown.\label{tab:ln2-averaging}}
\end{table}

The partial sums have errors close to $\pm1/(2n)$, as \vcref{1}{sec:absolute} observed; one average reduces them to about $1/(4n^2)$, and the full averaging $E_n$ gains a little more than a factor of $2$ per partial sum used, so that thirty partial sums give almost twelve correct digits. Like Ces\`aro's means, $E_n$ is a regular summation method: its weights are positive, their sum is $1$, and each of them tends to $0$ as $n\to\infty$, and such means converge to $s$ whenever $s_n\to s$~\cite{knopp1990}. Its speed here depends on the structure of the terms $1/k$, which are the values of a smooth function, and no transformation accelerates every series; Volume~XII treats the acceleration of series systematically.

\subsection{Chord Slopes of the Takagi Function}

The proof of \cref{thm:takagi} showed that the chord slopes $D_m$ of the Takagi function over the dyadic intervals $[k2^{-m},(k+1)2^{-m}]$ that contain a point $x$ are sums of $m$ signs $\pm1$. We computed them from the definition, evaluating $T$ at the ends of the intervals by \eqref{eq:takagi}. At a dyadic point $u=k2^{-m}$ all terms with $n\ge m$ vanish, and the others are multiples of $2^{-m}$, so in binary64 arithmetic the values $T(u)$ and the slopes are exact for $m\le48$. The computed slopes agreed with the walk of the binary digits described after the theorem. \Cref{fig:takagi-slopes} shows them for the binary64 number nearest to $1/3$, whose first $53$ binary digits are those of $1/3$, and for the point $x=0.6462854957694724$, the first number drawn by the uniform random number generator of the library with the seed $20261127$.

\begin{figure}
\sdfloat{%
\begin{tikzpicture}
\begin{axis}[sdplot, width=0.6\textwidth, height=0.36\textwidth, xmin=0, xmax=49, ymin=-9, ymax=4.2, xtick={1,8,16,24,32,40,48}, ytick={-8,-6,-4,-2,0,2}, xlabel={$m$}, sdlegendtop]
\addplot[black!45, line width=0.6pt, mark=o, mark size=1.1pt] coordinates {(1,1) (2,0) (3,1) (4,0) (5,1) (6,0) (7,1) (8,0) (9,1) (10,0) (11,1) (12,0) (13,1) (14,0) (15,1) (16,0) (17,1) (18,0) (19,1) (20,0) (21,1) (22,0) (23,1) (24,0) (25,1) (26,0) (27,1) (28,0) (29,1) (30,0) (31,1) (32,0) (33,1) (34,0) (35,1) (36,0) (37,1) (38,0) (39,1) (40,0) (41,1) (42,0) (43,1) (44,0) (45,1) (46,0) (47,1) (48,0)};
\addlegendentry{$x\approx1/3$}
\addplot[black, line width=0.6pt, mark=*, mark size=1.1pt] coordinates {(1,-1) (2,0) (3,-1) (4,0) (5,1) (6,0) (7,1) (8,0) (9,1) (10,0) (11,-1) (12,-2) (13,-1) (14,0) (15,-1) (16,0) (17,-1) (18,-2) (19,-3) (20,-4) (21,-3) (22,-4) (23,-5) (24,-6) (25,-5) (26,-6) (27,-5) (28,-6) (29,-7) (30,-8) (31,-7) (32,-6) (33,-5) (34,-4) (35,-5) (36,-6) (37,-5) (38,-6) (39,-5) (40,-6) (41,-7) (42,-6) (43,-5) (44,-4) (45,-3) (46,-2) (47,-3) (48,-4)};
\addlegendentry{$x=0.64628\ldots$}
\end{axis}
\end{tikzpicture}}{The slopes $D_m$ of the chords of the Takagi function over the dyadic intervals of length $2^{-m}$ that contain $x$, for $m=1,\ldots,48$, at a point close to $1/3$ (grey circles) and at a random point (dots). Each step changes the slope by $\pm1$, so the slopes cannot converge; at $1/3$ they alternate between $1$ and $0$, at the random point they perform a random walk.\label{fig:takagi-slopes}}
\end{figure}

At the point close to $1/3$ the slopes alternate between $1$ and $0$, and at the random point they wander between $-8$ and $1$; at $m=48$ the slope was $-4$. For a point whose binary digits behave like independent fair coin tosses, the slope $D_m$ is a simple random walk, which by the central limit theorem of Volume~IV is typically of the order $\sqrt m$, here about $7$ for $m=48$. Difference quotients of $T$ with small steps $h$ behave in the same erratic way, and no refinement of $h$ settles them. A numerical differentiation routine applied to such a function returns numbers, and only the theory can say that they approximate nothing.

\bigskip
A sequence of functions converges uniformly if the supremum of the errors tends to zero, which for bounded functions is convergence in the space $\ell^\infty(S)$; a uniformly Cauchy sequence with values in a complete space converges uniformly, and a sequence that converges only pointwise can be tested with one point for each index. Uniform convergence allows two limits to be exchanged when the inner limits exist, and so uniform limits of continuous functions are continuous; on compact sets, monotone pointwise convergence to a continuous limit is uniform, by Dini's theorem, and it gives polynomials that converge uniformly to $\sqrt x$ and to $|x|$. Uniform limits of integrable functions are integrable, with the limit of the integrals; improper integrals need a dominating function; and derivatives pass to the limit when the derivatives converge uniformly, which makes $C^k[a,b]$ a Banach space. For series, the M-test and the tests of Abel and Dirichlet establish uniform convergence, and the rules for term-by-term continuity, integration and differentiation follow, as do the smoothness of the zeta function and the convergence of infinite products. Nonnegative double series may be summed in any order, and so may absolutely summable ones, while other double series can give different iterated sums. Power series converge uniformly on compact subsets of their interval of convergence, their sums are analytic there and are determined by their values on any set with an accumulation point, and by Abel's theorem they are continuous up to an end at which the series converges; Ces\`aro summability implies Abel summability, and Tauber's theorem recovers convergence from Abel summability for terms of order $o(1/k)$. Finally, uniform limits can be continuous and nowhere differentiable, as the functions of Takagi and Weierstrass show, and a uniform limit of curves can fill a square.

The next chapter turns the question around: instead of asking which properties survive uniform limits, it asks which functions are uniform limits of simple ones. Every continuous function on a closed bounded interval is a uniform limit of polynomials, by the Weierstrass approximation theorem, and the Stone--Weierstrass theorem, which uses the polynomials that converge to $|x|$ found above, extends this to continuous functions on compact metric spaces. The same chapter characterises the compact sets of continuous functions by equicontinuity, and proves with Baire's theorem that nowhere differentiable functions are the typical continuous functions.

\chapternotes

Cauchy's theorem on the continuity of sums of series~\cite{cauchy1821} and Abel's counterexample~\cite{abel1826} started a discussion that lasted a generation. In 1847 Seidel~\cite{seidel1847} and Stokes~\cite{stokes1847} independently located the gap in Cauchy's argument: near a point of discontinuity of the sum, the convergence becomes arbitrarily slow, in Seidel's words, which is nonuniform convergence in modern terms. Weierstrass used uniform convergence in his work on power series of the 1840s and made it, together with the M-test, one of the foundations of his Berlin lectures, through which it became standard; Bottazzini~\cite{bottazzini1986} and Hairer and Wanner~\cite{hairer1996analysis} tell this history. Dini's theorem is in his \emph{Fondamenti} of 1878~\cite{dini1878}, and the theorem on increasing functions with a continuous limit is a lemma in P\'olya's paper of 1920 on the central limit theorem~\cite{polya1920}. The theorem on bounded convergence of Riemann integrals is due to Arzel\`a~\cite{arzela1885}, and Osgood~\cite{osgood1897} rediscovered it for continuous functions in his study of term-by-term integration; Hawkins~\cite{hawkins1975} describes how such results led to Lebesgue's integral. The theorem on the exchange of limits is known as the Moore--Osgood theorem. The observation that differentiation is a closed operator is the starting point of the theory of unbounded operators, for which Kreyszig~\cite[Sect.~4.13]{kreyszig1978} gives an introduction.

Abel's paper of 1826 on the binomial series contains both his theorem on the continuity of power series at the end of the interval of convergence and his theorem on products of series~\cite{abel1826}. Ces\`aro introduced his means in a paper on the multiplication of series~\cite{cesaro1890}, and Frobenius had shown ten years earlier that convergence of these means implies Abel summability~\cite{frobenius1880}. Tauber's theorem is from 1897~\cite{tauber1897}, and Littlewood's extension, which opened the modern theory, from 1911~\cite{littlewood1911}. Hardy's \emph{Divergent Series}~\cite{hardy1949divergent} is the classical account of summation methods and Tauberian theorems, and Korevaar's book~\cite{korevaar2004} follows their development through the twentieth century. Knopp~\cite{knopp1990} treats double series, infinite products and summability with many examples, and Krantz and Parks~\cite{krantz2002primer} develop the theory of real analytic functions, including the change of centre and the identity theorem.

Weierstrass presented his nowhere differentiable function to the Berlin Academy in 1872, and du Bois-Reymond published it with its proof in 1875~\cite{duboisreymond1875}. Hardy's paper of 1916~\cite{hardy1916} settled the case $ab\ge1$. Takagi's function appeared in 1903~\cite{takagi1903}, and van der Waerden's variant with base $10$ in 1930~\cite{vanderwaerden1930}; Allaart and Kawamura survey the many properties of the Takagi function that have been found since~\cite{allaart2012}. The first space-filling curves are Peano's of 1890~\cite{peano1890curve} and Hilbert's of 1891~\cite{hilbert1891}; the analytic construction used here is Schoenberg's~\cite{schoenberg1938}, and Sagan's book~\cite{sagan1994} collects the subject. Banach's theorem that nowhere differentiable functions are typical~\cite{banach1931} is proved in Chapter~5.

The material of this chapter is the core of every course in real analysis. Rudin~\cite[Chap.~7]{rudin1976} treats uniform convergence and its consequences concisely, Apostol~\cite[Chaps.~8 and~9]{apostol1974} at greater length, with summability methods and Tauber's theorem, and Abbott~\cite[Chap.~6]{abbott2015} with a careful discussion of the counterexamples. Titchmarsh~\cite[Chap.~1]{titchmarsh1939} gives the classical treatment of uniform convergence, double series and the exchange of limits that underlies the theory of functions.

\chapterreferences
